% !TeX program = xelatex
% 数学 LaTeX 排版 Skill v4.0.0；版式保持 v3 金标准. XeLaTeX 编译至少两次. 不分发字体.
% WZ-LOCKED-CLASS-BEGIN
\begin{filecontents*}[overwrite]{elegantbook-original-adapter.cls}
\NeedsTeXFormat{LaTeX2e}
\ProvidesClass{elegantbook-original-adapter}[2026/09/23 v3.0 fixed mathematical book environments]
\LoadClass[10pt,letterpaper,oneside]{book}
\RequirePackage[UTF8,scheme=plain,fontset=fandol]{ctex}
\RequirePackage{fontspec}
\setmainfont{cmunrm.otf}[BoldFont=cmunbx.otf,ItalicFont=cmunti.otf,BoldItalicFont=cmunbi.otf]
\setCJKfamilyfont{sourceheader}[ItalicFont=FandolKai-Regular.otf]{FandolKai-Regular.otf}
\RequirePackage{amsmath,amssymb,mathrsfs}
\RequirePackage{amsthm,etoolbox}
\RequirePackage{xcolor,geometry,setspace,fancyhdr,enumitem,needspace,xurl}
\definecolor{structurecolor}{HTML}{880E4F}
\definecolor{main}{HTML}{9C27B0}
\geometry{paperwidth=612bp,paperheight=792bp,left=20mm,right=20mm,top=95.04bp,bottom=20mm,headheight=12pt,headsep=25pt,footskip=30pt}
\setstretch{1.6}
\setlength{\parindent}{2em}
\setlength{\parskip}{0pt}
\setlist{nosep,leftmargin=2em}
\AtBeginDocument{\setlength{\abovedisplayskip}{4pt plus 1pt minus 1pt}\setlength{\belowdisplayskip}{4pt plus 1pt minus 1pt}\setlength{\abovedisplayshortskip}{2pt}\setlength{\belowdisplayshortskip}{3pt}}
\fancyhf{}
\fancyhead[L]{{\itshape\CJKfamily{sourceheader}\rightmark}}
\fancyhead[R]{{\itshape\CJKfamily{sourceheader}\leftmark}}
\fancyfoot[C]{\thepage}
\renewcommand{\headrulewidth}{0.4pt}
\renewcommand{\headrule}{{\color{structurecolor}\hrule height\headrulewidth width\headwidth\vskip-\headrulewidth}}
\pagestyle{fancy}
\fancypagestyle{cover}{\fancyhf{}\fancyfoot[C]{\thepage}\renewcommand{\headrulewidth}{0pt}\renewcommand{\headrule}{}}
\RequirePackage{hyperref}
\hypersetup{unicode=true,colorlinks=true,linkcolor=structurecolor,urlcolor=structurecolor,citecolor=structurecolor,linktoc=section,bookmarksnumbered=true,bookmarksopen=true,pdfborder={0 0 0},pdfstartview=FitH}
\setcounter{tocdepth}{1}
\renewcommand*\l@section{\@dottedtocline{1}{1.5em}{2.4em}}
\renewcommand{\thesection}{\thechapter.\arabic{section}}
\newcommand{\originalchapter}[1]{\par\addvspace{14pt}\Needspace{5\baselineskip}\refstepcounter{chapter}\setcounter{section}{0}\addcontentsline{toc}{chapter}{\protect\numberline{\thechapter}#1}\markboth{\thechapter\quad #1}{}\begingroup\centering\color{structurecolor}\bfseries\fontsize{14.4}{20}\selectfont\thechapter\quad #1\par\endgroup\nobreak\vspace{8pt}}
\newcommand{\originalsection}[1]{\par\addvspace{9pt}\Needspace{4\baselineskip}\refstepcounter{section}\addcontentsline{toc}{section}{\protect\hspace*{1.5em}\protect\numberline{\protect\textcolor{structurecolor}{\thesection}}#1}\markright{\thesection\quad #1}\noindent{\fontsize{12}{17}\selectfont\color{structurecolor}\thesection\quad}{\bfseries\fontsize{12}{17}\selectfont #1}\par\nobreak\vspace{4pt}}

% Semantic environments are registered once here, not re-designed in manuscripts.
\newtheoremstyle{wzstatement}{6pt}{5pt}
  {\normalfont\kaishu}{0pt}{\normalfont\bfseries\color{structurecolor}}
  {}{.7em}{\thmname{#1}\thmnumber{ #2}\thmnote{\normalfont\ (#3)}}
\newtheoremstyle{wzdefinition}{6pt}{5pt}
  {\normalfont\songti}{0pt}{\normalfont\bfseries\color{structurecolor}}
  {}{.7em}{\thmname{#1}\thmnumber{ #2}\thmnote{\normalfont\ (#3)}}
\newtheoremstyle{wzproblem}{7pt}{5pt}
  {\normalfont\songti}{2em}{\normalfont\bfseries\color{main}}
  {}{.7em}{\thmname{#1}\thmnumber{ #2}}
\newtheoremstyle{wzremark}{4pt}{4pt}
  {\normalfont\songti}{0pt}{\normalfont\bfseries}
  {}{.7em}{\thmname{#1}}
\theoremstyle{wzstatement}
\newtheorem{theorem}{定理}[chapter]
\newtheorem{lemma}[theorem]{引理}
\newtheorem{proposition}[theorem]{命题}
\newtheorem{corollary}[theorem]{推论}
\theoremstyle{wzdefinition}
\newtheorem{definition}[theorem]{定义}
\theoremstyle{wzproblem}
\newtheorem{example}{例题}[chapter]
\newtheorem{exercise}{习题}[chapter]
\theoremstyle{wzremark}
\newtheorem*{remark}{注}
\renewcommand{\proofname}{证明}
% Retain the native amsthm QED stack; only adjust the fixed typography.
\expandafter\patchcmd\csname\string\proof\endcsname{\normalfont \topsep6\p@\@plus6\p@\relax}
  {\normalfont\songti \topsep5\p@\@plus1\p@\relax}
  {}{\ClassError{elegantbook-original-adapter}{Cannot patch proof spacing}{Check the amsthm version.}}
\expandafter\patchcmd\csname\string\proof\endcsname{\itshape}{\normalfont\bfseries}
  {}{\ClassError{elegantbook-original-adapter}{Cannot patch proof heading}{Check the amsthm version.}}
\expandafter\patchcmd\csname\string\proof\endcsname{\ignorespaces}
  {\setlength{\parindent}{2em}\ignorespaces}
  {}{\ClassError{elegantbook-original-adapter}{Cannot patch proof paragraphs}{Check the amsthm version.}}
\AtBeginEnvironment{proof}{\par\Needspace{3\baselineskip}}
\renewcommand{\qedsymbol}{\ensuremath{\square}}
\newenvironment{solution}{\begin{proof}[解答]}{\end{proof}}
\newcommand{\wzcheckchapter}{\ifnum\value{chapter}<1
  \ClassError{elegantbook-original-adapter}{Numbered mathematics before chapter 1}{Move it after the first originalchapter.}\fi}

\newcommand{\wzregisterguard}[1]{\AtBeginEnvironment{#1}{\wzcheckchapter\par\Needspace{3\baselineskip}}}
\forcsvlist{\wzregisterguard}{theorem,lemma,proposition,corollary,definition,example,exercise}
% Front matter and bibliography are not numbered mathematical chapters.
\newcommand{\originalpreface}{\par\addvspace{10pt}\Needspace{5\baselineskip}
  \noindent{\bfseries\fontsize{12}{17}\selectfont 前言}\par\nobreak\vspace{4pt}}
\newcommand{\originalcontents}{\par\addvspace{14pt}
  \begin{center}{\color{structurecolor}\bfseries\fontsize{14.4}{20}\selectfont 目录}\end{center}
  \vspace{-10pt}\@starttoc{toc}}
\renewcommand{\bibname}{参考文献}
\renewenvironment{thebibliography}[1]{
  \par\addvspace{12pt}\Needspace{3\baselineskip}\phantomsection
  \addcontentsline{toc}{chapter}{\bibname}
  \markboth{\bibname}{}
  \noindent{\color{structurecolor}\bfseries\fontsize{14.4}{20}\selectfont\bibname}\par\nobreak\vspace{6pt}
  \list{\@biblabel{\@arabic\c@enumiv}}
    {\settowidth\labelwidth{\@biblabel{#1}}\leftmargin\labelwidth
     \advance\leftmargin\labelsep\usecounter{enumiv}
     \let\p@enumiv\@empty\renewcommand\theenumiv{\@arabic\c@enumiv}
     \setlength{\itemsep}{3pt}\setlength{\parsep}{0pt}}
  \sloppy\clubpenalty4000\widowpenalty4000\sfcode`\.=1000\relax
}{\def\@noitemerr{\@latex@warning{Empty bibliography}}\endlist}

\numberwithin{equation}{chapter}
\renewcommand{\theequation}{\thechapter.\arabic{equation}}
\allowdisplaybreaks[2]
\emergencystretch=2em
\clubpenalty=6000
\widowpenalty=6000
\raggedbottom
\endinput
\end{filecontents*}
% WZ-LOCKED-CLASS-END
\documentclass{elegantbook-original-adapter}
% ===== 元数据内容槽 =====
\newcommand{\BookTitle}{有限群表示论}
\hypersetup{pdftitle={有限群表示论: MIT 18.712 中文重构本},pdfauthor={Pavel Etingof 等原讲义; 中文重构}}
\usepackage{tikz,booktabs,longtable}
\usetikzlibrary{arrows.meta,positioning,calc}
\newcommand{\C}{\mathbb C}
\newcommand{\R}{\mathbb R}
\newcommand{\Z}{\mathbb Z}
\newcommand{\Hom}{\operatorname{Hom}}
\newcommand{\End}{\operatorname{End}}
\newcommand{\GL}{\operatorname{GL}}
\newcommand{\tr}{\operatorname{tr}}
\newcommand{\Ind}{\operatorname{Ind}}
\newcommand{\Res}{\operatorname{Res}}
\newcommand{\Irr}{\operatorname{Irr}}
\newcommand{\Cl}{\operatorname{Cl}}
\newcommand{\sgn}{\operatorname{sgn}}
\newcommand{\Sym}{\operatorname{Sym}}
\newcommand{\im}{\operatorname{im}}
\newcommand{\id}{\operatorname{id}}
\newcommand{\inner}[2]{\langle #1,#2\rangle}
\begin{document}
\frontmatter
\begin{center}
{\color{structurecolor}\bfseries\fontsize{18}{25}\selectfont\BookTitle\par}
{\bfseries MIT 18.712 中文重构本\par}
Pavel Etingof 等原讲义; 中文编排与增补\par
Introduction to Representation Theory, Fall 2010\par
2026年10月8日\par
\end{center}
\clearpage
\phantomsection\addcontentsline{toc}{chapter}{课程信息与教学进度}
\noindent{\bfseries\fontsize{12}{17}\selectfont 课程信息与教学进度}\par
官方课程为MIT数学系18.712《表示论导论》(Introduction to Representation Theory), 2010年秋季本科课程, 教师Pavel Etingof. 课程目标是介绍群、李代数与结合代数的表示论, 重点以线性空间中的对称及具体例子建立理解. 每周两次讲课, 每次1.5小时. 官网页面未列具体上课日期、钟点或教室, 此处只保留周次, 不推造日历.

官方先修为18.701、18.702(代数I、II), 或18.700、18.703(线性代数与现代代数), 要求扎实的线性代数及群、环、域基础. 本中文本前九章有限群主线补充必要群论; 后部课程原题涉及更广的代数、李代数与箭图, 所需前置在相应题解中补明. 以下大纲是原课程范围, 不等于本书正文全部覆盖.

原大纲包括: 结合代数、生成元与关系、群代数、路径代数、李代数与包络代数、不可约与不可分解、Schur引理和 $\mathfrak{sl}_2$; 稠密性、半单代数、Jordan--Hölder与Krull--Schmidt分解、表示扩张; 有限群的Maschke、平方和、对偶与张量、特征标与矩阵元素正交、酉表示和特征标表; Frobenius--Schur指标、群行列式、代数整数与整除性、Burnside定理、诱导、互反、有限一般线性群、对称群、Schur--Weyl对偶及不变量理论; 箭图、子空间三元组、根系、反射函子和Gabriel定理.

\begin{longtable}{c p{.22\textwidth} p{.60\textwidth}}
\toprule 周次&官方旧讲义阅读范围&内容参考\\\midrule\endhead
1&§1.1--1.5&代数、表示、理想与商\\
2&§1.6--1.13&生成元关系、箭图、李代数、张量与对偶\\
3&§1.14--2.8&$\mathfrak{sl}_2$、一般表示论、分解定理\\
4&§2.9--3.5&一般练习、有限群完全可约、特征标正交\\
5&§3.6--4.1&酉表示、矩阵系数、特征标表与指标\\
6&§4.2--4.9&群行列式、整除性、Burnside、诱导\\
7&§4.10--4.19&互反、对称群、钩长与Schur--Weyl对偶\\
8&§4.20--4.24&Schur多项式、一般线性群表示\\
9&§4.25--5.2&Artin、半直积、箭图入门\\
10&§5.3--5.5&$D_4$、根系与Gabriel定理\\
11&§5.6--5.9&反射函子、Coxeter元素及证明\\
12&第6章&范畴、函子、伴随与阿贝尔范畴\\\bottomrule
\end{longtable}
阅读书目除旧讲义外, 官网列 William Fulton与Joe Harris, \textit{Representation Theory: A First Course}, Springer, 1991, GTM 129; Jean-Pierre Serre, \textit{Linear Representations of Finite Groups}, Springer-Verlag, 1977, GTM 42, 均使用开头部分. 这些是课程参考书目, 本项目不归档其受版权保护的书籍, 也不以它们替换公开旧讲义.

官方考核为作业75\%、带回家期末作业25\%. 作业每周四布置, 下一周四截止; 允许合作, 要求学生实际参与并理解所写内容. 作业页列11次作业, 共51道必做原题及1道选做原题1.58; 期末PDF有4道大题. 未列独立测验、期中卷、练习考试或官方参考答案. 本书将这些公开题目全部中文化纳入同一主PDF, 保留原编号; 解答是编者独立推导, 与旧讲义给出的提示分别标明. 官方页面: \href{https://ocw.mit.edu/courses/18-712-introduction-to-representation-theory-fall-2010/pages/syllabus/}{大纲}, \href{https://ocw.mit.edu/courses/18-712-introduction-to-representation-theory-fall-2010/pages/readings/}{阅读进度}, \href{https://ocw.mit.edu/courses/18-712-introduction-to-representation-theory-fall-2010/pages/assignments/}{作业与期末}.

\clearpage
\originalpreface
群的表示把抽象对称写成线性变换. 对有限群, 平均运算使不变子空间总有不变补空间, 从而把研究任意表示的问题化为研究不可约表示. 特征标把这些线性变换压缩为共轭类上的数值, 正交关系又使这些数值足以恢复表示的分解. 诱导与限制则把一个群的表示和其子群的表示连接起来.

本书假定读者会复向量空间、基与矩阵、特征值、秩、迹和有限维内积空间. 群、子群、陪集与共轭类所需的事实在第一章补齐. 除特别说明外, 群有限, 表示空间为有限维复向量空间. Maschke定理另外说明一般域中的条件, 并给出模特征反例; 本书不展开模表示理论.

主要依据为 MIT OpenCourseWare \href{https://ocw.mit.edu/courses/18-712-introduction-to-representation-theory-fall-2010/pages/lecture-notes/}{18.712, Fall 2010} 的公开旧讲义, 完整PDF共109页, 首页标明2011年2月1日. 原作者为 Pavel Etingof、Oleg Golberg、Sebastian Hensel、Tiankai Liu、Alex Schwendner、Elena Yudovina 与 Dmitry Vaintrob. 本书选取旧讲义第一、二章中必要基础、第三章有限群表示基本理论, 以及第四章中的直积、诱导、Mackey分解、Frobenius互反等内容, 重新安排依赖关系, 以中文重写定义、证明、计算和题解. 英文原件仅在仓库sources目录归档, 不嵌入本PDF, 也不收入中文源码ZIP.

循环群、对称群 $S_3$ 与 $S_4$、二面体群、四元数群和 $A_4$ 提供贯穿例子. 正则表示、矩阵系数、群代数的矩阵分块及Frobenius行列式均给出证明. 对称平方、外平方、置换表示与平均投影的展开, 以及部分练习解答, 是编者为衔接主线补充的论证, 不冒称原课完整题解.

本书不是旧讲义全部章节的译本. 正文围绕有限群表示展开; 作业与期末题解按需要补出李代数与 $\mathfrak{sl}_2$、箭图、Young模、Schur--Weyl交换代数及 $\GL_2(\mathbb F_q)$ 的局部论证, 不把它们称为这些专题的完整教材. 范畴论、一般有限维代数的根基与投射覆盖、Frobenius整除性、Burnside可解性与Artin定理均未作为系统专题纳入. 群代数的半单分块是有限群主线的一部分, 不表示已经覆盖一般有限维代数. 书末给出逐节来源与取舍表.

旧讲义末页为MIT OCW许可说明; 默认 \href{https://creativecommons.org/licenses/by-nc-sa/4.0/}{CC BY-NC-SA 4.0} 及 \href{https://ocw.mit.edu/pages/privacy-and-terms-of-use/}{OCW使用条款}适用, 单独权利标记优先. 本中文改编同许可非商业共享, 并注明翻译、重组、增补及自绘图等改动; 不表示MIT或作者认可本项目. 官方另链的2016年AMS出版书明示禁止上传到非作者网站, 本项目未下载归档或改编该版本. 新增课程信息、11次作业的全部52道题(含选做)和4道期末题均置于本主PDF. 这些题解会涉及若干正文未系统展开的主题, 不把所需局部论证称为该主题的完整教材. 来源核查与文件校验见同目录source-manifest.json.
\clearpage
\originalcontents
\clearpage
\mainmatter
\originalchapter{群与线性表示}
\originalsection{有限群的必要语言}
一个群 $G$ 是带有结合乘法的集合, 有单位元 $e$, 每个元素 $g$ 有逆元 $g^{-1}$. 我们把乘积 $gh$ 理解为先作 $h$ 再作 $g$, 这与线性映射复合的记号相同. 群的阶 $|G|$ 是元素数. 子群 $H\le G$ 包含单位元并对乘法、取逆封闭. 左陪集 $gH$ 是集合 $\{gh:h\in H\}$; 两个左陪集或者相同, 或者不相交: 若 $g_1h_1=g_2h_2$, 则 $g_2^{-1}g_1=h_2h_1^{-1}\in H$, 从而 $g_1H=g_2H$. 因此陪集把 $G$ 分成等大的块, 每块有 $|H|$ 个元素, 所以
\[
 |G|=[G:H]|H|,
\]
其中 $[G:H]$ 为陪集数. 特别地, 元素 $g$ 的阶是最小的正整数 $m$ 使 $g^m=e$, 并且 $m$ 整除 $|G|$.

群作用在集合 $X$ 上, 指对每个 $g\in G$ 有置换 $x\mapsto gx$, 且 $ex=x$, $(gh)x=g(hx)$. 元素 $x$ 的轨道为 $Gx$, 稳定子为 $G_x=\{g:gx=x\}$. 映射 $gG_x\mapsto gx$ 良定义且为双射, 因为 $g_1x=g_2x$ 等价于 $g_2^{-1}g_1\in G_x$. 这证明轨道大小为 $[G:G_x]$.

把作用取为共轭 $g\cdot x=gxg^{-1}$, 就得到共轭类. 元素 $x$ 的中心化子 $C_G(x)=\{g:gx=xg\}$ 是其稳定子, 所以共轭类的大小是 $|G|/|C_G(x)|$. 群中心 $Z(G)$ 是与所有元素交换的元素组成的子群. 单元素共轭类恰由中心元素构成. 后面的特征标只在共轭类上取值, 因而先弄清这些类能大幅减少计算量.

\begin{example}\label{ex:s3classes}
求三次对称群 $S_3$ 的共轭类与类大小.
\end{example}
\begin{solution}
$S_3$ 是集合 $\{1,2,3\}$ 的所有置换. 对任意置换 $g$, 有 $g(i\ j)g^{-1}=(g(i)\ g(j))$, 三循环也按同样方式更换标签. 因此三个换位 $(12),(13),(23)$ 共轭, 两个三循环 $(123),(132)$ 共轭, 单位元自成一类. 各类大小为 $1,3,2$, 总和为6. 相应中心化子的阶为 $6,2,3$.
\end{solution}

\originalsection{表示与矩阵}
\begin{definition}
群 $G$ 的复表示是复向量空间 $V$ 及群同态 $\rho:G\to\GL(V)$. 即 $\rho(e)=I$, $\rho(gh)=\rho(g)\rho(h)$. 其维数是 $\dim V$. 简写 $gv=\rho(g)v$.
\end{definition}
选基后, 表示就是一组可逆矩阵 $\rho(g)$. 改变基使所有矩阵同时共轭. 所以某个矩阵的具体元素依赖基, 但维数、迹和整个表示的分解是内在对象. 只知道一个群元素的矩阵通常不够; 群的关系必须全部满足.

\begin{definition}
表示 $V,W$ 之间的同态, 或交织映射, 是线性映射 $T:V\to W$, 满足 $T\rho_V(g)=\rho_W(g)T$ 对一切 $g\in G$ 成立. 所有这种映射组成向量空间 $\Hom_G(V,W)$. 可逆交织映射称同构; 此时两个表示等价, 记 $V\cong W$.
\end{definition}
在同一组基中, 同构等价于存在一个可逆矩阵 $T$ 对所有 $g$ 同时满足 $\rho_W(g)=T\rho_V(g)T^{-1}$. 每个表示都可取直和: 在 $V\oplus W$ 上让 $g(v,w)=(gv,gw)$, 其矩阵为分块对角矩阵.

\begin{example}
设 $C_n=\langle r:r^n=e\rangle$. 描述其一维复表示.
\end{example}
\begin{solution}
一维空间上的可逆线性变换是乘以非零复数. 令 $\rho(r)=z$, 群关系要求 $z^n=1$. 因此 $z=\zeta^j$, 其中 $\zeta=\exp(2\pi i/n)$, $j=0,\ldots,n-1$. 记此表示为 $L_j$, 则 $\rho_j(r^a)=\zeta^{ja}$. 任意一维交织同构是非零标量, 所以不同 $j$ 的表示不等价. 第五章将证明这些就是 $C_n$ 的全部不可约表示.
\end{solution}

\originalsection{不变子空间与商表示}
\begin{definition}
子空间 $U\subseteq V$ 若满足 $gU\subseteq U$ 对所有 $g\in G$ 成立, 称不变子空间或子表示. 非零表示 $V$ 若只有 $0,V$ 两个子表示, 称不可约. 若 $V$ 不能写成两个非零子表示的直和, 称不可分解. 若 $V$ 是不可约表示的直和, 称完全可约.
\end{definition}
因为 $g^{-1}$ 也在群中, 条件 $gU\subseteq U$ 实际蕴含 $gU=U$. 当 $U$ 不变时, 在商空间上定义 $g(v+U)=gv+U$; 若更换代表元 $v$ 为 $v+u$, 两个结果只差 $gu\in U$, 所以这是良定义的商表示 $V/U$.

\begin{proposition}\label{prop:kernelimage}
交织映射的核与像都是子表示. 非零有限维表示总有不可约子表示.
\end{proposition}
\begin{proof}
若 $Tv=0$, 则 $T(gv)=gTv=0$; 若 $w=Tv$, 则 $gw=T(gv)$ 仍在像中. 对后一结论, 从所有非零子表示中选取维数最小者 $U$. 若 $U$ 有非零真子表示, 其维数更小, 矛盾.
\end{proof}
不可约必不可分解, 反向结论一般不成立. 对有限群复表示, 第二章的Maschke定理会保证反向成立. 定理的价值恰在这里: 一个任意选取的向量空间补空间未必不变, 需要利用群的有限性把它修正.

\begin{example}\label{ex:permstandard}
$S_3$ 置换 $\C^3$ 的三个坐标. 找到一个一维子表示及其不变补空间.
\end{example}
\begin{solution}
令 $e_1,e_2,e_3$ 为标准基, 定义 $g e_j=e_{g(j)}$. 直线 $L=\C(e_1+e_2+e_3)$ 上作用恒等. 平面 $U=\{(a,b,c):a+b+c=0\}$ 也不变, 因为置换不改变坐标和. 任意 $v=(a,b,c)$ 可唯一写成
\[
 v=\frac{a+b+c}{3}(1,1,1)+\left(v-\frac{a+b+c}{3}(1,1,1)\right),
\]
从而 $\C^3=L\oplus U$. 在基 $u_1=e_1-e_2,u_2=e_2-e_3$ 中, $s=(12)$ 与 $r=(123)$ 的矩阵是
\[
 \rho_U(s)=\begin{pmatrix}-1&1\\0&1\end{pmatrix},\qquad
 \rho_U(r)=\begin{pmatrix}0&-1\\1&-1\end{pmatrix}.
\]
直接相乘可核对 $s^2=e$, $r^3=e$, $srs=r^{-1}$. 第五章再证明 $U$ 不可约.
\end{solution}

\originalsection{群代数}
\begin{definition}
群代数 $\C[G]$ 是以符号 $[g]$ ($g\in G$) 为基的向量空间, 乘法由 $[g][h]=[gh]$ 双线性延拓. 其单位元是 $[e]$. 下文常把 $[g]$ 简写为 $g$.
\end{definition}
群代数把群的乘法与线性组合放进一个对象. 它是含单位元的结合代数: 结合律在基元素上由群结合律保证, 再由双线性推广. $\dim\C[G]=|G|$. 给定群表示, 公式
\[
 \rho\left(\sum_{g\in G}a_g g\right)=\sum_{g\in G}a_g\rho(g)
\]
定义一个保单位元的代数同态 $\C[G]\to\End(V)$. 反过来, 这样的代数同态限制到 $g$ 就是群表示, 因为 $g\cdot g^{-1}=e$ 保证像可逆. 因而群表示和群代数的左模是同一结构的两种语言. 本书只用群代数所需的代数知识, 不预先调用一般半单代数结构定理.

\begin{exercise}
设 $N\trianglelefteq G$ 是正规子群. 证明 $G/N$ 的表示可看成 $N$ 作用恒等的 $G$ 表示, 且二者的不可约性一致.
\end{exercise}
\begin{solution}
若 $\bar\rho:G/N\to\GL(V)$, 令 $\rho(g)=\bar\rho(gN)$, 则 $N$ 中元素作用恒等. 反过来, 若 $N\subseteq\ker\rho$, 则 $gN=hN$ 蕴含 $h^{-1}g\in N$, 所以 $\rho(g)=\rho(h)$; 因而可定义 $\bar\rho(gN)=\rho(g)$. 两种作用给出的线性变换集合相同, 不变子空间也相同, 所以不可约性一致.
\end{solution}

\originalchapter{交织映射与完全可约性}
\originalsection{Schur引理}
\begin{theorem}[Schur引理]\label{thm:schur}
设 $V,W$ 为不可约表示. 非零交织映射 $T:V\to W$ 必为同构. 如果 $V$ 是有限维复不可约表示, 则 $\End_G(V)=\C I_V$. 因此
\[
 \dim\Hom_G(V,W)=\begin{cases}1,&V\cong W,\\0,&V\not\cong W.\end{cases}
\]
\end{theorem}
\begin{proof}
由命题\ref{prop:kernelimage}, $\ker T$ 与 $\im T$ 都不变. 非零性排除 $\ker T=V$ 与 $\im T=0$; 不可约性迫使 $\ker T=0$, $\im T=W$, 所以 $T$ 同构.

若 $T\in\End_G(V)$, 复数域保证其特征多项式有根 $\lambda$, 因而 $T-\lambda I$ 的核非零. 它仍为交织映射, 若非零就应可逆, 与非零核矛盾. 所以 $T=\lambda I$. 若 $V\cong W$, 固定一个同构 $S:V\to W$, 则任意 $T$ 满足 $S^{-1}T\in\C I$, 即 $T$ 为 $S$ 的标量倍; 若不等价, 第一部分排除所有非零 $T$.
\end{proof}
这里两个条件发挥不同作用. 不可约性控制核与像, 复数域和有限维性提供特征值. 例如 $\R^2$ 上90度旋转生成的 $C_4$ 实表示不可约, 但其交织算子包括旋转矩阵本身, 并不都是实标量. 第一部分的核像论证仍然有效.

\begin{corollary}\label{cor:abelian}
有限交换群的每个复不可约表示都是一维的. 对任意群, 中心元素在不可约表示上作用为标量.
\end{corollary}
\begin{proof}
若 $z\in Z(G)$, $\rho(z)$ 与所有 $\rho(g)$ 交换, 所以Schur引理适用. 当整个 $G$ 交换时, 每个 $g$ 都作用为标量, 任意直线都不变. 不可约性迫使维数为1.
\end{proof}
中心作用的标量随 $z$ 乘法相乘, 从而给出 $Z(G)\to\C^\times$ 的同态, 称中心特征标. 它通常不能决定整个不可约表示, 因为中心以外还保留信息.

\originalsection{平均投影与Maschke定理}
\begin{theorem}[Maschke定理]\label{thm:maschke}
设 $G$ 为有限群, $k$ 为域, 且 $|G|$ 在 $k$ 中非零. 对任意有限维 $k$ 表示 $V$ 及子表示 $U$, 存在子表示 $W$ 使 $V=U\oplus W$. 因而每个有限维表示都完全可约. 特别地, 结论对 $k=\C$ 成立.
\end{theorem}
\begin{proof}
先取普通线性投影 $P:V\to U$, 使 $P|_U=I_U$; 它来自任意向量空间补空间, 未必交织. 我们希望修正投影而保留它在 $U$ 上的恒等性质. 对任意 $g$, 共轭算子 $\rho(g)P\rho(g)^{-1}$ 的像仍在 $U$ 内, 在 $U$ 上仍恒等. 因此可以平均:
\[
 Q=\frac1{|G|}\sum_{g\in G}\rho(g)P\rho(g)^{-1}.
\]
有限性使和式有意义, 底域条件使除以 $|G|$ 合法. 对任意 $h\in G$, 左乘 $h$ 置换求和指标, 所以 $\rho(h)Q\rho(h)^{-1}=Q$. 又 $Q(V)\subseteq U$, $Q|_U=I_U$, 因而 $Q^2=Q$. 令 $W=\ker Q$, 其不变性由 $Q\rho(h)=\rho(h)Q$ 保证. 每个 $v$ 写成 $Qv+(v-Qv)$, 两项分别在 $U,W$; 交集为0, 所以 $V=U\oplus W$.

若 $V\ne0$, 先选一个不可约子表示, 取其不变补空间, 再对较低维补空间归纳. 维数严格下降, 过程有限步终止, 得到不可约直和分解.
\end{proof}
不能仅把 $P$ 的矩阵各元素平均后就认定它还是投影; 平均一般不保持幂等性. 此处幂等性成立是因为每个被平均的算子像都在 $U$ 内、在 $U$ 上都恒等, 这两个性质合起来给出 $Q^2=Q$.

\begin{example}\label{ex:modular}
设 $k$ 的特征为素数 $p$, 在 $k^2$ 上令 $C_p$ 的生成元作用为 $J=\begin{pmatrix}1&1\\0&1\end{pmatrix}$. 证明这确为表示, 但不完全可约.
\end{example}
\begin{solution}
令 $N=J-I$, 则 $N^2=0$, 由二项式公式 $J^p=I+pN=I$, 所以群关系成立. 直线 $U=ke_1$ 不变. 若存在不变补直线, 它由 $v=ae_1+be_2$, $b\ne0$ 生成. $Jv=(a+b)e_1+be_2$ 若是 $v$ 的标量倍, 比较 $e_2$ 系数得标量为1, 再比较 $e_1$ 系数得 $b=0$, 矛盾. 因此 $U$ 无不变补空间. 这也展示不可分解与不可约在一般底域中不同.
\end{solution}

\originalsection{不变内积}
\begin{proposition}\label{prop:unitary}
有限群的每个复表示可选取正定Hermite内积, 使所有群元素的作用为酉算子.
\end{proposition}
\begin{proof}
取任意正定Hermite内积 $(\ ,\ )_0$, 本书约定其第一变量线性. 定义
\[
 (v,w)_G=\frac1{|G|}\sum_{g\in G}(gv,gw)_0.
\]
这是Hermite形式. 非零 $v$ 的每个 $gv$ 非零, 因而 $(v,v)_G>0$. 对任意 $h$, 求和指标 $g\mapsto gh$ 是双射, 所以 $(hv,hw)_G=(v,w)_G$.
\end{proof}
如果 $U$ 不变, 则 $U^\perp$ 也不变, 因为对 $w\in U^\perp$, $u\in U$, 有 $(gw,u)_G=(w,g^{-1}u)_G=0$. 这给出Maschke定理在复数域上的第二个证明. 平均投影证明适用更多底域; 内积证明则对后面的矩阵系数正交非常方便.

\originalsection{分解的唯一性}
\begin{proposition}\label{prop:multiplicity}
若 $V$ 完全可约, $S$ 为不可约表示, 则任意不可约分解中 $S$ 的出现次数都等于 $\dim\Hom_G(S,V)$, 也等于 $\dim\Hom_G(V,S)$.
\end{proposition}
\begin{proof}
对有限直和, 线性映射由各分量唯一决定, 所以
$\Hom_G(S,\bigoplus_j V_j)\cong\bigoplus_j\Hom_G(S,V_j)$. Schur引理使每个与 $S$ 等价的分量贡献维数1, 其余贡献0. 反向Hom同样逐分量计算. 因而重数不依赖所选分解.
\end{proof}
这里的唯一性是指不可约同构类型与重数唯一, 而不是具体子空间唯一. 例如两个平凡表示之和是平凡作用的 $\C^2$, 任意两条不同直线都可作为分量.

\begin{proposition}
令 $V_S$ 为 $V$ 中所有与 $S$ 等价的不可约子表示之和. 则这些等型分量给出典范分解 $V=\bigoplus_S V_S$, 且
\[
 S\otimes\Hom_G(S,V)\longrightarrow V_S,\qquad s\otimes T\longmapsto T(s)
\]
是同构, 其中群只作用在第一个因子.
\end{proposition}
\begin{proof}
选一个不可约直和分解, 按同构类型把分量分组. 若 $S$ 与另一类型 $T$ 不等价, 任何 $S$ 子表示投影到 $T$ 分量的映射都为0; 所以所有 $S$ 子表示恰落在对应组内. 这证明组的空间就是 $V_S$, 与选择无关. 若 $V_S\cong S^{\oplus m}$, 选相应同构 $T_1,\ldots,T_m$ 作为 $\Hom_G(S,V)$ 的基, 上述评价映射就在各组上是该直和同构, 因而是双射.
\end{proof}
张量积将在第六章具体构造; 这里的 $S\otimes\C^m$ 只需理解为 $m$ 份 $S$ 的直和. 这一表述把不可约空间与重数空间分开, 解释了为什么重数是一个Hom空间的维数.

\begin{exercise}
若 $V\cong\bigoplus_i S_i^{\oplus m_i}$, 其中 $S_i$ 两两不等价, 求 $\End_G(V)$ 的结构与维数.
\end{exercise}
\begin{solution}
交织映射的各块 $S_i\to S_j$ 在 $i\ne j$ 时为0; 同型块为标量. 因此 $m_i$ 份 $S_i$ 之间的所有映射由一个 $m_i\times m_i$ 复矩阵描述, 复合对应矩阵乘法. 所以 $\End_G(V)\cong\bigoplus_i\operatorname{Mat}_{m_i}(\C)$, 维数是 $\sum_i m_i^2$.
\end{solution}

\originalchapter{特征标与正交关系}
\originalsection{特征标的基本性质}
\begin{definition}
表示 $V$ 的特征标是函数 $\chi_V:G\to\C$, $\chi_V(g)=\tr\rho_V(g)$. 在每个共轭类上取常值的函数称类函数; 所有复值类函数构成向量空间 $\Cl(G)$.
\end{definition}
迹不因换基而改变. 又 $\rho(hgh^{-1})=\rho(h)\rho(g)\rho(h)^{-1}$, 因而特征标为类函数. $\chi_V(e)=\dim V$, 直和满足 $\chi_{V\oplus W}=\chi_V+\chi_W$.

若 $g$ 的阶为 $m$, 则 $\rho(g)^m=I$. 多项式 $t^m-1$ 在复数域中没有重根, 所以 $\rho(g)$ 可对角化, 所有特征值都是 $m$ 次单位根. 因此
\[
 \chi_V(g^{-1})=\overline{\chi_V(g)},\qquad |\chi_V(g)|\le\dim V.
\]
后一不等式来自单位复数之和的三角不等式. 若 $\chi_V(g)=\dim V$, 各特征值实部至多为1而实部总和已等于维数, 所以它们全为1; 对角化再给出 $\rho(g)=I$. 于是
\begin{equation}\label{eq:character-kernel}
 \ker\rho_V=\{g\in G:\chi_V(g)=\chi_V(e)\}.
\end{equation}
特征标能辨别哪些群元素实际作用恒等, 但并不单凭一个值恢复整套矩阵.

\begin{example}
设 $G$ 作用在有限集合 $X$ 上, 在基 $\{e_x:x\in X\}$ 张成的空间 $\C[X]$ 上令 $g e_x=e_{gx}$. 求其特征标.
\end{example}
\begin{solution}
矩阵每列只含一个1. 对角线上对应 $x$ 的元素等于1恰在 $gx=x$ 时, 否则为0. 因此 $\chi_{\C[X]}(g)=|X^g|$, 其中 $X^g=\{x:gx=x\}$. 例题\ref{ex:permstandard}的三维置换表示在 $S_3$ 三类上的值是 $3,1,0$; 减去平凡分量的特征标1, 得到标准平面表示的值 $2,0,-1$.
\end{solution}

\originalsection{类函数内积与不变向量}
\begin{definition}
在复值函数空间上定义Hermite内积
\[
 \inner{f}{h}_G=\frac1{|G|}\sum_{g\in G}f(g)\overline{h(g)}.
\]
若两者为类函数, 选各类代表 $c_j$, 类大小为 $n_j$, 则内积等于 $|G|^{-1}\sum_j n_j f(c_j)\overline{h(c_j)}$.
\end{definition}
注意这是加权内积, 不是直接对特征标表每列求平均; 不同共轭类含的元素数不同.

\begin{proposition}\label{prop:reynolds}
对表示 $V$, 算子 $P_G=|G|^{-1}\sum_{g\in G}\rho_V(g)$ 是投影到不变向量空间 $V^G=\{v:gv=v\ \forall g\}$ 的投影. 因而
\[
 \dim V^G=\frac1{|G|}\sum_{g\in G}\chi_V(g)=\inner{\chi_V}{1}_G.
\]
\end{proposition}
\begin{proof}
对任意 $h$, $\rho(h)P_G=P_G$, 所以 $P_G(V)\subseteq V^G$; 对 $v\in V^G$, 所有求和项都是 $v$, 所以 $P_Gv=v$. 两者合起来给出 $P_G^2=P_G$, 且像正是 $V^G$. 在 $\im P_G\oplus\ker P_G$ 的基中矩阵为 $\operatorname{diag}(I,0)$, 迹等于像的维数. 迹的线性性给出所列等式.
\end{proof}

\begin{corollary}[轨道计数公式]\label{cor:burnside-count}
有限群 $G$ 作用在有限集合 $X$ 上, 轨道数为 $|G|^{-1}\sum_{g\in G}|X^g|$.
\end{corollary}
\begin{proof}
$\C[X]$ 中不变向量的系数在每个轨道上必须相同, 所以每个轨道贡献一个自由系数, 不变空间的维数就是轨道数. 应用命题\ref{prop:reynolds}与置换特征标公式即可.
\end{proof}
这是通常称为Burnside引理的计数公式, 与本书未纳入的Burnside $p^aq^b$ 可解性定理是两个不同结论.

\originalsection{交织空间的特征标计算}
我们希望把分解重数写成特征标的内积. 命题\ref{prop:multiplicity}已经把重数写为交织空间维数; 所以先把普通线性映射空间本身变成表示.

对 $T\in\Hom_\C(V,W)$, 定义 $g\cdot T=\rho_W(g)T\rho_V(g)^{-1}$. 这满足群作用法则, 而不变向量恰是交织映射. 若 $A,B$ 分别作用在 $W,V$, 线性算子 $T\mapsto ATB$ 在矩阵单位 $E_{ij}$ 基下的迹是 $\sum_{i,j}A_{ii}B_{jj}=\tr A\tr B$. 因此此Hom表示的特征标为 $\chi_W(g)\chi_V(g^{-1})$.

\begin{theorem}\label{thm:hom-character}
任意有限群复表示 $V,W$ 满足
\[
 \dim\Hom_G(V,W)=\inner{\chi_W}{\chi_V}_G.
\]
\end{theorem}
\begin{proof}
把命题\ref{prop:reynolds}用于上述Hom表示, 得
\[
 \dim\Hom_G(V,W)=\frac1{|G|}\sum_g\chi_W(g)\chi_V(g^{-1})
 =\frac1{|G|}\sum_g\chi_W(g)\overline{\chi_V(g)}.
\]
这里每个等式均发生在有限维空间与有限和中, 无须额外极限交换条件.
\end{proof}

\begin{corollary}[不可约特征标的行正交]\label{cor:roworth}
若 $S_i$ 两两不等价且不可约, 则 $\inner{\chi_{S_i}}{\chi_{S_j}}_G=\delta_{ij}$.
\end{corollary}
\begin{proof}
由定理\ref{thm:hom-character}, 内积是 $\dim\Hom_G(S_j,S_i)$, 再应用Schur引理.
\end{proof}
\begin{corollary}\label{cor:char-decomposition}
若 $V\cong\bigoplus_i S_i^{\oplus m_i}$, 则
\[
 m_i=\inner{\chi_V}{\chi_{S_i}}_G,\qquad
 \inner{\chi_V}{\chi_V}_G=\sum_i m_i^2.
\]
特别地, 非零表示不可约当且仅当其特征标的内积平方范数为1. 两个表示同构当且仅当特征标相同.
\end{corollary}
\begin{proof}
特征标的可加性与行正交给出两式. 正整数重数的平方和等于1恰表示只有一个分量且重数为1. 若特征标相同, 与每个不可约特征标取内积得到相同重数, 分解唯一性保证同构. 反向由迹在同构下不变得到.
\end{proof}
范数判别只适用于已知是表示特征标的函数. 任意类函数即使范数为1, 也可能不是表示特征标; 第六章会给虚表示版本的精确补充条件.

\originalsection{分解计算}
\begin{example}
$S_3$ 作用在集合 $X=\{(i,j):1\le i,j\le3,\ i\ne j\}$ 上. 分解对应的六维置换表示.
\end{example}
\begin{solution}
一个有序对被 $g$ 固定, 当且仅当其两个坐标都被固定. 单位元固定6个, 换位只固定一个数字, 因而没有互异固定对, 三循环无固定点. 特征标为 $(6,0,0)$. 三个候选不可约特征标为平凡 $(1,1,1)$、符号 $(1,-1,1)$ 和标准 $(2,0,-1)$; 标准表示的范数为 $(4+0+2)/6=1$, 所以不可约. 与六维表示取内积得重数 $1,1,2$. 三者维数之和为 $1+1+2\cdot2=6$, 所以本表示的分解为 $1\oplus\sgn\oplus U^{\oplus2}$. 第四章将证明该候选列表也穷尽了 $S_3$ 的全部不可约表示.
\end{solution}

\begin{exercise}
设置换表示 $\C[X]$ 中平凡表示出现 $a$ 次. 证明 $a$ 等于 $G$ 在 $X$ 上的轨道数. 若作用传递, 是否必有 $\C[X]\cong1\oplus$ 一个不可约表示?
\end{exercise}
\begin{solution}
平凡表示重数为 $\dim\C[X]^G$, 由轨道计数论证即为轨道数. 后一断言不成立: $C_3$ 在自身上的左乘作用传递, 其三维置换表示分解成三个不同一维表示. 其中平凡表示确只出现一次, 但剩余二维空间仍可约.
\end{solution}

\originalchapter{正则表示与群代数结构}
\originalsection{正则表示与平方和}
群代数 $R=\C[G]$ 上的左乘作用 $h\cdot[g]=[hg]$ 称左正则表示. 它也是 $G$ 在自身上左乘的置换表示, 所以
\[
 \chi_R(g)=\begin{cases}|G|,&g=e,\\0,&g\ne e.\end{cases}
\]
非单位元左乘不可能固定群元素, 这解释了所有非单位元处的零值.

\begin{theorem}\label{thm:regular}
有限群只有有限多个复不可约同构类型 $S_1,\ldots,S_r$. 令 $d_i=\dim S_i$, 则
\[
 R\cong\bigoplus_{i=1}^r S_i^{\oplus d_i},\qquad |G|=\sum_{i=1}^r d_i^2.
\]
\end{theorem}
\begin{proof}
Maschke定理使有限维 $R$ 分解为有限个不可约分量. 对任意不可约 $S$, 定理\ref{thm:hom-character}给出
\[
 \dim\Hom_G(S,R)=\inner{\chi_R}{\chi_S}_G=\overline{\chi_S(e)}=\dim S>0.
\]
因此 $S$ 一定同构于 $R$ 的某个分量, 所有不可约类型都已在这个有限列表中. 同一公式给出 $S_i$ 在 $R$ 中的重数为 $d_i$. 再比较总维数得到平方和.
\end{proof}
平方和是检验候选列表是否完整的有效办法: 若已构造一批两两不等价的不可约, 其维数平方和等于 $|G|$, 就不可能再有新不可约类型. 只看到平方和正确, 却未验证候选不可约与互不等价, 还不能宣布分类完成.

\originalsection{矩阵分块与共轭类数}
\begin{theorem}\label{thm:groupalgebra}
映射
\[
 \Phi:\C[G]\longrightarrow\bigoplus_{i=1}^r\End_\C(S_i),\qquad
 a\longmapsto(\rho_i(a))_{i=1}^r
\]
是保单位元的代数同构.
\end{theorem}
\begin{proof}
各 $\rho_i$ 都为代数同态, 所以 $\Phi$ 保乘法与单位元. 若 $\Phi(a)=0$, 则 $a$ 在所有不可约上作用为零, 因此在它们组成的正则表示上也作用为零. 可是正则作用中 $a\cdot e=a$, 所以 $a=0$. 这证明单射. 两边维数分别为 $|G|$ 与 $\sum_i d_i^2$, 定理\ref{thm:regular}说明相等, 因而单射亦为满射.
\end{proof}
这里不调用一般有限维代数的结构定理. 完全可约性与正则表示已经提供证明所需的全部信息. 该同构到抽象 $\End(S_i)$ 不需选基; 把各块写成具体矩阵代数时才需选择 $S_i$ 的基. 分块大小 $d_i$ 是内在的.

\begin{theorem}\label{thm:classes}
不可约表示的同构类型数等于共轭类数. 不可约特征标构成 $\Cl(G)$ 的正交规范基.
\end{theorem}
\begin{proof}
对 $a=\sum_g a_g g$, 中心条件等价于对所有 $h$ 有 $hah^{-1}=a$, 即系数在共轭类上相同. 因此每个类的元素和构成群代数中心的一组基, 中心维数为类数.

矩阵代数 $\operatorname{Mat}_d(\C)$ 的中心只有标量矩阵: 与所有对角矩阵交换迫使非对角项为零, 再与每个矩阵单位 $E_{ij}$ 交换迫使所有对角项相同. 直和代数的中心是各块中心的直和, 所以由定理\ref{thm:groupalgebra}, 群代数中心维数为 $r$. 两种计算得类数等于 $r$. 类函数空间的维数也等于类数, 而 $r$ 个不可约特征标已由推论\ref{cor:roworth}证明正交规范, 故构成基.
\end{proof}

\originalsection{列正交与矩阵系数}
\begin{corollary}[特征标列正交]\label{cor:columnorth}
若 $g,h$ 不共轭, 则 $\sum_i\chi_i(g)\overline{\chi_i(h)}=0$. 若 $g,h$ 共轭, 此和为 $|C_G(g)|$.
\end{corollary}
\begin{proof}
令类代表为 $c_1,\ldots,c_r$, 类大小为 $n_j$. 矩阵 $A_{ij}=\sqrt{n_j/|G|}\chi_i(c_j)$ 由行正交满足 $AA^*=I$. 它是方阵, 所以 $A$ 可逆且 $A^*=A^{-1}$, 进而 $A^*A=I$. 列 $j,k$ 的关系给出 $j\ne k$ 时和为0, $j=k$ 时和为 $|G|/n_j=|C_G(c_j)|$. 同类的特征标值相同, 因而得到一般形式.
\end{proof}
行正交确定分解重数, 列正交则约束不同共轭类的值. 在计算特征标表时, 两者结合常能恢复缺失项, 但仍要构造或论证这些行确实来自表示.

\begin{theorem}[矩阵系数正交]\label{thm:matrixorth}
对每个不可约 $S_i$ 取不变内积的正交规范基, 写 $\rho_i(g)_{ab}$ 为矩阵元素. 则
\[
 \frac1{|G|}\sum_{g\in G}\rho_i(g)_{ab}\overline{\rho_j(g)_{cd}}
 =\begin{cases}d_i^{-1}\delta_{ac}\delta_{bd},&i=j,\\0,&i\ne j.\end{cases}
\]
所以所有函数 $\sqrt{d_i}\rho_i(\cdot)_{ab}$ 构成 $G$ 上全部复值函数空间的正交规范基.
\end{theorem}
\begin{proof}
对 $T:S_j\to S_i$, 平均算子
$\mathcal P(T)=|G|^{-1}\sum_g\rho_i(g)T\rho_j(g)^{-1}$ 为交织映射. 若 $i\ne j$, Schur引理给 $\mathcal P(T)=0$. 若 $i=j$, 其为标量算子, 且平均不改变迹, 所以 $\mathcal P(T)=(\tr T/d_i)I$.

令 $T$ 为只有第 $(b,d)$ 项为1的矩阵单位. 因各矩阵酉, $\rho_j(g)^{-1}_{dc}=\overline{\rho_j(g)_{cd}}$, 故 $\mathcal P(T)$ 的 $(a,c)$ 项就是所求平均. 当 $i=j$ 时, $\tr T=\delta_{bd}$; 这给出公式. 正交规范函数的总数为 $\sum_i d_i^2=|G|$, 与全部函数空间的维数相同, 因而是一组基.
\end{proof}

\originalsection{中心投影与傅里叶反演}
\begin{proposition}\label{prop:central-idempotent}
定义
\[
 e_i=\frac{d_i}{|G|}\sum_{g\in G}\chi_i(g^{-1})g\in\C[G].
\]
则 $e_i$ 在 $S_j$ 上作用为 $\delta_{ij}I$, 并满足 $e_i^2=e_i$, $e_ie_j=0$ ($i\ne j$), $\sum_i e_i=e$. 在任意表示 $V$ 上, $\rho_V(e_i)$ 是到 $S_i$ 等型分量的投影.
\end{proposition}
\begin{proof}
系数为类函数, 所以 $e_i$ 在群代数中心, 在 $S_j$ 上为标量 $\lambda_{ij}I$. 取迹得到
\[
 d_j\lambda_{ij}=\frac{d_i}{|G|}\sum_g\chi_i(g^{-1})\chi_j(g)
 =d_i\delta_{ij}.
\]
当 $i=j$ 时 $d_i=d_j$, 从而 $\lambda_{ij}=\delta_{ij}$. 各乘积与总和在每个矩阵块上分别满足所述关系, 定理\ref{thm:groupalgebra}的单射性把关系拉回群代数. 最后按 $V$ 的不可约分解检查作用, 就得到等型投影.
\end{proof}

\begin{proposition}
若 $a=\sum_g a_g g$, 则
\[
 a_g=\frac1{|G|}\sum_i d_i\tr\bigl(\rho_i(g^{-1})\rho_i(a)\bigr).
\]
\end{proposition}
\begin{proof}
对任意 $b\in\C[G]$, 左乘 $b$ 在正则表示上的迹为 $|G|b_e$, 因为每个对角元素都为单位元系数 $b_e$. 又正则表示包含 $d_i$ 份 $S_i$, 所以该迹为 $\sum_i d_i\tr\rho_i(b)$. 取 $b=g^{-1}a$, 其单位元系数为 $a_g$, 即得公式.
\end{proof}
这是有限非交换群上的傅里叶反演: 标量系数 $a_g$ 与一族矩阵 $\rho_i(a)$ 相互决定. 对交换群, 所有块一维, 就回到通常的离散傅里叶变换.

\originalsection{Frobenius行列式}
\begin{theorem}\label{thm:frobdet}
为每个 $g\in G$ 引入独立变量 $x_g$, 将行、列都按群元素排列. 矩阵 $X$ 的 $(g,h)$ 项为 $x_{gh^{-1}}$. 则
\[
 \det X=\prod_i\det\left(\sum_{g\in G}x_g\rho_i(g)\right)^{d_i}.
\]
\end{theorem}
\begin{proof}
令 $a=\sum_g x_g g$, 在多项式系数环上考虑正则表示. $a h=\sum_g x_g gh$, 输出基元素 $u$ 的系数为 $x_{uh^{-1}}$, 所以左乘矩阵就是 $X$. 正则表示的复数域分解使用一个固定的常系数换基矩阵, 因而在加入变量后仍成立. 每个 $S_i$ 块出现 $d_i$ 次, 分块对角矩阵的行列式就是右侧乘积.
\end{proof}
该公式说明因子次数和出现重数都由不可约维数控制. 本节证明行列式的表示论分解, 不另外证明这些多项式因子在多项式环中的不可约性.

\begin{exercise}
对 $C_2=\{e,s\}$ 直接核对Frobenius行列式及中心投影.
\end{exercise}
\begin{solution}
两个一维表示在 $s$ 上取 $1,-1$. 矩阵 $X=\begin{pmatrix}x_e&x_s\\x_s&x_e\end{pmatrix}$ 的行列式是 $(x_e+x_s)(x_e-x_s)$, 与两块对应. 中心投影为 $e_+=(e+s)/2$, $e_-=(e-s)/2$; 用 $s^2=e$ 直接得 $e_\pm^2=e_\pm$, $e_+e_-=0$, $e_++e_-=e$.
\end{solution}

\originalchapter{经典群的特征标表}
\originalsection{循环群与一维表示}
由推论\ref{cor:abelian}, $C_n$ 的不可约复表示都是一维, 第一章构造的 $L_0,\ldots,L_{n-1}$ 已穷尽所有可能. 特征标表是 $\chi_j(r^a)=\zeta^{ja}$. 正交关系也可直接由几何级数核对: 当 $j\ne k\pmod n$ 时,
\[
 \sum_{a=0}^{n-1}\zeta^{(j-k)a}=\frac{1-\zeta^{(j-k)n}}{1-\zeta^{j-k}}=0;
\]
当 $j=k$ 时和为 $n$. 这说明特征标正交确实推广了离散傅里叶基的正交.

对一般群, 一维表示是同态 $G\to\C^\times$. 由于目标群交换, 每个交换子 $ghg^{-1}h^{-1}$ 必映成1. 令 $[G,G]$ 为交换子生成的正规子群, 则所有一维表示都经交换化 $G/[G,G]$ 分解. 反过来, 交换化的每个一维表示拉回后仍是一维表示. 因此一维表示只看见群的交换部分, 更高维表示记录其余对称结构.

\originalsection{三次对称群}
$S_3$ 的三类已在例题\ref{ex:s3classes}中求出. 平凡表示、置换矩阵行列式给出的符号表示以及标准平面表示 $U$ 的特征标如下.
\begin{center}
\begin{tabular}{c|rrr}
\toprule
代表元 &$e$&$(12)$&$(123)$\\
类大小 &1&3&2\\\midrule
$1$&1&1&1\\
$\sgn$&1&$-1$&1\\
$U$&2&0&$-1$\\\bottomrule
\end{tabular}
\end{center}
符号表示为同态, 因为置换矩阵的行列式在乘积下相乘. $U$ 的范数已算为1, 所以不可约. 三行不同, 维数平方和 $1+1+4=6$ 穷尽全部不可约. 特征标表既可由具体矩阵计算, 又可用正交关系检查; 两种做法彼此支持.

\begin{example}
分解 $U\otimes U$, 并求 $S_3$ 正则表示中 $U$ 的重数.
\end{example}
\begin{solution}
张量积特征标逐点相乘, 其证明见第六章. 所以 $U\otimes U$ 的特征标为 $(4,0,1)$. 与三行取加权内积分别得 $(4+2)/6=1$, $(4+2)/6=1$, $(8-2)/6=1$. 因此 $U\otimes U\cong1\oplus\sgn\oplus U$. 正则表示的重数等于不可约维数, 所以 $U$ 出现两次.
\end{solution}

\originalsection{正方形的二面体群}
记 $D_8=\langle r,s:r^4=s^2=e,\ srs=r^{-1}\rangle$, 阶为8; 本书的下标表示群的阶. 可把 $r$ 看成正方形旋转90度, $s$ 看成反射. 元素为 $r^a$ 与 $r^as$, $a=0,1,2,3$. 关系给出五个共轭类:
\[
 \{e\},\quad\{r^2\},\quad\{r,r^3\},\quad\{s,r^2s\},\quad\{rs,r^3s\}.
\]
例如 $r(r^as)r^{-1}=r^{a+2}s$, 所以反射按指数奇偶分成两类. $r^2$ 在中心.

一维表示中 $r\mapsto a$, $s\mapsto b$. 关系 $srs=r^{-1}$ 给 $a=a^{-1}$, 所以 $a,b\in\{1,-1\}$; 四种选择都成立, 记为 $L_{a,b}$. 二维几何表示 $E$ 给出矩阵
\[
 \rho_E(r)=\begin{pmatrix}0&-1\\1&0\end{pmatrix},\qquad
 \rho_E(s)=\begin{pmatrix}1&0\\0&-1\end{pmatrix}.
\]
其特征标表可合并写成
\begin{center}
\begin{tabular}{c|rrrrr}
\toprule
代表元 &$e$&$r^2$&$r$&$s$&$rs$\\
类大小 &1&1&2&2&2\\\midrule
$L_{a,b}$&1&1&$a$&$b$&$ab$\\
$E$&2&$-2$&0&0&0\\\bottomrule
\end{tabular}
\end{center}
$E$ 的范数为 $(4+4)/8=1$, 所以不可约; 平方和 $4\cdot1^2+2^2=8$, 分类完整. 这说明二维实几何空间在复数域中仍不可约. 对照仅旋转的循环群, 旋转矩阵可分成两个复特征直线, 但加入反射后两条直线被交换, 不再分别为子表示.

\begin{center}
\begin{tikzpicture}[scale=.8,>=Stealth]
\draw (-1,-1) rectangle (1,1);
\foreach \p/\t in {(1,1)/1,(-1,1)/2,(-1,-1)/3,(1,-1)/4}{\fill \p circle (2pt);}
\node[above right] at (1,1){1};\node[above left] at (-1,1){2};
\node[below left] at (-1,-1){3};\node[below right] at (1,-1){4};
\draw[dashed] (-1.6,0)--(1.6,0);\node[left] at (-1.6,0){反射轴};
\draw[->] (1.7,0) arc[start angle=0,end angle=75,radius=1.7];\node[right] at (1.7,.7){$r$};
\end{tikzpicture}
\end{center}
图中 $r$ 取逆时只改变生成元选择, 不改变特征标表; 这里的矩阵约定为逆时针90度旋转.

\originalsection{四元数群}
四元数群 $Q_8=\{\pm1,\pm i,\pm j,\pm k\}$ 满足 $i^2=j^2=k^2=-1$, $ij=k=-ji$. 其五个共轭类依次为 $\{1\}$、$\{-1\}$、$\{\pm i\}$、$\{\pm j\}$ 与 $\{\pm k\}$. 中心为 $\{\pm1\}$, 交换化为 $C_2\times C_2$, 所以有四个一维表示, 由 $i\mapsto a,j\mapsto b$, $a,b=\pm1$ 决定; $-1$ 必映到1, $k$ 映到 $ab$.

令 $i$ 与 $j$ 在 $\C^2$ 上作用为
\[
 I=\begin{pmatrix}i&0\\0&-i\end{pmatrix},\qquad
 J=\begin{pmatrix}0&1\\-1&0\end{pmatrix}.
\]
矩阵满足 $I^2=J^2=-I_2$, $IJ=-JI$, 所以构成 $Q_8$ 表示 $F$. 其特征标如下.
\begin{center}
\begin{tabular}{c|rrrrr}
\toprule
代表元 &1&$-1$&$i$&$j$&$k$\\
类大小 &1&1&2&2&2\\\midrule
$L_{a,b}$&1&1&$a$&$b$&$ab$\\
$F$&2&$-2$&0&0&0\\\bottomrule
\end{tabular}
\end{center}
同样由范数1与平方和8得完整性. 作为带有列大小的数值表, 它与 $D_8$ 的表相同; 可是群并不同构, 因为 $Q_8$ 中只有一个二阶元素, $D_8$ 有五个. 特征标表并不总能恢复群. 第九章将用包含 $g^2$ 信息的Frobenius--Schur指标区别二者的二维表示.

\originalsection{四次交错群}
$A_4$ 为 $S_4$ 中的偶置换, 阶为12. 令
$N=\{e,(12)(34),(13)(24),(14)(23)\}$, 它是正规四元子群, 商群为 $C_3$. 三个非单位双换位构成一个共轭类. 八个三循环在 $A_4$ 中分成两个大小4的类: 例如 $(123)$ 的中心化子在 $A_4$ 中阶为3, 所以类大小为4; 其逆不在同一类, 因为把一个三循环共轭成逆的 $S_4$ 置换为奇置换. 选两类代表 $t,t^{-1}$, 使 $tN$ 生成商群.

商群给出三个一维表示 $1,L,L^2$, 在 $t$ 上取 $1,\omega,\omega^2$, 其中 $\omega=\exp(2\pi i/3)$. 四点置换表示的常数向量补空间 $W$ 为三维. 固定点计数减1给出特征标 $3,-1,0,0$, 所以
\begin{center}
\begin{tabular}{c|rrrr}
\toprule
代表元 &$e$&$(12)(34)$&$t$&$t^{-1}$\\
类大小 &1&3&4&4\\\midrule
$1$&1&1&1&1\\
$L$&1&1&$\omega$&$\omega^2$\\
$L^2$&1&1&$\omega^2$&$\omega$\\
$W$&3&$-1$&0&0\\\bottomrule
\end{tabular}
\end{center}
$W$ 范数为 $(9+3)/12=1$, 所以不可约. 平方和 $1+1+1+9=12$, 分类完整. 三维表示具有实矩阵, 但两个非平凡一维表示互为复共轭, 这说明特征标的值可能不是实数.

\begin{exercise}
分解 $A_4$ 的 $W\otimes W$, 并由特征标检查所得总维数.
\end{exercise}
\begin{solution}
张量特征标为 $(9,1,0,0)$. 与三个一维特征标内积都为 $(9+3)/12=1$, 与 $W$ 内积为 $(27-3)/12=2$. 故 $W\otimes W\cong1\oplus L\oplus L^2\oplus W^{\oplus2}$, 维数 $1+1+1+2\cdot3=9$.
\end{solution}

\originalchapter{对偶、张量积与新表示}
\originalsection{对偶表示}
$V^*=\Hom_\C(V,\C)$ 的元素是线性泛函. 要使自然配对 $\lambda(v)$ 在同时作用后不变, 必须定义
\[
 (g\lambda)(v)=\lambda(g^{-1}v).
\]
这满足 $(gh)\lambda=g(h\lambda)$, 在对偶基中的矩阵为 $\rho_{V^*}(g)=\rho_V(g^{-1})^{\mathsf T}$, 所以
\[
 \chi_{V^*}(g)=\chi_V(g^{-1})=\overline{\chi_V(g)}.
\]
若 $V$ 不可约, $V^*$ 也不可约: 由对偶特征标与原特征标的平方范数相同, 推论\ref{cor:char-decomposition}即可判断. 也可取非零真子表示的零化子获得反证.

有自然线性同构 $V^*\otimes W\to\Hom_\C(V,W)$, 把 $\lambda\otimes w$ 映到 $v\mapsto\lambda(v)w$. 在基与对偶基中这些映射就是矩阵单位, 因而构成一组基; 对群作用的兼容性由上述逆元公式保证. 这解释第三章Hom特征标乘积公式的结构来源.

\originalsection{张量积与重数}
\begin{definition}
$V\otimes W$ 是由符号 $v\otimes w$ 张成、满足对每一变量的线性关系的向量空间. 若 $\{v_i\},\{w_j\}$ 为基, 则 $\{v_i\otimes w_j\}$ 为基, 维数为 $\dim V\dim W$. 两个 $G$ 表示的张量积作用定义为 $g(v\otimes w)=gv\otimes gw$.
\end{definition}
基结论可通过坐标直接核实: 双线性关系使任意纯张量写成 $\sum_{i,j}a_i b_j v_i\otimes w_j$; 反向把这些基符号映到独立坐标向量, 给出线性独立性. 因而张量积不是只由纯张量组成的集合, 一般元素是纯张量的和.

\begin{proposition}\label{prop:tensorcharacter}
$\chi_{V\otimes W}(g)=\chi_V(g)\chi_W(g)$. 若 $S_i$ 为全部不可约, 则 $S_k$ 在 $V\otimes W$ 中的重数为
\[
 \frac1{|G|}\sum_{g\in G}\chi_V(g)\chi_W(g)\overline{\chi_k(g)}.
\]
\end{proposition}
\begin{proof}
若 $g$ 在两空间的矩阵为 $A,B$, 在张量基的对角元素为 $A_{ii}B_{jj}$, 总迹为 $\sum_{i,j}A_{ii}B_{jj}=\tr A\tr B$. 重数公式再由特征标分解定理得到.
\end{proof}

张量积可以产生新表示, 也可以暴露已有表示的对偶关系. 例如 $1$ 在 $V^*\otimes W$ 中的重数为 $\dim\Hom_G(V,W)$. 当 $V,W$ 不可约时, 这个数恰为0或1, 而不是张量积本身通常不可约.

\originalsection{对称平方与外平方}
在 $V\otimes V$ 上定义交换算子 $\tau(v\otimes w)=w\otimes v$. 它满足 $\tau^2=I$ 并与群作用交换. 因为底域为复数, 投影 $(I\pm\tau)/2$ 给出直和
\[
 V\otimes V=\Sym^2V\oplus\bigwedge^2V,
\]
两者分别为 $\tau$ 的 $+1,-1$ 特征空间. 在基中, 对称部分由 $v_i\otimes v_i$ 与 $v_i\otimes v_j+v_j\otimes v_i$ ($i<j$) 张成, 外平方由 $v_i\otimes v_j-v_j\otimes v_i$ ($i<j$) 张成. 维数分别为 $d(d+1)/2$ 与 $d(d-1)/2$.

\begin{proposition}\label{prop:symext}
对任意 $g\in G$, 有
\[
 \chi_{\Sym^2V}(g)=\frac{\chi_V(g)^2+\chi_V(g^2)}2,\qquad
 \chi_{\bigwedge^2V}(g)=\frac{\chi_V(g)^2-\chi_V(g^2)}2.
\]
\end{proposition}
\begin{proof}
令 $\rho_V(g)$ 的特征值为 $\lambda_1,\ldots,\lambda_d$, 对称平方的特征值是 $\lambda_i\lambda_j$ ($i\le j$), 外平方是同样乘积但只取 $i<j$. 平方 $(\sum_i\lambda_i)^2$ 把所有 $i\ne j$ 项计两次, 而 $\sum_i\lambda_i^2=\chi_V(g^2)$ 只计对角项. 相加或相减再除以2, 得到两式.
\end{proof}

\begin{example}
对 $S_3$ 标准表示 $U$, 求 $\Sym^2U$ 与 $\bigwedge^2U$.
\end{example}
\begin{solution}
三类上 $\chi_U=(2,0,-1)$, 平方映射把换位送到单位元、三循环送到另一三循环, 所以 $\chi_U(g^2)=(2,2,-1)$. 两公式给出对称平方特征标 $(3,1,0)$, 外平方特征标 $(1,-1,1)$. 前者是 $1\oplus U$, 后者是符号表示. 也可由 $\bigwedge^2U$ 的作用就是二维矩阵行列式直接看出后一结论.
\end{solution}

\originalsection{直积群}
若 $V$ 是 $G$ 表示, $W$ 是 $H$ 表示, 在 $V\otimes W$ 上定义 $(g,h)(v\otimes w)=gv\otimes hw$, 记为外张量积 $V\boxtimes W$. 此处群的两个因子独立作用, 与同一个群的内部张量积有所区别.

\begin{theorem}\label{thm:products}
若 $S_i,T_j$ 分别为 $G,H$ 的全部复不可约表示, 则 $S_i\boxtimes T_j$ 两两不等价且不可约, 并穷尽 $G\times H$ 的全部不可约.
\end{theorem}
\begin{proof}
特征标为 $\chi_i(g)\psi_j(h)$. 在直积群上的内积是两重有限和, 可分成因子的内积乘积:
\[
 \inner{\chi_i\psi_j}{\chi_k\psi_\ell}_{G\times H}
 =\inner{\chi_i}{\chi_k}_G\inner{\psi_j}{\psi_\ell}_H
 =\delta_{ik}\delta_{j\ell}.
\]
由范数判别知不可约, 由不同特征标正交知互不等价. 又维数平方和为
$\sum_{i,j}(d_i\dim T_j)^2=(\sum_i d_i^2)(\sum_j(\dim T_j)^2)=|G||H|$, 所以没有遗漏.
\end{proof}

\originalsection{虚表示}
\begin{definition}
不可约表示的整数线性组合 $\sum_i n_i S_i$ ($n_i\in\Z$) 称虚表示, 特征标为 $\sum_i n_i\chi_i$. 实际表示对应所有 $n_i\ge0$ 的情形.
\end{definition}
虚表示允许在构造中先减去已知分量. 例如置换表示减去平凡表示是合法虚表示; 若存在不变补空间, 它又实际由该补空间实现. 不能只凭某个函数在单位元处为正, 就把任意差写成实际表示.

\begin{lemma}\label{lem:virtualnorm}
若虚特征标 $f$ 满足 $\inner f f_G=1$ 且 $f(e)>0$, 则 $f$ 是某个不可约表示的特征标.
\end{lemma}
\begin{proof}
写 $f=\sum_i n_i\chi_i$, 行正交给 $\inner f f_G=\sum_i n_i^2$. 整数平方和为1迫使恰有一个系数为 $1$ 或 $-1$, 其他全为0. 此时 $f(e)=\pm d_i$, 正值条件选出 $+1$.
\end{proof}
“虚”这一条件不可省略: 它保证系数为整数, 才能从平方和为1推出单一系数. 期末作业中会利用这个引理构造二元二十面体群的一个不可约表示.

\begin{exercise}
证明用任意一维表示 $L$ 张量, 会置换不可约表示的同构类型. 说明其逆操作.
\end{exercise}
\begin{solution}
$L$ 的每个群矩阵为非零标量. $L\otimes V$ 的子空间不变条件因此与 $V$ 相同, 所以保持不可约性. 若 $L\otimes V\cong L\otimes W$, 再张量 $L^*$ 并用 $L^*\otimes L\cong1$, 得 $V\cong W$. 逆操作正是张量 $L^*$, 从而给出同构类型集合的置换.
\end{solution}

\Needspace{9\baselineskip}
\originalchapter{限制与诱导}
\originalsection{限制与构造问题}
若 $H\le G$, 在同一个空间上仅让 $H$ 中元素作用, 得限制表示 $\Res_H^G V$. 一个 $G$ 不可约表示限制后可以分裂: $S_3$ 的标准表示限制到 $C_3=\langle(123)\rangle$ 时, 三循环矩阵有两个不同特征值 $\omega,\omega^2$, 所以分解为 $L_1\oplus L_2$. 反过来, 给定 $H$ 表示 $W$, 怎样系统构造 $G$ 表示? 不能随意给 $G\setminus H$ 的元素指定矩阵, 因为全部群关系必须相容. 诱导表示通过为每个陪集安放一份 $W$ 来解决这个问题.

\originalsection{群代数张量构造}
$\C[G]$ 可在左边受 $G$ 左乘、右边受 $H$ 右乘. 定义平衡张量积
\[
 \Ind_H^G W=\C[G]\otimes_{\C[H]}W,
\]
它是普通 $\C[G]\otimes_\C W$ 的商空间, 增添关系 $(ah)\otimes w=a\otimes hw$ ($a\in\C[G],h\in H$). 这里下标表示可以把 $H$ 的作用从第一因子移到第二因子, 而不只是复标量之间的移动. $G$ 的作用为 $g(a\otimes w)=ga\otimes w$; 左乘与右乘交换, 所以作用保持平衡关系, 在商空间上良定义.

选左陪集代表 $T\subset G$, 每个元素唯一写成 $th$, $t\in T,h\in H$. 因而群代数作为右 $\C[H]$ 模是 $\bigoplus_{t\in T}t\C[H]$, 诱导空间相应为
\begin{equation}\label{eq:ind-cosets}
 \Ind_H^G W=\bigoplus_{t\in T}t\otimes W,\qquad
 \dim\Ind_H^G W=[G:H]\dim W.
\end{equation}
要看这不是仅有生成性, 可将 $th\otimes w$ 映到第 $t$ 份中的 $hw$. 此映射保持全部平衡关系, 与 $w\mapsto t\otimes w$ 的直和互逆, 所以没有隐藏的额外关系.

若 $g t=t'h$ ($t'\in T,h\in H$), 则作用公式为
\[
 g(t\otimes w)=t'\otimes hw.
\]
所以群作用置换陪集块, 同时在每个块内应用一个 $H$ 矩阵. 这个公式可直接生成诱导表示的具体矩阵.

\originalsection{等变函数模型}
另一种模型为所有函数 $f:G\to W$, 满足 $f(xh)=h^{-1}f(x)$, 群作用为 $(g f)(x)=f(g^{-1}x)$. 条件沿右乘 $H$ 的方向决定函数值, 所以每个左陪集的一个值决定整个函数. 左平移保持此条件, 且满足表示法则.

把 $t\otimes w$ 映成支撑在 $tH$ 上的函数 $f_{t,w}$, 在 $th$ 处取值 $h^{-1}w$. 若 $t$ 改为 $th_0$, $w$ 改为 $h_0^{-1}w$, 所得函数不变; 这正对应平衡张量关系. 逐陪集的取值映射证明它是同构, 且 $g f_{t,w}=f_{gt,w}$ 保证交织. 因而这两种构造描述同一个诱导表示.

\begin{remark}
旧讲义§4.8使用条件 $F(hx)=hF(x)$ 与作用 $(gF)(x)=F(xg)$, 即右陪集函数模型. 令 $F(x)=f(x^{-1})$ 就得到该约定. 因此旧源公式中的 $xgx^{-1}$ 与本书公式中的 $x^{-1}gx$ 不冲突, 只是代表元方向不同. 用公式计算时不能混用两个约定.
\end{remark}

\originalsection{诱导特征标}
\begin{theorem}\label{thm:ind-char}
设 $\chi_W$ 为 $H$ 表示 $W$ 的特征标, $T$ 为左陪集代表. 则
\begin{align*}
 \chi_{\Ind_H^G W}(g)
 &=\sum_{\substack{t\in T\\t^{-1}gt\in H}}\chi_W(t^{-1}gt)\\
 &=\frac1{|H|}\sum_{\substack{x\in G\\x^{-1}gx\in H}}\chi_W(x^{-1}gx).
\end{align*}
\end{theorem}
\begin{proof}
依照式\eqref{eq:ind-cosets}写作用矩阵. 若 $gtH\ne tH$, 原块 $t\otimes W$ 被送到另一块, 对总迹没有贡献. 若 $gtH=tH$, 令 $h=t^{-1}gt\in H$, 则 $g(t\otimes w)=t\otimes hw$, 这一对角块的迹为 $\chi_W(h)$. 相加得到第一式. 同一陪集代表换成 $th_0$ 时, $h$ 变成 $h_0^{-1}hh_0$, 特征标不变; 每个陪集有 $|H|$ 个元素, 从而得到第二式.
\end{proof}
若 $g$ 的共轭类不与 $H$ 相交, 诱导特征标在 $g$ 处为0. 在单位元处公式给 $[G:H]\dim W$, 与空间维数一致. 若 $W=1_H$, 诱导表示就是 $G$ 在 $G/H$ 上的置换表示, 特征标正是固定陪集数.

\originalsection{Frobenius互反}
\begin{theorem}[Frobenius互反]\label{thm:reciprocity}
若 $W$ 为 $H$ 表示, $V$ 为 $G$ 表示, 则有自然线性同构
\[
 \Hom_G(\Ind_H^G W,V)\cong\Hom_H(W,\Res_H^G V).
\]
因此在特征标层面,
\[
 \inner{\chi_{\Ind_H^G W}}{\chi_V}_G
 =\inner{\chi_W}{\chi_{\Res_H^G V}}_H.
\]
\end{theorem}
\begin{proof}
对 $G$ 交织映射 $T$, 定义 $\alpha(w)=T(e\otimes w)$. 由 $e\otimes hw=h\otimes w=h(e\otimes w)$ 得 $\alpha(hw)=h\alpha(w)$, 所以 $\alpha$ 为 $H$ 交织映射.

反过来, 对 $H$ 交织映射 $\alpha$, 定义 $T_\alpha(g\otimes w)=g\alpha(w)$, 再线性延拓. 需核对平衡关系: $T_\alpha(gh\otimes w)=gh\alpha(w)=g\alpha(hw)=T_\alpha(g\otimes hw)$, 所以映射在商空间上良定义. 又 $T_\alpha(a(g\otimes w))=ag\alpha(w)=aT_\alpha(g\otimes w)$, 因而交织. 在 $e\otimes w$ 上两构造互逆; 这类向量的全部 $G$ 平移生成诱导空间, 所以确是互逆同构.

取维数得Hom版本的数值等式. 定理\ref{thm:hom-character}把相应维数写成 $\inner{\chi_V}{\chi_{\Ind W}}_G$ 与 $\inner{\chi_{\Res V}}{\chi_W}_H$. 两边是实非负整数, 共轭不变, 所以也得到陈述中顺序的特征标等式.
\end{proof}
“自然”指这两个公式无需选基或陪集代表, 并与交织映射的前后复合相容. 互反把一个较大群上的分解计算变为子群上的限制计算, 是诱导表示最常用的工具.

对有限群还可证明另一个方向
$\Hom_G(V,\Ind_H^G W)\cong\Hom_H(\Res_H^G V,W)$. 在函数模型中, $T\mapsto(v\mapsto T(v)(e))$ 是所需映射. 反向公式为 $T_\alpha(v)(x)=\alpha(x^{-1}v)$. 因 $\alpha$ 交织,
$T_\alpha(v)(xh)=h^{-1}T_\alpha(v)(x)$; 又
$T_\alpha(gv)(x)=\alpha(x^{-1}gv)=T_\alpha(v)(g^{-1}x)$, 所以得到良定义交织映射. 取单位元值和利用交织性即可逐项核对互逆. 有限性保证函数模型与上节的诱导空间同构; 在无限群情形, 两种方向不能不加条件地混称同一个构造.

\originalsection{诱导的传递性与投影公式}
\begin{proposition}\label{prop:transitivity}
若 $K\le H\le G$, 则 $\Ind_H^G(\Ind_K^H W)\cong\Ind_K^G W$.
\end{proposition}
\begin{proof}
定义 $g\otimes(h\otimes w)\mapsto gh\otimes w$. 将中间 $H$ 元素从外因子移到内因子, 或将 $K$ 元素移到 $w$, 都不改变输出的平衡张量, 所以良定义且交织. 逆映射为 $g\otimes w\mapsto g\otimes(e\otimes w)$. 对 $k\in K$, 两种写法 $gk\otimes w$ 与 $g\otimes kw$ 的像相同, 因为可把 $k$ 先移入 $\C[H]$ 再移到 $W$. 两个复合在纯张量上为恒等, 所以是同构.
\end{proof}
\begin{proposition}\label{prop:projectionformula}
$G$ 表示 $V$ 与 $H$ 表示 $W$ 满足
\[
 \Ind_H^G(\Res_H^G V\otimes W)\cong V\otimes\Ind_H^G W.
\]
\end{proposition}
\begin{proof}
把 $g\otimes(v\otimes w)$ 映到 $gv\otimes(g\otimes w)$. 对 $h\in H$, 源中的 $gh\otimes(v\otimes w)$ 等于 $g\otimes(hv\otimes hw)$; 两者输出都为 $ghv\otimes(g\otimes hw)$. 因而映射良定义. 它与 $G$ 作用兼容. 逆映射把 $v\otimes(g\otimes w)$ 映为 $g\otimes(g^{-1}v\otimes w)$; 替换 $g$ 为 $gh$ 时, 相应 $h^{-1}$ 恰由平衡关系消去. 在生成元上核对两个复合为恒等即可.
\end{proof}

\begin{example}\label{ex:s3induction}
令 $H=C_3\trianglelefteq S_3$. 分解 $\Ind_H^{S_3}L_j$, $j=0,1,2$.
\end{example}
\begin{solution}
$1$ 与 $\sgn$ 限制到 $H$ 均为平凡表示; 标准表示限制为 $L_1\oplus L_2$. 互反给出 $\Ind L_0$ 在三种不可约中的重数为 $1,1,0$, 所以 $\Ind L_0\cong1\oplus\sgn$. $\Ind L_1$ 的重数为 $0,0,1$, 所以等价于标准表示; $L_2$ 同理. 各诱导表示维数为2, 与所得分解一致. 也可由诱导特征标计算: 两个非平凡诱导在三类上取 $2,0,\omega+\omega^2= -1$.
\end{solution}

\begin{exercise}
求 $S_3$ 从 $H=\langle(12)\rangle\cong C_2$ 诱导平凡表示及非平凡一维表示的分解.
\end{exercise}
\begin{solution}
标准表示在换位上的特征值为 $1,-1$, 所以限制为平凡加非平凡. $\sgn$ 限制为非平凡. 互反得 $\Ind_H^{S_3}1\cong1\oplus U$, $\Ind_H^{S_3}\varepsilon\cong\sgn\oplus U$. 两者都为三维; 前者就是三个陪集上的置换表示.
\end{solution}

\originalchapter{双陪集与二面体群表示}
\originalsection{Mackey分解}
诱导后再限制到另一个子群 $K$, 陪集块不一定仍被 $K$ 传递地置换. 应把这些块按 $K$ 的轨道分组. $K$ 在 $G/H$ 上的轨道恰由双陪集 $KtH$ 参数化. 双陪集指所有 $kth$ ($k\in K,h\in H$) 组成的集合, 两个这样的集合或者相同或者不交.

\begin{theorem}[Mackey分解]\label{thm:mackey}
取 $K\backslash G/H$ 的代表 $t$. 令 $L_t=K\cap tHt^{-1}$, 在 $W$ 的原空间上定义 $L_t$ 表示 $W_t$, 使 $\ell$ 的作用为 $\rho_W(t^{-1}\ell t)$. 则
\[
 \Res_K^G\Ind_H^G W\cong
 \bigoplus_{t\in K\backslash G/H}\Ind_{L_t}^K W_t.
\]
\end{theorem}
\begin{proof}
依照 $K$ 在左陪集上的轨道, 把 $\Ind_H^G W$ 的陪集直和分成若干 $K$ 不变子空间 $M_t$, 分别由 $kt\otimes w$ 张成. 定义
\[
 \C[K]\otimes_{\C[L_t]} W_t\longrightarrow M_t,\qquad
 k\otimes w\longmapsto kt\otimes w.
\]
若 $\ell\in L_t$, 则 $k\ell t=kt(t^{-1}\ell t)$, 因而 $(k\ell)t\otimes w=kt\otimes\rho_W(t^{-1}\ell t)w$. 这正是源中的平衡关系, 所以映射良定义且交织. 它显然满射. 左陪集 $kL_t$ 与轨道中的陪集 $ktH$ 一一对应: 相同条件都是 $k_2^{-1}k_1\in K\cap tHt^{-1}$. 所以两边维数均为 $[K:L_t]\dim W$, 满射即同构. 所有轨道的直和给出结论.
\end{proof}
旧讲义§4.9把诱导特征标的陪集求和称为Mackey公式; 这里进一步给出双陪集版本, 是编者为计算诱导不可约性补充的结构论证. 两者均来自同一个陪集分块.

\originalsection{诱导不可约性的判据}
\begin{corollary}\label{cor:mackeycriterion}
若 $W$ 是 $H$ 不可约表示, 则 $\Ind_H^G W$ 不可约, 当且仅当对每个非单位双陪集代表 $t$ 有
\[
 \Hom_{H\cap tHt^{-1}}(\Res_{H\cap tHt^{-1}}^H W,W_t)=0.
\]
\end{corollary}
\begin{proof}
由互反,
$\dim\End_G(\Ind W)=\dim\Hom_H(W,\Res_H^G\Ind W)$. 再用Mackey分解与另一方向的互反, 右侧为所有双陪集对应的上述Hom维数之和. 单位双陪集可取 $t=e$, 对应 $\End_H(W)$, 维数为1. 其余项均为非负整数, 所以总维数为1恰当且仅当它们都为0. 完全可约表示的交织代数维数为重数平方和, 等于1恰对应不可约.
\end{proof}
当 $H\trianglelefteq G$, 每个 $L_t=H$, 所以判据简化为 $W$ 与各非平凡商陪集给出的共轭表示互不等价. 若某个共轭表示与 $W$ 等价, 诱导中就存在额外交织算子, 因而会分裂.

\originalsection{一般二面体群}
令 $D_{2n}=\langle r,s:r^n=s^2=e,\ srs=r^{-1}\rangle$, $n\ge3$, 阶为 $2n$. 旋转子群 $H=\langle r\rangle\cong C_n$ 正规且指数为2. 令 $\zeta=\exp(2\pi i/n)$, $L_j(r)=\zeta^j$, 并记 $E_j=\Ind_H^{D_{2n}}L_j$.

选陪集代表 $e,s$. 在相应两份空间的基中,
\[
 \rho_{E_j}(r)=\begin{pmatrix}\zeta^j&0\\0&\zeta^{-j}\end{pmatrix},\qquad
 \rho_{E_j}(s)=\begin{pmatrix}0&1\\1&0\end{pmatrix}.
\]
限制到 $H$ 是 $L_j\oplus L_{-j}$; $s$ 交换两条特征直线. 当 $j\not\equiv-j\pmod n$ 时, 二面体作用没有不变直线, 因为任何 $r$ 不变直线都只能是两条特征直线之一, 而它们被 $s$ 交换. 因而 $E_j$ 不可约. 这也是上述正规子群判据的具体情形.

若 $j=0$, $r$ 作用恒等, $s$ 的 $\pm1$ 特征线给出两个一维表示. 若 $n$ 为偶数且 $j=n/2$, $r$ 作用为 $-I$, 同样分成 $s$ 取 $\pm1$ 的两个一维表示. 其他 $E_j$ 与 $E_{-j}$ 等价, 用交换两个基向量实现; 若两组 $\{j,-j\}$ 不同, 限制到 $H$ 的特征标不同, 故不等价.

\begin{theorem}
若 $n$ 为奇数, $D_{2n}$ 有两个一维表示, 以及 $E_j$, $1\le j\le(n-1)/2$. 若 $n$ 为偶数, 它有四个一维表示, 以及 $E_j$, $1\le j<n/2$. 这些列表穷尽全部复不可约表示. 二维表示的特征标为
\[
 \chi_{E_j}(r^a)=\zeta^{ja}+\zeta^{-ja}=2\cos(2\pi ja/n),\qquad
 \chi_{E_j}(r^as)=0.
\]
\end{theorem}
\begin{proof}
一维表示满足 $r\mapsto a=a^{-1}$, $s\mapsto b=\pm1$, 且 $a^n=1$. 奇数 $n$ 只能 $a=1$, 偶数 $n$ 可取 $a=\pm1$, 所以分别有二个或四个. 上面已构造并验证二维不可约及互不等价. 奇数时维数平方和为 $2+4(n-1)/2=2n$; 偶数时为 $4+4(n/2-1)=2n$. 所以分类完整. 特征标由所列矩阵直接取迹得到.
\end{proof}
取 $n=4$ 即回到第五章的 $D_8$; 取 $n=3$ 得 $D_6\cong S_3$. 一般分类把这些小群的表统一起来.

\originalsection{限制与分支计算}
\begin{example}
设 $m\mid n$ 且 $m\ge3$, $D_{2m}=\langle r^{n/m},s\rangle\le D_{2n}$. 求 $E_j$ 限制到这个子群的分解.
\end{example}
\begin{solution}
$u=r^{n/m}$ 的矩阵对角元为 $\exp(2\pi i j/m)$ 与其逆, $s$ 仍交换两个基向量. 因此若 $j\not\equiv-j\pmod m$, 限制就是 $D_{2m}$ 的二维表示 $E_{j\bmod m}$. 若 $j\equiv0\pmod m$, 它分成 $u$ 取1、$s$ 分别取 $\pm1$ 的两个一维表示. 若 $m$ 为偶数且 $j\equiv m/2\pmod m$, 则分成 $u$ 取 $-1$、$s$ 分别取 $\pm1$ 的两个一维表示. 这说明大群的不可约可以在子群上出现完整或分裂两种行为.
\end{solution}

\begin{exercise}
设 $H\trianglelefteq G$, $W$ 为 $H$ 不可约表示, 各商陪集给出的共轭表示互不等价. 证明 $\Res_H^G\Ind_H^G W$ 是这些共轭表示各出现一次的直和, 并求诱导表示的维数.
\end{exercise}
\begin{solution}
正规性给 $H\cap tHt^{-1}=H$. Mackey分解每项都为 $W_t$, 代表 $t$ 遍历 $G/H$. 假设保证它们两两不等价, 所以限制中各出现一次, 同时判据说明诱导不可约. 维数为 $[G:H]\dim W$.
\end{solution}

\originalchapter{不变形式与四次对称群}
\originalsection{不变双线性形式}
Hermite内积在一个变量上线性、另一个变量上共轭线性. 本节的双线性形式 $B:V\times V\to\C$ 则在两个变量上都复线性. 二者不能混同. 若 $B(gv,gw)=B(v,w)$ 对所有 $g$ 成立, 称 $G$ 不变. 它对应交织映射 $T_B:V\to V^*$, 定义 $T_B(v)(w)=B(v,w)$: 不变性等价于 $T_B(gv)(w)=T_B(v)(g^{-1}w)$.

\begin{proposition}\label{prop:invariantbilinear}
若 $V$ 复不可约, 则不变双线性形式空间在 $V\not\cong V^*$ 时为0; 在 $V\cong V^*$ 时为一维. 后一情形每个非零不变形式非退化, 且必为对称或交错形式之一.
\end{proposition}
\begin{proof}
形式空间同构于 $\Hom_G(V,V^*)$. $V^*$ 也不可约, 所以Schur引理给出维数0或1. 非零 $T_B$ 必为同构, 因而 $B$ 非退化. 转置形式 $B^{\mathsf T}(v,w)=B(w,v)$ 也不变, 所以 $B^{\mathsf T}=cB$. 再转置得 $B=c^2B$, 所以 $c=\pm1$. $c=1$ 即对称; $c=-1$ 给 $B(v,v)=0$, 即交错.
\end{proof}
交错非退化形式只能在偶数维空间存在: 其矩阵满足 $A^{\mathsf T}=-A$, 若维数 $d$ 为奇数, 则 $\det A=\det A^{\mathsf T}=\det(-A)=(-1)^d\det A=-\det A$, 从而行列式为0.

\originalsection{Frobenius--Schur指标}
\begin{theorem}\label{thm:fs}
对不可约复表示 $V$ 定义 $\nu(V)=|G|^{-1}\sum_g\chi_V(g^2)$. 则
\[
 \nu(V)=\begin{cases}
 0,&V\not\cong V^*,\\
 1,&V\text{有非零不变对称双线性形式},\\
 -1,&V\text{有非零不变交错双线性形式}.
 \end{cases}
\]
\end{theorem}
\begin{proof}
把双线性形式看成 $V^*\otimes V^*$ 中的张量, 其对称和交错部分分别为 $\Sym^2(V^*)$ 与 $\bigwedge^2(V^*)$. 交换两变量的算子与群作用交换, 所以不变形式空间也按这两部分分解. 由命题\ref{prop:symext}与平均投影公式, 两个不变空间维数之差为
\[
 \frac1{|G|}\sum_g\bigl(\chi_{\Sym^2 V^*}(g)-\chi_{\bigwedge^2 V^*}(g)\bigr)
 =\frac1{|G|}\sum_g\chi_{V^*}(g^2)=\overline{\nu(V)}.
\]
命题\ref{prop:invariantbilinear}说明: 非自对偶时两维数都是0; 自对偶时只有一种形式, 对称情况的差为1, 交错情况为 $-1$. 因而差是实数, 等于其共轭, 得到结论.
\end{proof}
本书用指标辨别三种不变形式类型. 与实表示、四元数表示的完整对应在课程习题原题4.1中进一步说明; 这里的证明不借助该对应.

\begin{example}
比较第五章 $D_8$ 的二维表示 $E$ 与 $Q_8$ 的二维表示 $F$ 的指标.
\end{example}
\begin{solution}
$D_8$ 中 $e,r^2$ 及四个反射平方为 $e$, 两个奇数次旋转平方为 $r^2$. 因此 $\nu(E)=(6\cdot2+2\cdot(-2))/8=1$. $E$ 的实正交矩阵也直接保持 $v^{\mathsf T}w$ 这一对称双线性形式.

$Q_8$ 中 $\pm1$ 的平方为1, 其余六个元素平方为 $-1$, 所以 $\nu(F)=(2\cdot2+6\cdot(-2))/8=-1$. 矩阵 $A=\begin{pmatrix}0&1\\-1&0\end{pmatrix}$ 定义交错形式 $B(v,w)=v^{\mathsf T}Aw$. 第五章的两生成元矩阵均行列式为1, 而任意二阶矩阵 $M$ 满足 $M^{\mathsf T}AM=(\det M)A$, 所以此形式不变. 二群的特征标表相同, 平方映射不同, 指标因而不同.
\end{solution}

\originalsection{四次对称群的标准表示}
$S_4$ 的共轭类按循环类型分为 $e,(12),(12)(34),(123),(1234)$, 类大小分别为 $1,6,3,8,6$. 换位为选择两个数字, 共6个; 双换位为把四个数字分成两对, 共3个; 三循环先选三个数字再选两个方向, 共8个; 四循环共 $(4-1)!=6$ 个. 类大小之和为24.

四点置换表示分成平凡直线与标准空间 $W=\{(a_1,\ldots,a_4):\sum a_i=0\}$, 所以 $\chi_W(g)$ 为固定点数减1, 得 $3,1,-1,0,-1$. 平方范数是 $(9+6+3+0+6)/24=1$, 故 $W$ 不可约. 张量符号表示给 $W'=W\otimes\sgn$, 特征标为 $3,-1,-1,0,1$, 也不可约.

\originalsection{商群构造与完整特征标表}
考虑三种把四点分成两对的划分
\[
 P_1=12|34,\qquad P_2=13|24,\qquad P_3=14|23.
\]
$S_4$ 置换这三个对象, 给同态 $S_4\to S_3$. 换位 $(12)$ 交换 $P_2,P_3$, 三循环在三划分上给三循环, 所以像包含 $S_3$ 的生成元, 因而满射. 核包含正规四元群 $N=\{e,(12)(34),(13)(24),(14)(23)\}$, 而核的阶为 $24/6=4$, 所以恰为 $N$. 把 $S_3$ 的二维标准表示沿此商映射拉回, 得不可约 $E$: 不变子空间与商群上的不变子空间相同.

在三划分上的置换特征标减1, 得 $\chi_E=2,0,2,-1,0$. 所以全部候选列表为
\begin{center}
\begin{tabular}{c|rrrrr}
\toprule
代表元 &$e$&$(12)$&$(12)(34)$&$(123)$&$(1234)$\\
类大小 &1&6&3&8&6\\\midrule
$1$&1&1&1&1&1\\
$\sgn$&1&$-1$&1&1&$-1$\\
$E$&2&0&2&$-1$&0\\
$W$&3&1&$-1$&0&$-1$\\
$W'$&3&$-1$&$-1$&0&1\\\bottomrule
\end{tabular}
\end{center}
已构造的五个不可约两两特征标不同, 平方和为 $1+1+4+9+9=24$, 因而穷尽. 全部具有实矩阵, 并保持平均后的实正定内积; 其复线性延拓为对称不变双线性形式, 所以指标均为1. 这不需要一般Young图理论.

\begin{example}
设 $H=S_3\le S_4$ 是第四个点的稳定子, $U$ 为 $H$ 的二维标准表示. 分解 $\Ind_H^{S_4}U$.
\end{example}
\begin{solution}
$W$ 限制到 $H$ 的特征标为 $(3,1,0)$, 等于 $1_H\oplus U$ 的特征标; 因此 $\Res W\cong1_H\oplus U$. 同理 $\Res W'\cong\sgn_H\oplus U$, 因为 $U\otimes\sgn_H\cong U$. 二维 $E$ 限制到 $H$ 的值是 $(2,0,-1)$, 即 $U$. 两个一维 $S_4$ 表示限制后都不含 $U$. 互反得到 $\Ind_H^{S_4}U\cong E\oplus W\oplus W'$. 维数为 $2+3+3=8=[S_4:H]\cdot2$.
\end{solution}

\begin{exercise}
求 $W\otimes W$、$\Sym^2W$ 与 $\bigwedge^2W$ 的分解.
\end{exercise}
\begin{solution}
张量特征标为 $(9,1,1,0,1)$, 与五行取内积给重数 $1,0,1,1,1$, 所以 $W\otimes W\cong1\oplus E\oplus W\oplus W'$. 平方映射的类顺序给 $\chi_W(g^2)=(3,3,3,0,-1)$, 从而对称平方为 $(6,2,2,0,0)$, 外平方为 $(3,-1,-1,0,1)$. 与表比较得 $\Sym^2W\cong1\oplus E\oplus W$, $\bigwedge^2W\cong W'$.
\end{solution}

\originalchapter{课程习题与考试}
\label{ch:course-problems}
以下题目取自MIT 18.712, Fall 2010官方作业清单及期末作业, 与前九章编者练习分开编号. 所有解答均由AI辅助独立编写并作数学复核, 不是MIT公布的标准答案. 旧讲义已有的提示在题解中用于参考并核查; 原题笔误及本稿采用的条件会明确说明. 官方未提供单独测验、期中卷或练习考试.
% MIT 18.712 Fall 2010, Homework 1--4. 中文题面据2011-02-01旧讲义逐页核对.
% 解答由编者独立推导, 并非MIT官方解答. 不新增排版宏.
\originalsection{Homework 1: 子表示与循环表示}

以下题目中的代数均有单位, 表示保持单位. 依旧讲义第7页的总约定, 未另作说明时底域 $k$ 代数闭. 一个非零表示若不能写成两个非零子表示的直和, 称为不可分解表示. 这个条件弱于不可约性. 本节至 Homework 4 的题解均为编者独立编写, 不标作官方答案.

\begin{exercise}\label{hw:1.20}
MIT 18.712 (2010), Homework 1, 原题1.20. 设 $V$ 是代数 $A$ 的非零有限维表示. 证明 $V$ 含有不可约子表示, 并举例说明无限维表示未必有此性质.
\end{exercise}
\begin{solution}
在 $V$ 的所有非零子表示中, 取维数最小的一个 $W$. 这样的选择存在, 因为这些维数是 $1,\ldots,\dim V$ 中的整数. 若 $W$ 有非零真子表示 $U$, 则 $\dim U<\dim W$, 与选择矛盾. 所以 $W$ 不可约.

无限维反例取 $A=k[x]$, $V=A$ 为左正则表示. 它的子表示正是 $k[x]$ 的理想. 若 $J$ 是非零理想, 则 $xJ$ 是非零子理想; 而 $xJ\ne J$, 因为在 $J$ 中选取次数最小的非零多项式 $f$, $xJ$ 中所有非零多项式的次数都严格大于 $\deg f$. 因此任一非零子表示仍含非零真子表示, $V$ 没有不可约子表示.
\end{solution}

\begin{exercise}\label{hw:1.21}
MIT 18.712 (2010), Homework 1, 原题1.21. 代数 $A$ 的中心定义为 $Z(A)=\{z\in A:za=az\text{ 对所有 }a\in A\}$. 若 $A$ 交换, 则 $Z(A)=A$.
\begin{enumerate}[label=(\alph*)]
\item 证明: 若 $V$ 是有限维不可约表示, 则每个 $z\in Z(A)$ 在 $V$ 上以某标量 $\chi_V(z)$ 作用, 而 $\chi_V:Z(A)\to k$ 是代数同态, 称为中心特征.
\item 证明: 若 $V$ 是有限维不可分解表示, 则 $\rho(z)$ 只有一个特征值 $\chi_V(z)$, 它等于 $z$ 在任一不可约子表示上作用的标量; 这些标量仍给出同态 $\chi_V:Z(A)\to k$.
\item (b)中的 $\rho(z)$ 是否一定是标量算子?
\end{enumerate}
\end{exercise}
\begin{solution}
(a) 因为 $z$ 在中心, $\rho(z)$ 与所有 $\rho(a)$ 交换. 取它的一个特征值 $c\in k$. 非零特征空间 $\ker(\rho(z)-c\id)$ 是子表示, 故不可约性迫使它等于 $V$. 于是 $\rho(z)=c\id$. 由表示的线性、乘法及单位性质, 有
\[
\chi_V(z+z')=\chi_V(z)+\chi_V(z'),\quad
\chi_V(zz')=\chi_V(z)\chi_V(z'),\quad
\chi_V(cz)=c\chi_V(z),\quad \chi_V(1)=1.
\]

(b) $\rho(z)$ 的各广义特征空间都被 $A$ 保持, 因为 $\rho(a)$ 与 $(\rho(z)-c\id)^r$ 交换. 代数闭域上的广义特征空间分解是直和分解. 不可分解性因此保证 $\rho(z)$ 恰有一个特征值. 由上一题取不可约子表示 $W\subset V$. (a)说明 $z$ 在 $W$ 上以标量 $\chi_W(z)$ 作用, 而该标量必须是 $\rho(z)$ 的唯一特征值. 因此对所有 $z$ 有 $\chi_V(z)=\chi_W(z)$, (a)的同态性质立即传给 $\chi_V$. 这也说明选取不同的 $W$ 不会改变所得同态.

(c) 不一定. 取 $A=k[t]$, 在 $V=k^2$ 上令 $t$ 以非零二阶幂零Jordan块 $N$ 作用. 若 $V$ 分解为两个一维子表示, 则 $N$ 在两条直线上都只能以特征值0作用, 从而 $N=0$, 矛盾. 故 $V$ 不可分解, 但中心元 $t$ 的作用并非标量. 此题依赖代数闭性: 例如把 $\C$ 视作实代数, 其一维复正则模是不可约实表示, $\mathrm i$ 的作用却不是实标量算子.
\end{solution}

\begin{exercise}\label{hw:1.22}
MIT 18.712 (2010), Homework 1, 原题1.22. 对表示 $V$, 记 $\End_A(V)$ 为所有 $A$-线性自同态组成的代数. 证明左正则表示满足 $\End_A(A)\cong A^{\mathrm{op}}$, 其中 $A^{\mathrm{op}}$ 的乘法与 $A$ 相反.
\end{exercise}
\begin{solution}
若 $T:A\to A$ 是 $A$-线性的, 则 $T(a)=T(a\cdot1)=aT(1)$. 因而 $T$ 正是右乘 $b=T(1)$ 的算子 $R_b$. 反过来, $R_b(ac)=a(cb)$, 所以右乘算子总是 $A$-线性的. 映射 $b\mapsto R_b$ 是向量空间同构, 且 $R_bR_c(a)=(ac)b=a(cb)$, 即 $R_bR_c=R_{cb}$. 这正是 $A^{\mathrm{op}}$ 的乘法.
\end{solution}

\begin{exercise}\label{hw:1.23}
MIT 18.712 (2010), Homework 1, 原题1.23. 证明Dixmier的无限维Schur引理: 若 $A$ 是复代数, $V$ 是具有至多可数基的不可约表示, 则每个表示同态 $\phi:V\to V$ 都是标量算子. 原题提示先考察除代数 $D=\End_A(V)$ 及其子域 $\C(\phi)$.
\end{exercise}
\begin{solution}
非零 $T\in D$ 的核与像都是子表示, 因而核为0、像为 $V$; 逆映射也是 $A$-线性的. 故 $D$ 是除代数. 固定 $0\ne v\in V$. 评价映射 $D\to V$, $T\mapsto Tv$ 是单射: 如果 $Tv=0$, 则 $T(av)=aTv=0$, 而不可约性给出 $Av=V$. 因此 $D$ 的复维数至多可数.

假设 $\phi$ 不是标量. 若有非零多项式 $p$ 满足 $p(\phi)=0$, 在 $\C$ 上分解 $p(t)=c\prod_i(t-a_i)$. 除代数没有零因子, 于是某个 $\phi-a_i\id=0$, 矛盾. 所以 $\phi$ 对 $\C$ 超越, 每个非零多项式 $p(\phi)$ 均可逆. 它们和 $\phi$ 彼此交换, 给出嵌入 $\C(t)\hookrightarrow D$.

但 $\{(t-a)^{-1}:a\in\C\}$ 是不可数线性无关族. 确实, 若有限和 $\sum_{i=1}^r c_i/(t-a_i)=0$ 且 $a_i$ 两两不同, 乘以 $\prod_i(t-a_i)$ 再令 $t=a_j$, 得到 $c_j\prod_{i\ne j}(a_j-a_i)=0$, 即所有 $c_j=0$. 这与 $D$ 至多可数维矛盾. 故 $\phi$ 是标量.
\end{solution}

\begin{exercise}\label{hw:1.24}
MIT 18.712 (2010), Homework 1, 原题1.24. 设 $A=k[x_1,\ldots,x_n]$, 真理想 $I\ne A$ 包含所有次数 $\ge N$ 的齐次多项式. 证明 $A/I$ 是不可分解表示.
\end{exercise}
\begin{solution}
记 $B=A/I$, $\mathfrak m=(\bar x_1,\ldots,\bar x_n)\subset B$. 因为 $\mathfrak m^N=0$, 每个 $b\in B$ 可写成 $c+u$, 其中 $c\in k$, $u\in\mathfrak m$. 若 $c\ne0$, 则
\[
(c+u)^{-1}=c^{-1}\sum_{j=0}^{N-1}(-c^{-1}u)^j.
\]
若 $I$ 包含常数项非零的多项式, 其在 $B$ 中的像为0又可逆, 就会得到 $B=0$, 与 $I\ne A$ 矛盾. 所以 $B/\mathfrak m\cong k$, $\mathfrak m$ 是唯一极大理想.

若 $B=M\oplus L$ 是两个子表示的直和, 到 $M$ 的投影 $P$ 属于 $\End_A(B)$. 它由 $e=P(1)$ 唯一决定, 因为 $P(\bar a)=\bar a e$, 且 $P^2=P$ 给出 $e^2=e$. 在 $B/\mathfrak m=k$ 中, $e$ 的像是0或1. 前者使 $e$ 幂零, 结合 $e^r=e$ 得 $e=0$; 后者同样使幂等元 $1-e$ 为0, 即 $e=1$. 投影只能是零或恒等, 因此不存在非平凡直和分解. 这里未假设 $I$ 自身齐次; 所需条件是它包含全部高次齐次多项式.
\end{solution}

\begin{exercise}\label{hw:1.25}
MIT 18.712 (2010), Homework 1, 原题1.25. 对非零表示 $V$, 若 $Av=V$, 则称 $v$ 为循环向量; 有循环向量的表示称为循环表示.
\begin{enumerate}[label=(\alph*)]
\item 证明 $V$ 不可约当且仅当每个非零向量都是循环向量.
\item 证明 $V$ 循环当且仅当 $V\cong A/I$, 其中 $I$ 是某左理想.
\item 举出不可分解但不循环的表示. 原题提示取 $A=\C[x,y]/(x,y)^2$, $V=A^*$, 令 $(af)(b)=f(ba)$.
\end{enumerate}
\end{exercise}
\begin{solution}
(a) 对非零 $v$, $Av$ 是包含 $v=1v$ 的非零子表示. 不可约性保证它等于 $V$. 反之, 若 $U$ 是非零子表示, 取 $0\ne v\in U$, 则 $V=Av\subset U$, 所以 $U=V$.

(b) 若 $v$ 循环, 则 $A\to V$, $a\mapsto av$ 是满射表示同态; 其核 $I$ 是左理想, 因而诱导 $A/I\cong V$. 反过来, $1+I$ 在 $A/I$ 中是循环向量.

(c) $A$ 的基为 $1,x,y$. 取对偶基 $\varepsilon_0,\varepsilon_x,\varepsilon_y$, 则
\[
x\varepsilon_x=\varepsilon_0,\quad y\varepsilon_y=\varepsilon_0,
\qquad x\varepsilon_0=x\varepsilon_y=y\varepsilon_0=y\varepsilon_x=0.
\]
对 $f=a\varepsilon_0+b\varepsilon_x+c\varepsilon_y$, 子空间 $Af$ 包含在 $\operatorname{span}(f,\varepsilon_0)$ 中, 维数至多2, 所以三维的 $V$ 不循环. 任一非零子表示 $U$ 都包含 $\varepsilon_0$: 如果它有 $b$ 或 $c$ 非零的向量, 作用 $x$ 或 $y$ 即可得到非零倍的 $\varepsilon_0$; 否则它已经包含在 $\C\varepsilon_0$ 中. 两个非零子表示因此交非零, 不能构成直和, $V$ 不可分解.
\end{solution}

\originalsection{Homework 2: 生成关系与张量运算}

\begin{exercise}\label{hw:1.26}
MIT 18.712 (2010), Homework 2, 原题1.26. Weyl代数 $A$ 由 $x,y$ 生成, 满足 $yx-xy=1$.
\begin{enumerate}[label=(\alph*)]
\item 若 $\operatorname{char}k=0$, 求所有有限维表示及所有双边理想.
\item 以下设 $\operatorname{char}k=p>0$. 求 $A$ 的中心.
\item 求所有有限维不可约表示. 原题提示从 $y$ 的特征向量 $v$ 出发, 证明 $v,xv,\ldots,x^{p-1}v$ 构成基.
\end{enumerate}
\end{exercise}
\begin{solution}
先说明计算所用的基. 关系 $yx=xy+1$ 把每个单词化为 $x^iy^j$ 的线性组合, $i,j\ge0$. 这些正规单词也线性无关. 可在自由代数中按此规则约化: 每次交换降低逆序数, 出现的常数项则降低单词长度, 所以约化终止. 两个可约的 $yx$ 不会有共同字母; 对不相交的两个位置, 先约化哪一个都给出相同结果. 对单词长度和逆序数归纳, 最终正规式唯一. 因而 $\{x^iy^j\}$ 确是基, 这在任意特征都成立.

(a) 若 $X,Y$ 是 $d$ 维表示中的矩阵, 则 $YX-XY=\id$ 的迹给出 $0=d$. 特征0时这迫使 $d=0$, 所以只有零表示. 对非零双边理想 $J$, 取非零 $u=\sum_{j=0}^r f_j(x)y^j\in J$, $f_r\ne0$. 由于
\[
[x,u]=-\sum_{j=1}^r j f_j(x)y^{j-1},
\quad (\operatorname{ad}x)^r u=(-1)^rr!f_r(x),
\]
理想 $J$ 含非零多项式 $f_r(x)$. 再用 $[y,f(x)]=f'(x)$, 连续取次数多次交换子, 得到非零常数, 即 $1\in J$. 所以双边理想只有0和 $A$.

(b) 在特征 $p$ 下, $[y,x^p]=p x^{p-1}=0$, $[x,y^p]=-p y^{p-1}=0$, 故 $k[x^p,y^p]\subset Z(A)$. 对一般元素 $u=\sum c_{ij}x^iy^j$, 正规基的独立性表明 $[y,u]=0$ 当且仅当所有出现的 $i$ 均被 $p$ 整除; 同理 $[x,u]=0$ 当且仅当所有 $j$ 均被 $p$ 整除. 因此 $Z(A)=k[x^p,y^p]$.

(c) 仍依旧讲义约定令 $k$ 代数闭. 中心元 $x^p,y^p$ 在不可约有限维 $V$ 中分别作用为标量 $a,b$. 取 $yv=\lambda v$, $v\ne0$, 则 $\lambda^p=b$. 由 $yx^i=x^iy+i x^{i-1}$, 空间
\[
U=\operatorname{span}(v,xv,\ldots,x^{p-1}v)
\]
被 $x,y$ 保持: 最后一个向量再乘 $x$ 得到 $av$. 所以 $U=V$. 若有关系 $\sum_{i=0}^r c_i x^iv=0$ 且 $r<p$, $c_r\ne0$, 作用 $(y-\lambda)^r$ 得 $r!c_rv=0$, 矛盾. 因而上述 $p$ 个向量是基.

每个 $(a,b)\in k^2$ 都给出一个表示:
\[
V_{a,b}=k[t]/(t^p-a),\qquad
x f=t f,\quad y f=\frac{df}{dt}+\lambda f,
\qquad \lambda^p=b.
\]
导数保持理想 $(t^p-a)$, 所以作用定义良好, 且 $yx-xy=1$, $x^p=a\id$, $y^p=b\id$. 非零不变子空间中取次数为 $r<p$ 的非零多项式, 作用 $(y-\lambda)^r$ 得到非零常数, 再乘 $x$ 得到全空间. 故它不可约. 上述基计算说明任何不可约表示都同构于其中一个, 而不同 $(a,b)$ 由中心作用区分. 代数闭特征 $p$ 的域中 $p$ 次根唯一, 无额外参数.
\end{solution}

\begin{exercise}\label{hw:1.27}
MIT 18.712 (2010), Homework 2, 原题1.27. 令 $q\in\C^\times$, 代数 $A$ 由可逆元 $x,y$ 生成, 满足 $xy=qyx$.
\begin{enumerate}[label=(\alph*)]
\item 对不同的 $q$ 求中心; 若 $q$ 不是单位根, 求所有双边理想.
\item 哪些 $q$ 使 $A$ 有非零有限维表示?
\item 对这些 $q$, 求所有有限维不可约表示.
\end{enumerate}
\end{exercise}
\begin{solution}
本题严格使用 $xy=qyx$ 的方向. 基为 $x^iy^j$, $i,j\in\Z$, 其乘法为
\[
(x^iy^j)(x^ry^s)=q^{-jr}x^{i+r}y^{j+s}.
\]
右侧本身定义一个结合代数并满足给定关系, 又每个单词可整理为这些单项式, 故它也证明基的独立性.

(a) 共轭给出 $x(x^iy^j)x^{-1}=q^j x^iy^j$, $y(x^iy^j)y^{-1}=q^{-i}x^iy^j$. 若 $q$ 不是单位根, 中心只有 $\C$. 若 $q$ 的精确阶为 $\ell\ge1$, 则中心为 $\C[x^{\pm\ell},y^{\pm\ell}]$.

当 $q$ 不是单位根时, 对非零双边理想 $J$, 取支持项数最少的非零元素 $u$. 左乘可逆单项式并乘标量后, 可设其常数项为1. 若 $u$ 有某项的 $y$ 指数 $j\ne0$, 则 $xux^{-1}-u$ 是非零的 $J$ 中元素, 其常数项消失, 支持更小, 矛盾. 所以所有项的 $j=0$. 再用 $yuy^{-1}-u$ 同理消去 $i\ne0$ 的可能性. 于是 $u=1$, $J=A$. 双边理想只有0和 $A$.

(b) 对 $d>0$ 维表示的可逆矩阵 $X,Y$, 行列式关系 $\det(XY)=q^d\det(YX)$ 给出 $q^d=1$. 因而 $q$ 必须是单位根. 下面的构造证明此条件充分.

(c) 设 $q$ 精确阶为 $\ell$. 中心作用使 $X^\ell=a\id$, $Y^\ell=b\id$, 其中 $a,b\ne0$. 取 $Yv=\beta v$, $\beta^\ell=b$. 因为 $YX=q^{-1}XY$, 向量 $X^iv$ 的 $Y$ 特征值是 $\beta q^{-i}$, $0\le i<\ell$, 两两不同. 它们线性无关, 且张成空间被 $X,Y$ 保持, 故不可约性使它们构成全空间的基. 于是表示维数为 $\ell$.

对任意 $a,b\in\C^\times$, 在基 $e_0,\ldots,e_{\ell-1}$ 上定义
\[
Xe_i=e_{i+1}\ (i<\ell-1),\quad Xe_{\ell-1}=ae_0,
\qquad Ye_i=\beta q^{-i}e_i,\quad \beta^\ell=b.
\]
直接验证 $XY=qYX$. 不变子空间被 $Y$ 的特征投影保持, 故包含某个 $e_i$, 再由可逆 $X$ 生成全部基, 所以表示不可约. 更换 $\beta$ 为另一 $\ell$ 次根相当于沿 $X$ 的循环重新选取起点, 不改变同构类. 因而同构类恰由中心参数 $(a,b)$ 决定. $\ell=1$ 即 $q=1$ 时, 这给出所有一维表示.
\end{solution}

\begin{exercise}\label{hw:1.33}
MIT 18.712 (2010), Homework 2, 原题1.33. 箭图 $Q$ 的顶点集为 $I$, 边集为 $E$; 边 $h$ 从 $h'$ 指向 $h''$. 路径代数 $P_Q$ 的基为所有有向路径, 包括各顶点的长度0路径 $p_i$, 乘法 $ab$ 表示先走 $b$ 再走 $a$, 无法连接时为0. 证明它由 $p_i$ 与单边路径 $a_h$ 生成, 定义关系为
\begin{enumerate}[label=\arabic*.]
\item $p_i^2=p_i$, $p_ip_j=0$ 当 $i\ne j$;
\item $a_hp_{h'}=a_h$, $a_hp_j=0$ 当 $j\ne h'$;
\item $p_{h''}a_h=a_h$, $p_i a_h=0$ 当 $i\ne h''$.
\end{enumerate}
\end{exercise}
\begin{solution}
每条正长度路径是对应单边路径按逆行进顺序的乘积, 所以这些元素生成 $P_Q$, 且三组关系直接由连接规则得到. 还需证明不存在遗漏的关系.

在仅由这些关系定义的代数中, 任意含 $p_i$ 的单词可把 $p_i$ 吸收或变为0. 两条不匹配的边满足
\[
a_h a_g=a_h p_{h'}p_{g''}a_g=0\qquad(h'\ne g'').
\]
所以剩下的非零单词恰是有向路径, 以及长度0路径. 它们的像在 $P_Q$ 中是不同的基向量, 不可能有额外线性关系. 因而由关系定义的代数与 $P_Q$ 同构.

若在有单位代数的范畴表述有限箭图, 还须写明 $1=\sum_i p_i$, 或把右侧直接定义为单位. 否则在自由的有单位代数中仅写原题三组关系, 会保留一个多余的单位分量. 旧讲义第13页已先说明 $\sum_i p_i$ 是路径代数的单位; 此处沿用该约定.
\end{solution}

\begin{exercise}\label{hw:1.38}
MIT 18.712 (2010), Homework 2, 原题1.38. 若 $A=\bigoplus_{n\ge0}A[n]$ 且 $A[n]A[m]\subset A[n+m]$, 称 $A$ 为非负整数分次代数. 各分次有限维时, Hilbert级数为 $h_A(t)=\sum_{n\ge0}\dim A[n]t^n$. 求以下代数的Hilbert级数:
\begin{enumerate}[label=(\alph*)]
\item 多项式代数 $k[x_1,\ldots,x_m]$, 按总次数分次;
\item 自由代数 $k\langle x_1,\ldots,x_m\rangle$, 按单词长度分次;
\item 外代数, 关系为 $x_ix_j+x_jx_i=0$ 且 $x_i^2=0$, 按次数分次;
\item 有限箭图的路径代数 $P_Q$, 令 $\deg p_i=0$, $\deg a_h=1$. 原题提示用邻接矩阵 $M_Q$ 表达.
\end{enumerate}
\end{exercise}
\begin{solution}
(a) 单项式对应指数 $r_1,\ldots,r_m\ge0$, 因而
\[
h_A(t)=\prod_{i=1}^m\left(\sum_{r_i\ge0}t^{r_i}\right)=\frac1{(1-t)^m}.
\]
这也给出 $\dim A[n]=\binom{n+m-1}{m-1}$, 当 $m>0$.

(b) 长度 $n$ 的单词有 $m^n$ 个, 故 $h_A(t)=\sum_{n\ge0}m^nt^n=(1-mt)^{-1}$.

(c) 换序和平方为0的关系使递增指标的单项式 $x_{i_1}\cdots x_{i_n}$, $i_1<\cdots<i_n$, 张成 $A[n]$. 它们独立, 例如在基为所有子集的向量空间上, 让 $x_i$ 以插入 $i$ 并按换序次数赋符号的算子作用: 若已有 $i$, 作用为0; 否则得到相应并集. 这些算子满足全部关系, 不同递增单项式作用于空集得到不同基向量. 这一证明在特征2也成立, 因为平方为0是单独给出的关系. 所以 $\dim A[n]=\binom mn$, $h_A(t)=(1+t)^m$.

(d) 规定 $(M_Q)_{ij}$ 为从顶点 $i$ 到 $j$ 的边数. 矩阵乘法归纳表明 $(M_Q^n)_{ij}$ 是从 $i$ 到 $j$ 的长度 $n$ 路径数. 令 $\mathbf1$ 为全1列向量, 则
\[
h_{P_Q}(t)=\sum_{n\ge0}\mathbf1^{\mathsf T}M_Q^n\mathbf1\,t^n
=\mathbf1^{\mathsf T}(I-tM_Q)^{-1}\mathbf1.
\]
这里是形式幂级数恒等式; 不需要先证明分析意义的收敛. 因为 $I-tM_Q$ 的行列式常数项为1, 逆矩阵给出的表达式也是有理函数.
\end{solution}

\begin{exercise}\label{hw:1.49}
MIT 18.712 (2010), Homework 2, 原题1.49. 张量积 $V\otimes W$ 是把形式符号 $v\otimes w$ 按双线性关系取商得到的向量空间.
\begin{enumerate}[label=(\alph*)]
\item 构造双线性映射 $V\times W\to U$ 与线性映射 $V\otimes W\to U$ 之间的自然双射.
\item 若 $(v_i)$ 与 $(w_j)$ 分别是 $V,W$ 的基, 证明 $(v_i\otimes w_j)$ 是张量积的基.
\item 当 $V$ 有限维时, 构造无须选基的自然同构 $V^*\otimes W\cong\Hom(V,W)$.
\item 定义 $S^nV$ 为 $V^{\otimes n}$ 对所有 $T-sT$ 张成的子空间的商, 其中 $s$ 为换位. 定义 $\bigwedge^nV$ 为对所有满足某个换位 $sT=T$ 的张量 $T$ 张成的子空间的商. 用 $V$ 的基给出这两个空间的基; 当 $\dim V=m$ 时求维数.
\item 特征0时, 将上述两个商自然地分别等同于所有换位下不变的张量子空间、所有换位下变号的张量子空间.
\item 线性算子 $A:V\to W$ 诱导 $S^nA$ 与 $\bigwedge^nA$. 当 $V=W$ 为 $N$ 维, $A$ 的特征值为 $\lambda_1,\ldots,\lambda_N$ 时, 求这两个诱导算子的迹.
\item 证明 $\bigwedge^N A=\det(A)\id$, 并据此证明 $\det(AB)=\det(A)\det(B)$.
\end{enumerate}
\end{exercise}
\begin{solution}
(a) 双线性 $b$ 对形式符号给出 $v\otimes w\mapsto b(v,w)$. 双线性恰保证它消去定义张量积的关系, 因而唯一降为线性映射 $\widetilde b$. 反过来, 对线性 $L$ 令 $b(v,w)=L(v\otimes w)$, 就得到双线性映射. 两种操作互逆, 且都只使用映射与张量积本身, 没有基的选择.

(b) 展开 $v,w$ 的有限基表示可知 $v_i\otimes w_j$ 张成张量积. 令 $v_i^*,w_j^*$ 为取相应坐标的线性泛函. 双线性映射 $(v,w)\mapsto v_i^*(v)w_j^*(w)$ 依(a)降为张量积上的泛函, 在 $v_r\otimes w_s$ 上的值为 $\delta_{ir}\delta_{js}$. 对任一有限线性关系应用这些泛函, 各系数都为0, 所以它们独立.

(c) 映射
\[
\Theta:V^*\otimes W\longrightarrow\Hom(V,W),\qquad
\Theta(f\otimes w)(v)=f(v)w
\]
由(a)定义, 无须选基. 为证明它可逆, 临时取有限基 $v_1,\ldots,v_m$ 及对偶基, 给 $L$ 配上 $\sum_i v_i^*\otimes L(v_i)$. 代入可见 $\Theta$ 将它送回 $L$; 对纯张量的同样计算也给出左逆. 因而 $\Theta$ 是同构, 其逆既唯一, 也与临时选基无关. 有限维条件用于保证上述和是有限和.

(d) 对 $S^nV$, 每个纯基张量可重新排序, 所以基候选为
\[
v_{i_1}\cdots v_{i_n},\qquad i_1\le\cdots\le i_n.
\]
把张量送到具有同名变量的交换单项式, 可证明这些候选独立. 对外幂, 重复指标的纯张量被交换这两个位置的换位固定, 故在商中为0. 对任一张量 $T$, $T+sT$ 被 $s$ 固定, 因而换位在商中变号. 所以外幂由 $v_{i_1}\wedge\cdots\wedge v_{i_n}$, $i_1<\cdots<i_n$, 张成.

为了在任意特征中核对独立性, 把纯基张量的不同指标按排序符号送到相应递增指标的形式基向量, 有重复指标则送到0. 若 $sT=T$, 在纯基张量基中, $s$ 的二元轨道上的系数相等, 两项经过此映射相消; 一元轨道含重复指标, 像为0. 因此映射确实消去原题给出的全部关系, 并把递增候选送到不同基向量. 于是
\[
\dim S^nV=\binom{m+n-1}{n},\qquad
\dim\bigwedge^nV=\binom mn.
\]
约定 $n=0$ 时两个空间均为 $k$, $n>m$ 时外幂为0; $m=0,n>0$ 时两个空间均为0.

(e) 特征0时定义平均算子
\[
P_+=\frac1{n!}\sum_{\sigma\in S_n}\sigma,
\qquad P_-=\frac1{n!}\sum_{\sigma\in S_n}\sgn(\sigma)\sigma.
\]
它们的像分别是不变张量与变号张量, 且在各自的像上为恒等. $P_+$ 消去 $T-sT$; 商映射 $q_+$ 满足 $q_+(P_+T)=q_+(T)$, 故它从商到不变子空间的映射是同构. 对外幂, 若 $sT=T$, 则 $P_-T=P_-sT=-P_-T$, 故 $P_-T=0$. 又商中 $\sigma T=\sgn(\sigma)T$, 所以 $q_-(P_-T)=q_-(T)$, 同理给出同构. 除以 $n!$ 和2在这里都合法.

(f) 在特征值所在的分裂域中将 $A$ 上三角化; 扩域不改变迹. 无须假设 $A$ 可对角化. 在(d)的有序单项式基上, 诱导算子仍为三角矩阵, 对角元是相应特征值的乘积, 故
\begin{align*}
\tr(S^nA)&=\sum_{\substack{r_1+\cdots+r_N=n\\ r_i\ge0}}
\lambda_1^{r_1}\cdots\lambda_N^{r_N},\\
\tr(\bigwedge^nA)&=\sum_{1\le i_1<\cdots<i_n\le N}
\lambda_{i_1}\cdots\lambda_{i_n}.
\end{align*}
第一式是完全齐次对称多项式, 第二式是初等对称多项式. 对 $n>N$, 第二式为空和, 值为0.

(g) 在 $V$ 的基 $v_1,\ldots,v_N$ 下展开 $Av_1\wedge\cdots\wedge Av_N$, 重复指标项为0, 剩余各排列项带符号, 系数恰为矩阵行列式. 由于 $\bigwedge^NV$ 是以 $v_1\wedge\cdots\wedge v_N$ 为基的一维空间, 得到所述等式. 再由作用于纯张量的定义,
\[
\det(AB)\id=\bigwedge^N(AB)
=(\bigwedge^NA)(\bigwedge^NB)=\det(A)\det(B)\id.
\]
这是从顶次外幂的复合性质得到的乘法公式.
\end{solution}

\originalsection{Homework 3: Dehn不变量与李代数表示}

本节为官方作业涉及的扩展内容, 不把有限群主线扩称为一般李代数教材. 李代数是带双线性交替括号的向量空间, 满足Jacobi恒等式 $[x,[y,z]]+[y,[z,x]]+[z,[x,y]]=0$. 其表示 $\rho$ 满足 $\rho([x,y])=\rho(x)\rho(y)-\rho(y)\rho(x)$. 下列题解补齐本题所需的张量作用、最高权和可解性论证.

\begin{exercise}\label{hw:1.51}
MIT 18.712 (2010), Homework 3, 原题1.51. 平面上同面积多边形可以有限次直线切割后重拼. Hilbert第三问题询问同体积多面体是否也有此性质, 特别是立方体和正四面体. 对多面体 $A$, 定义
\[
D(A)=\sum_a l(a)\otimes\overline{\frac{\beta(a)}\pi}
\quad\in\R\otimes_{\mathbb Q}(\R/\mathbb Q),
\]
其中 $a$ 遍历边, $l(a)$ 为边长, $\beta(a)$ 为内部二面角, 横线表示模 $\mathbb Q$ 的类.
\begin{enumerate}[label=(\alph*)]
\item 证明平面直切 $A$ 为 $B,C$ 时, $D(A)=D(B)+D(C)$.
\item 证明 $\alpha=\arccos(1/3)/\pi$ 非有理. 原题提示考察 $z+z^{-1}=2/3$ 及 $z^k+z^{-k}$ 的分母.
\item 计算正四面体和立方体的Dehn不变量, 据此否定上述问题.
\end{enumerate}
\end{exercise}
\begin{solution}
(a) 先把切割前后的全部边按交点细分为公共的小线段. 沿同一条边分成长度 $l_1,l_2$ 时, 张量积的双线性给出 $(l_1+l_2)\otimes\bar\theta=l_1\otimes\bar\theta+l_2\otimes\bar\theta$, 因而这种细分不改变不变量.

对每条小线段, 取垂直于它的局部截面, 观察各块占据的角扇形. 若线段来自原多面体的边, 各块的内部角之和是原二面角; 若它落在原面内部, 各块角之和为 $\pi$; 若落在原多面体内部, 总角为 $2\pi$. 后两类对角除以 $\pi$ 后是整数, 在 $\R/\mathbb Q$ 中为0. 因而各块新增边的贡献总和为0, 原边的贡献则按长度和角度可加. 逐线段求和即得结论. 这种局部扇形论证也允许切割平面经过顶点或原边, 无须用不变量并不保证连续的极限论证. 刚体运动保持长度和内部角, 所以 $D$ 也是重拼不变量.

(b) 假设 $\alpha=2m/n$, 取 $z=e^{\mathrm i\pi\alpha}$, 则 $z^n=1$, $z+z^{-1}=2\cos(\pi\alpha)=2/3$. 记 $u_k=z^k+z^{-k}$. 递推式为
\[
u_0=2,\quad u_1=2/3,\qquad u_{k+1}=(2/3)u_k-u_{k-1}.
\]
写 $u_k=a_k/3^k$, 得 $a_0=2$, $a_1=2$, $a_{k+1}=2a_k-9a_{k-1}$. 对 $k\ge1$, 归纳可知 $a_k\equiv2^k\pmod3$, 所以 $a_k$ 不被3整除. $u_k$ 的最简分母恰为 $3^k$. 但 $z^n=1$ 给出 $u_n=2$, 分母为1, 矛盾.

(c) 立方体的每个二面角为 $\pi/2$, 因而 $D(\text{立方体})=0$. 边长 $l>0$ 的正四面体每条边的内部二面角为 $\arccos(1/3)$, 所以 $D(\text{正四面体})=6l\otimes\bar\alpha$. 这里的角度可由正四面体两相邻面的向外单位法向量内积为 $-1/3$ 得到: 内二面角是法向量夹角的补角.

这个纯张量非零. 在 $\R$ 作为 $\mathbb Q$-向量空间时, 将非零向量 $6l$ 扩充为基, 取线性泛函 $f:\R\to\mathbb Q$ 满足 $f(6l)=1$. 则 $(f\otimes\id)(6l\otimes\bar\alpha)=\bar\alpha\ne0$, 非零性由(b)保证. 因而即使两者体积相同, 它们也不能有限平面切割后重拼.
\end{solution}

\begin{exercise}\label{hw:1.54}
MIT 18.712 (2010), Homework 3, 原题1.54. 设 $V,W,U$ 是李代数 $\mathfrak g$ 的有限维表示. 张量积上的作用为 $x(v\otimes w)=xv\otimes w+v\otimes xw$, 对偶作用为 $(xf)(w)=-f(xw)$. 证明
\[
\Hom_{\mathfrak g}(V\otimes W,U)\cong
\Hom_{\mathfrak g}(V,U\otimes W^*).
\]
\end{exercise}
\begin{solution}
先以题\ref{hw:1.49}(c)的自然同构把 $U\otimes W^*$ 识别为 $\Hom(W,U)$, 其作用为 $(xT)(w)=xT(w)-T(xw)$. 对线性 $F:V\otimes W\to U$, 定义 $\Phi(F)(v)(w)=F(v\otimes w)$. 反向映射把 $T:V\to\Hom(W,U)$ 送到 $v\otimes w\mapsto T(v)(w)$. 双线性保证反向映射定义良好, 两者互逆.

$F$ 是 $\mathfrak g$-线性的条件是
\[
F(xv\otimes w)+F(v\otimes xw)=xF(v\otimes w).
\]
移项后即 $\Phi(F)(xv)(w)=(x\Phi(F)(v))(w)$, 正好是 $\Phi(F)$ 的 $\mathfrak g$-线性条件. 若希望把同构写入 $U\otimes W^*$, 任取 $W$ 的基 $w_i$ 与对偶基 $w_i^*$, 则
\[
\Phi(F)(v)=\sum_i F(v\otimes w_i)\otimes w_i^*.
\]
这个式子的基无关性来自上面的自然同构, 而非额外的选择.
\end{solution}

\begin{exercise}\label{hw:1.55}
MIT 18.712 (2010), Homework 3, 原题1.55. 复李代数 $\mathfrak{sl}_2$ 的表示由算子 $E,F,H$ 描述, 满足
\[
[H,E]=2E,\qquad [H,F]=-2F,\qquad [E,F]=H.
\]
设 $V$ 有限维; 记 $\overline V(\lambda)$ 为 $H$ 对特征值 $\lambda$ 的广义特征空间.
\begin{enumerate}[label=(\alph*)]
\item 选 $H$ 的实部最大的特征值 $\lambda$, 证明 $E\overline V(\lambda)=0$.
\item 对任意表示 $W$ 中满足 $Ew=0$ 的非零向量 $w$, 求次数为 $k$ 的多项式 $P_k$ 使 $E^kF^kw=P_k(H)w$. 先算 $EF^kw$.
\item 对 $v\in\overline V(\lambda)$, 证明存在 $N>0$ 使 $F^Nv=0$.
\item 证明 $H$ 在 $\overline V(\lambda)$ 上可对角化.
\item 对非零 $v$ 令 $N_v$ 为(c)中最小的 $N$, 证明 $\lambda=N_v-1$.
\item 证明每个 $N>0$ 都有唯一同构类的 $N$ 维不可约表示, 并在适当的基下给出 $E,F,H$ 的矩阵. 记最高权为 $\lambda$ 的这种表示为 $V_\lambda$, 其维数为 $\lambda+1$.
\item 证明Casimir算子 $C=EF+FE+H^2/2$ 与 $E,F,H$ 交换, 并在 $V_\lambda$ 上等于 $\lambda(\lambda+2)\id/2$.
\item 假设存在不能分解为较小表示直和的可约表示, 取其中维数最小者 $V$. 证明 $C$ 在 $V$ 上只有一个特征值 $\lambda(\lambda+2)/2$, 其中 $\lambda$ 是非负整数.
\item 证明 $V$ 有子表示 $W\cong V_\lambda$, 且 $V/W\cong V_\lambda^{\oplus n}$.
\item 证明 $H$ 的特征空间 $V(\lambda)$ 维数为 $n+1$. 取基 $v_1,\ldots,v_{n+1}$, 证明全部 $F^jv_i$, $1\le i\le n+1$, $0\le j\le\lambda$, 独立且构成 $V$ 的基. 原题提示: 若 $Fx=0$, $Hx=\mu x$, 则 $Cx=\mu(\mu-2)x/2$, 并从谱范围推出 $\mu=-\lambda$.
\item 令 $W_i=\operatorname{span}(v_i,Fv_i,\ldots,F^\lambda v_i)$. 证明各 $W_i$ 是子表示, 由此否定上述最小反例.
\item 给定有限维复空间上的幂零算子 $A$, 证明存在一个在表示同构意义下唯一的 $\mathfrak{sl}_2$ 表示使 $E=A$. 原题标为Jacobson--Morozov引理.
\item 求 $V_\lambda\otimes V_\mu$ 的不可约直和分解. 原题提示用 $\chi_V(x)=\tr(e^{xH})$ 及其对直和、张量积的性质.
\item 设 $V=\C^M\otimes\C^N$, $A=J_M(0)\otimes\id_N+\id_M\otimes J_N(0)$, 其中 $J_r(0)e_i=e_{i-1}$, $J_r(0)e_1=0$. 用(l),(m)求 $A$ 的Jordan形.
\end{enumerate}
\end{exercise}
\begin{solution}
(a) 由 $HE=E(H+2)$, 得 $(H-\lambda-2)^rE=E(H-\lambda)^r$. 所以 $E$ 将 $\overline V(\lambda)$ 映入 $\overline V(\lambda+2)$. 后者为0, 因为 $\lambda$ 的实部最大. 故 $E$ 在最高广义特征空间上为0. 同理由 $HF=F(H-2)$, 得 $F\overline V(\nu)\subset\overline V(\nu-2)$.

(b) 交换子展开给出
\begin{align*}
[E,F^k]&=\sum_{i=0}^{k-1}F^iHF^{k-1-i}\\
&=\sum_{i=0}^{k-1}F^{k-1}(H-2(k-1-i))
=kF^{k-1}(H-k+1).
\end{align*}
所以 $EF^kw=kF^{k-1}(H-k+1)w$. 因为 $E(H-c)w=(H-c-2)Ew=0$, 多项式作用后的向量仍被 $E$ 消去. 对 $k$ 归纳便得到
\[
P_k(H)=k!\,H(H-1)\cdots(H-k+1).
\]
这些多项式的根 $0,1,\ldots,k-1$ 均为单根.

(c) (a)末尾的降权公式给出 $F^Nv\in\overline V(\lambda-2N)$. $H$ 的谱有限, 而这些复数两两不同, 所以充分大的 $N$ 使该广义特征空间为0. 因为 $\overline V(\lambda)$ 有限维, 也可选一个 $N$ 使 $F^N$ 在整个空间上为0.

(d) 对这样的 $N$, (a),(b)给出 $P_N(H)|_{\overline V(\lambda)}=E^NF^N=0$. 一个被无重根多项式消去的算子可对角化: 若 $p(t)=\prod_i(t-c_i)$, 不同因子互素, 多项式的Bezout恒等式把空间分解为各 $\ker(H-c_i)$ 的直和. 因此 $H$ 在该空间上可对角化. 它只有特征值 $\lambda$, 故实际上是 $\lambda\id$.

(e) 对非零 $v$, 已知 $Hv=\lambda v$. 用(b)的一次升权式及 $F^{N_v}v=0$ 得
\[
0=EF^{N_v}v=N_v(\lambda-N_v+1)F^{N_v-1}v.
\]
最后一个向量非零, 且特征0, 所以 $\lambda=N_v-1\in\Z_{\ge0}$.

(f) 设 $v$ 是非零最高权向量. (e)说明向量 $v_j=F^jv$, $0\le j\le\lambda$, 均非零, 具有不同的 $H$ 特征值. 它们张成的空间被 $E,F,H$ 保持, 所以在不可约表示中就是全空间. 其作用为
\[
Hv_j=(\lambda-2j)v_j,\quad
Fv_j=v_{j+1},\quad
Ev_j=j(\lambda-j+1)v_{j-1},
\]
其中 $v_{-1}=v_{\lambda+1}=0$. 因此 $H$ 的矩阵为 $\operatorname{diag}(\lambda,\lambda-2,\ldots,-\lambda)$; $F$ 的次对角线全为1; $E$ 的第 $j$ 列在第 $j-1$ 行有系数 $j(\lambda-j+1)$, 其余为0, 这里列行从0编号.

反过来, 对任意非负整数 $\lambda$ 以这些公式定义算子, 在每个 $v_j$ 上计算可验证三个交换关系, 因而确实定义表示. 非零子表示经 $H$ 的多项式特征投影包含某个 $v_j$; 连续作用 $E$ 得到 $v_0$, 因为中间系数均非零; 再作用 $F$ 得到全部基. 所以表示不可约. 公式也说明同一维数的不可约表示彼此同构. $N$ 维者对应 $\lambda=N-1$.

(g) 使用 $[AB,D]=A[B,D]+[A,D]B$, 有
\begin{align*}
[EF+FE,E]&=-EH-HE,&[H^2/2,E]&=HE+EH,\\
[EF+FE,F]&=HF+FH,&[H^2/2,F]&=-HF-FH.
\end{align*}
所以 $[C,E]=[C,F]=0$. 直接由权关系得到 $[C,H]=0$. 在最高权向量 $v_0$ 上, $FEv_0=0$, $EFv_0=\lambda v_0$, 所以 $Cv_0=\lambda(\lambda+2)v_0/2$. 因为 $C$ 与 $F$ 交换, 同一标量作用于全部 $F^jv_0$, 得到结论.

(h) 最小反例不可分解; 否则每个较小直和项依维数归纳都完全可约, 其直和也完全可约. $C$ 与表示作用交换, 所以其广义特征空间是子表示. 不可分解性迫使 $C$ 只有一个特征值 $c$. 任取不可约子表示 $W$, 它由(f)同构于某个 $V_\lambda$, 所以由(g)有 $c=\lambda(\lambda+2)/2$.

(i) 此 $W\cong V_\lambda$ 是真子表示. 商 $V/W$ 维数更小, 由最小性完全可约. 因为 $C$ 的唯一特征值在商中仍为 $c$, 每个不可约商分量 $V_\nu$ 满足 $\nu(\nu+2)=\lambda(\lambda+2)$. $\nu,\lambda\ge0$, 所以 $\nu=\lambda$, 得 $V/W\cong V_\lambda^{\oplus n}$, 其中 $n\ge1$.

(j) 取基使 $W$ 位于前面, 则 $H$ 的矩阵为分块上三角, 其特征多项式是 $W$ 和商的特征多项式之积. 因此谱恰在 $\lambda,\lambda-2,\ldots,-\lambda$ 中, 每个的代数重数为 $n+1$. 最高实部的广义特征空间 $\overline V(\lambda)$ 由(d)已经是特征空间, 故 $\dim V(\lambda)=n+1$.

固定 $0\le j\le\lambda$. 若 $\sum_i c_iF^jv_i=0$, 作用 $E^j$, (b)给出
\[
j!\lambda(\lambda-1)\cdots(\lambda-j+1)\sum_i c_i v_i=0.
\]
前面的系数非零, 所以所有 $c_i=0$. 不同 $j$ 的向量位于不同 $H$ 特征空间, 因而全部独立. 它们共有 $(n+1)(\lambda+1)=\dim V$ 个, 所以是基.

原题提示中的底权计算也成立: 若 $Fx=0$, $Hx=\mu x$, 则 $EFx=0$, $FEx=EFx-Hx=-\mu x$, 所以 $Cx=\mu(\mu-2)x/2$. 比较唯一特征值得到 $\mu=-\lambda$ 或 $\mu=\lambda+2$. 后一种不在刚求出的 $H$ 的谱中, 因而必须舍去. 不能只从二次方程未经谱检查就断言 $\mu=-\lambda$.

(k) 由(e), 每个最高权向量满足 $F^{\lambda+1}v_i=0$, 再由(f)的作用公式, $W_i$ 被 $E,F,H$ 保持, 且 $W_i\cong V_\lambda$. (j)的基给出 $V=\bigoplus_{i=1}^{n+1}W_i$, 与最小反例矛盾. 因此所有有限维复 $\mathfrak{sl}_2$ 表示均完全可约. 旧讲义(k)把末句中的 $W_i$ 误写作 $V_i$, 这里统一使用已定义的 $W_i$.

(l) 在 $V_\lambda$ 上, $E$ 的每个相邻升权系数均非零, 所以适当缩放基后, $E$ 是一个大小 $\lambda+1$ 的幂零Jordan块. 给定 $A$ 的Jordan块大小 $r_1,\ldots,r_s$, 取表示 $\bigoplus_iV_{r_i-1}$; 其中的 $E$ 与 $A$ 具有相同Jordan形. 用一个相似变换把此表示搬到给定空间, 使 $E$ 恰等于 $A$, 得到存在性.

任意满足 $E=A$ 的表示由(k)完全可约. 各不可约分量的维数恰是 $A$ 的Jordan块大小; Jordan形的唯一性因此决定表示的同构类. 这里的唯一性是在允许同时共轭 $E,F,H$ 的意义下, 并非固定 $E$ 后 $F,H$ 作为具体矩阵逐项唯一.

(m) $H$ 在 $V_\lambda$ 上的权为 $\lambda,\lambda-2,\ldots,-\lambda$. 令 $q=e^x$, 则
\[
\chi_{V_\lambda}(x)=q^\lambda+q^{\lambda-2}+\cdots+q^{-\lambda}.
\]
对直和, 迹相加; 对张量积, $H=H_V\otimes\id+\id\otimes H_W$, 两项交换, 所以 $e^{xH}=e^{xH_V}\otimes e^{xH_W}$, 迹相乘. 假设 $\lambda\le\mu$. 乘积中权 $\lambda+\mu-2r$ ($0\le r\le\lambda+\mu$)的系数是满足 $i+j=r$, $0\le i\le\lambda$, $0\le j\le\mu$ 的数目. 另一方面, 在
\[
\sum_{k=0}^{\lambda}
(q^{\lambda+\mu-2k}+q^{\lambda+\mu-2k-2}+\cdots+q^{-\lambda-\mu+2k})
\]
中该系数为 $\min(\lambda,r,\lambda+\mu-r)+1$, 正是前述计数. 两个Laurent多项式相等. 不同 $V_r$ 的特征标线性无关, 因为最大权项的系数可逐次读出. 结合(k)的完全可约性, 得
\[
V_\lambda\otimes V_\mu\cong
\bigoplus_{k=0}^{\min(\lambda,\mu)}V_{\lambda+\mu-2k}.
\]

(n) 由(l)把 $J_M(0)$ 与 $J_N(0)$ 分别看作 $V_{M-1}$ 与 $V_{N-1}$ 的 $E$. 张量积的 $E$ 正是题中的 $A$. 按(m)分解后, 每个 $V_r$ 上的 $E$ 是单个大小 $r+1$ 的块, 因而
\[
A\sim\bigoplus_{k=0}^{\min(M,N)-1}J_{M+N-1-2k}(0).
\]
各块大小从 $M+N-1$ 以步长2降到 $|M-N|+1$. 它们之和为 $MN$, 与原空间维数一致.
\end{solution}

\originalsection{Homework 4: 可解李代数与生成关系}

\begin{exercise}\label{hw:1.56}
MIT 18.712 (2010), Homework 4, 原题1.56. 李代数的导代数 $K(\mathfrak g)=[\mathfrak g,\mathfrak g]$ 是所有 $[x,y]$ 的线性包, 它是理想. 递归令 $K^0(\mathfrak g)=\mathfrak g$, $K^{r+1}(\mathfrak g)=[K^r(\mathfrak g),K^r(\mathfrak g)]$. 若某个 $K^r(\mathfrak g)=0$, 则称 $\mathfrak g$ 可解. 证明Lie定理: 有限维复可解李代数的任一有限维不可约表示均为一维. 原题提示对李代数维数归纳, 并用迹证明导代数的共同特征空间被 $\mathfrak g$ 保持.
\end{exercise}
\begin{solution}
Jacobi恒等式给出 $[z,[x,y]]=[[z,x],y]+[x,[z,y]]$, 所以 $K(\mathfrak g)$ 是理想. 非零可解李代数的导代数 $\mathfrak a=K(\mathfrak g)$ 是真子代数, 否则导出列永不变为0. 对 $\dim\mathfrak g$ 归纳, 同时证明其每个非零有限维表示有共同特征向量. 当 $\dim\mathfrak g=0$ 时任取非零向量即可. 对较小的 $\mathfrak a$, 归纳给出 $0\ne v\in V$ 和线性泛函 $\chi:\mathfrak a\to\C$ 使 $av=\chi(a)v$.

固定 $x\in\mathfrak g$, 取 $U=\operatorname{span}(v,xv,x^2v,\ldots)$, 这是有限维的 $x$-不变空间. 令 $U_j=\operatorname{span}(v,\ldots,x^jv)$. 对 $a\in\mathfrak a$, 归纳得到
\[
a x^jv\equiv\chi(a)x^jv\pmod{U_{j-1}}.
\]
事实上 $ax^j=x(ax^{j-1})+[a,x]x^{j-1}$, 且 $[a,x]\in\mathfrak a$, 因而第二项落在 $U_{j-1}$, 第一项的最高项由归纳确定. 这也证明每个 $U_j$ 被 $\mathfrak a$ 保持. 取第一个线性相关前的 $v,xv,\ldots,x^{d-1}v$ 为 $U$ 的基, 则 $a|_U$ 的对角线均为 $\chi(a)$, 故 $\tr_U(a)=d\chi(a)$.

特别地, $[x,a]\in\mathfrak a$, 其在 $U$ 上的迹为 $d\chi([x,a])$. 又因为 $x,a$ 均保持 $U$, 交换子矩阵的迹为0. 特征0使 $d\ne0$, 所以 $\chi([x,a])=0$.

令 $V_\chi=\{w:aw=\chi(a)w\text{ 对所有 }a\in\mathfrak a\}$, 它非零. 对 $w\in V_\chi$, 有
\[
a(xw)=x(aw)+[a,x]w=\chi(a)xw+\chi([a,x])w=\chi(a)xw.
\]
因此 $V_\chi$ 被 $\mathfrak g$ 保持. 在该空间上, 每个 $[x,y]\in\mathfrak a$ 作用为标量, 其迹又是交换子的迹0, 所以该标量为0. 从而 $\mathfrak g$ 在 $V_\chi$ 上的所有算子彼此交换. 有限个交换复矩阵有共同特征向量: 逐次取一个算子的非零特征空间, 其余算子保持它, 直至遍历 $\mathfrak g$ 的一组有限基即可. 这完成归纳中的较强结论. 若原表示不可约, 共同特征向量张成的非零一维子表示就是全空间, 故 $\dim V=1$.
\end{solution}

\begin{exercise}\label{hw:1.57}
MIT 18.712 (2010), Homework 4, 原题1.57. 对基为 $X,Y$、关系为 $[X,Y]=Y$ 的二维李代数, 分别在特征0与正特征中分类有限维不可约表示. 正特征中Lie定理是否仍成立? 底域依旧讲义总约定代数闭.
\end{exercise}
\begin{solution}
在特征0中, 导代数为 $kY$, 再次取导代数为0, 所以它可解. 上一题的证明适用于任意代数闭特征0域: 只使用特征值的存在和正整数可逆, 未使用复数的分析性质. 不可约表示均一维. 一维中交换子为0, 所以 $Y$ 作用为0, 而 $X$ 作用为任意标量 $a\in k$. 这些 $a$ 给出互不同构的全部表示.

以下令 $\operatorname{char}k=p$. 表示中的同名算子满足 $XY-YX=Y$. 由此 $Y^p$ 与 $X,Y$ 都交换, 以及 $X^p-X$ 与 $X,Y$ 都交换. 后一断言可由 $XY=Y(X+1)$ 推出: $X^pY=Y(X+1)^p=Y(X^p+1)$, 故 $[X^p-X,Y]=0$. 由有限维Schur引理,
\[
Y^p=b\id,\qquad X^p-X=c\id.
\]
若 $b=0$, 则 $Y$ 幂零, $\ker Y\ne0$. 由 $YX=XY-Y$, 此核也被 $X$ 保持. 不可约性给出 $\ker Y=V$, 即 $Y=0$. 剩下单个算子 $X$ 的不可约有限维表示必为一维, 与特征0的族相同.

若 $b\ne0$, 则 $Y$ 可逆. 多项式 $t^p-t-c$ 的导数为 $-1$, 没有重根, 故 $X$ 可对角化. 取 $Xv=av$, 则
\[
XY^iv=(a+i)Y^iv,\qquad 0\le i<p.
\]
这 $p$ 个特征值两两不同, 相应向量独立. 它们的张成空间被 $X,Y$ 保持, 因为 $Y^pv=bv$, 所以它就是 $V$. 全部这类表示因此是 $p$ 维, 具有基 $e_0,\ldots,e_{p-1}$ 和作用
\[
Xe_i=(a+i)e_i,\qquad Ye_i=e_{i+1}\ (i<p-1),
\qquad Ye_{p-1}=be_0.
\]
其中 $b\in k^\times$. 这些公式满足 $[X,Y]=Y$, 包括最后一个基向量上的关系, 因为 $p=0$ 于 $k$. 非零不变子空间由 $X$ 的特征投影包含某个 $e_i$, 再用可逆 $Y$ 得到全部基, 故表示不可约.

参数 $a$ 可以沿循环改为 $a+i$, $i\in\mathbb F_p$, 而不改变同构类. 完整的同构参数为 $(c,b)=(a^p-a,b)\in k\times k^\times$: 多项式 $t^p-t-c$ 的根恰为 $a+\mathbb F_p$, 所以相同的中心参数给出同构表示, 不同参数则不能同构. 可解的二维李代数已有 $p>1$ 维不可约表示, 因而Lie定理在正特征中不成立. 这里所分类的是一般李代数表示, 未另加限制性表示条件 $Y^p=0$.
\end{solution}

\begin{exercise}\label{hw:1.58}
MIT 18.712 (2010), Homework 4, 原题1.58, 官方标为难题, 作业中为选做. 对李代数元素 $x$, 记 $\operatorname{ad}(x)(z)=[x,z]$. 令 $\mathfrak g_n$ 由两个生成元 $x,y$ 和定义关系
\[
(\operatorname{ad}x)^2(y)=0,\qquad
(\operatorname{ad}y)^{n+1}(x)=0
\]
生成.
\begin{enumerate}[label=(\alph*)]
\item 证明 $\mathfrak g_1,\mathfrak g_2,\mathfrak g_3$ 有限维, 并求维数.
\item 证明 $\mathfrak g_4$ 无限维, 并显式构造它的一组基.
\end{enumerate}
原题未在此处指定底域. 以下证明给出复数底域的版本, 并在末尾说明正特征下原陈述的条件问题. (b)中的全空间基采用确定的有限消元算法形式, 而不调用仿射李代数分类或未证明的表示完备性定理.
\end{exercise}
\begin{solution}
(a) 令 $D=\operatorname{ad}y$, $x_i=D^ix$. Jacobi恒等式说明 $D$ 是导子: $D[u,v]=[Du,v]+[u,Dv]$. 两条定义关系等价于 $[x_0,x_1]=0$ 和 $x_{n+1}=0$.

当 $n=1$ 时, $x_0,x_1$ 交换, $Dx_0=x_1$, $Dx_1=0$, 所以 $y,x_0,x_1$ 张成一个对括号封闭的空间, 包含两个生成元, 因而张成整个 $\mathfrak g_1$. 反过来, 在三个独立基向量上定义唯一非零基本括号 $[y,x_0]=x_1$, 便得到满足原关系且由 $x_0,y$ 生成的三维李代数. 因此原代数维数既不大于也不小于3, 恰为3.

当 $n=2$ 时, 依次对 $[x_0,x_1]=0$ 作用 $D$, 得到 $[x_0,x_2]=0$, 再作用 $D$ 得 $[x_1,x_2]=0$. 因而 $y,x_0,x_1,x_2$ 张成全代数. 以 $x_0,x_1,x_2$ 为交换的三维空间, 让 $D$ 依次把 $x_0$ 送到 $x_1$, $x_1$ 送到 $x_2$, $x_2$ 送到0, 构成与 $ky$ 的半直积; 这给出满足原关系的四维模型. 所以 $\dim\mathfrak g_2=4$.

当 $n=3$ 时, 从 $[x_0,x_1]=0$ 连续取导, 利用 $x_4=0$, 得
\[
[x_0,x_2]=0,\quad [x_1,x_2]+[x_0,x_3]=0,
\quad 2[x_1,x_3]=0,\quad 2[x_2,x_3]=0.
\]
在特征0下可除以2. 令 $z=[x_0,x_3]=-[x_1,x_2]$. $z$ 与全部 $x_i$ 交换: 对 $x_0,x_3$ 用表达式 $-[x_1,x_2]$ 和已得零括号, 对 $x_1,x_2$ 用表达式 $[x_0,x_3]$ 和Jacobi恒等式. 另外 $Dz=[x_1,x_3]=0$. 因而
\[
y,x_0,x_1,x_2,x_3,z
\]
张成全代数. 为证明独立, 先以 $x_0,x_1,x_2,x_3,z$ 为基定义一个两步幂零李代数, 唯一非零基本括号为 $[x_0,x_3]=z$, $[x_1,x_2]=-z$, $z$ 中心. 再令 $Dx_i=x_{i+1}$, $Dx_3=Dz=0$. 逐对检查可知 $D$ 保持这些括号: 唯一可能产生非零和的一对是 $(x_0,x_2)$, 其导数括号和为 $[x_1,x_2]+[x_0,x_3]=-z+z=0$; 其余直接为0. 因此 $D$ 是导子, 与 $ky$ 的半直积是六维李代数, 满足原关系且由 $x_0,y$ 生成. 得到 $\dim\mathfrak g_3=6$.

(b) 先以具体矩阵证明无限维, 再构造覆盖全空间的基. 在 $\mathfrak{sl}_3(\C[t])$ 中令 $e=E_{12}+E_{23}$, $P_0=E_{31}$, $P_i=(\operatorname{ad}e)^iP_0$. 用 $E_{ij}E_{rs}=\delta_{jr}E_{is}$ 计算得到
\begin{align*}
P_1&=E_{21}-E_{32},&P_2&=\operatorname{diag}(1,-2,1),\\
P_3&=3(E_{23}-E_{12}),&P_4&=6E_{13},\qquad P_5=0.
\end{align*}
并且 $[P_0,P_1]=0$. 所以 $x\mapsto tP_0$, $y\mapsto e$ 满足两条定义关系, 给出李代数同态 $\Psi:\mathfrak g_4\to\mathfrak{sl}_3(\C[t])$.

为在像中产生无限多个独立元素, 递归定义
\[
a_1=x,\qquad
a_{m+2}=\frac16
\Bigl[\bigl[[y,x],[y,[y,a_m]]\bigr],[y,x]\Bigr]
\quad(m=1,3,5,\ldots).
\]
设 $Q_0=E_{21}+E_{32}$. 直接矩阵计算给出 $[P_1,P_2]=3Q_0$, $[Q_0,P_1]=2P_0$. 若 $\Psi(a_m)=t^mP_0$, 则递推式的像为
\[
\frac16\bigl[[tP_1,t^mP_2],tP_1\bigr]
=\frac16\,t^{m+2}[3Q_0,P_1]=t^{m+2}P_0.
\]
归纳知全部奇数 $m$ 的像是 $t^mE_{31}$. 不同 $t$ 次数的矩阵多项式独立, 故对应的 $a_m$ 在原代数中也独立. 因而 $\mathfrak g_4$ 无限维. 此论证只使用同态的存在, 没有把该矩阵像擅自等同于整个原代数.

下面给出一组真正覆盖原代数的显式算法基. 为使选择完全确定, 对所有二叉括号词使用字符串编码: $x,y$ 为长度1词, 若 $u,v$ 已编码, 则 $[u,v]$ 的编码就是左方括号、$u$、逗号、$v$、右方括号的连接. 按ASCII字符串字典序排列, 即字符顺序为逗号、左方括号、右方括号、$x$、$y$. 记 $W_d$ 为具有 $d$ 个叶子的全部这种词, $F_d$ 为以 $W_d$ 为基的有理向量空间. 每个 $W_d$ 是有限集, 递归由
\[
W_1=\{x,y\},\qquad
W_d=\bigcup_{a=1}^{d-1}\{[u,v]:u\in W_a,\ v\in W_{d-a}\}
\]
生成. 这里尚不使用反交换或Jacobi关系; $[u,v]$ 与 $[v,u]$ 作为形式词是不同基向量.

逐个次数建立以 $W_d$ 为列的有理关系矩阵 $R_d$. 行的加入规则如下:
\begin{enumerate}[label=\arabic*.]
\item 加入所有本次数的反交换关系 $[u,v]+[v,u]$.
\item 加入所有本次数的Jacobi关系 $[u,[v,w]]+[v,[w,u]]+[w,[u,v]]$.
\item 次数3时加入 $[x,[x,y]]$; 次数6时加入 $[y,[y,[y,[y,[y,x]]]]]$.
\item 对全部 $a<d$, 取已计算的 $R_a$ 的非零行最简阶梯形作为关系空间的固定基. 对其每行 $r$ 和每个 $u\in W_{d-a}$, 加入 $[r,u]$、$[u,r]$; 这里把括号按线性展开为 $W_d$ 的坐标.
\end{enumerate}
所有枚举都有限, 只需整数和有理数运算. 将 $R_d$ 化为最简行阶梯形, 主元取当前最左的非零列. 定义 $B_d$ 为非主元列所对应的括号词的余类. 则
\[
\mathcal B=\bigcup_{d\ge1}B_d
\]
就是所构造的基. 它的定义不涉及选择未知子空间、不以猜测维数为停止准则; 给定任一括号词, 按其次数作有限计算便能判断它是否属于 $\mathcal B$, 并算出其在 $B_d$ 上的坐标.

证明完备性如下. 令 $T_d$ 为 $R_d$ 的行空间, $T=\bigoplus_dT_d$, $F=\bigoplus_dF_d$. 第4步保证 $[T,F]\subset T$, $[F,T]\subset T$. 前两步保证 $F/T$ 的括号反交换且满足Jacobi; 在特征0中, 反交换也给出 $[u,u]=0$. 第3步给出两条定义关系. 反过来, 以上每一行都由这些反交换、Jacobi和定义关系以及取括号生成, 所以 $T$ 正是这些关系在自由非结合括号空间中生成的理想. 因此 $(F/T)\otimes_\mathbb Q\C$ 具有原定义代数的普遍性质: 任何两个满足原关系的李代数元素都唯一决定从它出发的同态. 于是它就是 $\mathfrak g_4$.

有限矩阵的行最简阶梯形中, 主元列的向量可由关系表达为非主元列的组合, 而非主元列之间不存在非零关系. 所以 $B_d$ 是每个次数商空间的基. 两条定义关系都是齐次的, 各次数不混合, 因而它们的并集是全代数的基. 扩充标量到 $\C$ 保持这些基和关系的独立性.

为给出可直接核算的开头几项, 令 $v=[x,D^3x]$. 各次数可另取如下基:
\begin{center}
\begin{tabular}{ccl}
\toprule
总次数 $d$ & 维数 & 一组基\\
\midrule
1&2&$x,y$\\
2&1&$Dx$\\
3&1&$D^2x$\\
4&1&$D^3x$\\
5&2&$v,D^4x$\\
6&1&$Dv$\\
7&2&$[x,Dv],D^2v$\\
\bottomrule
\end{tabular}
\end{center}
表格给出有限次数的一组基; 全空间算法基的正确性由前面的理想和消元证明保证. 此种显式基是算法形式, 计算量随次数很快增加, 不把它声称为闭式的仿射根基.

以下不依赖消元程序, 直接证明表中各组向量是基. 令 $x_i=D^ix$, 此时 $x_5=0$. 从 $[x_0,x_1]=0$ 及其前四次导数得到
\begin{align*}
[x_0,x_2]&=0,&[x_1,x_2]&=-[x_0,x_3]=-v,\\
2[x_1,x_3]+[x_0,x_4]&=0,&
2[x_2,x_3]+3[x_1,x_4]&=0.
\end{align*}
Jacobi还给出 $[x_0,v]=-[x_0,[x_1,x_2]]=0$, 其导数为 $[x_1,v]=-[x_0,Dv]$. 又
\[
Dv=\tfrac12[x_0,x_4],\qquad
D^2v=\tfrac12[x_1,x_4],\qquad [x_2,x_3]=-3D^2v.
\]
从次数1逐次取括号, 上述关系把次数2至7的全部词分别约化到表列的张成集. 原因是每个括号词最外层都分成两个次数更小的词; 用已经列出的较低次数张成集代入, 这些等式涵盖总次数至7的所有分拆. 矩阵同态则给出相应非零像:
\[
\Psi(x_i)=tP_i,\quad
\Psi(v)=-3t^2Q_0,\quad
\Psi(Dv)=-3t^2\operatorname{diag}(1,0,-1),
\]
以及 $\Psi(D^2v)=3t^2e$, $\Psi([x,Dv])=-6t^3P_0$. 在每个列出两个词的次数中, 它们的像有不同 $t$ 次数, 因而独立; 其余次数的单个像非零. 这从上界和下界两方面证明了表中的初始维数, 没有从数值实验外推无限次数的结论.

最后核对底域遗漏. (a)对 $n=3$ 的证明用到2可逆. 若把原题按全局约定解释为任意代数闭域, 则特征2时其有限维断言确实不成立. 在特征2的 $\mathfrak{sl}_3(k[t])$ 中仍取 $x\mapsto tE_{31}$, $y\mapsto e$. 此时上述 $P_3=e$, $P_4=0$, 所以满足 $\mathfrak g_3$ 的关系. 像中含 $te$; 若已含 $t^m e$, 则
\[
[tE_{31},t^me]=t^{m+1}(E_{21}+E_{32}),\qquad
(\operatorname{ad}e)^2(E_{21}+E_{32})=e.
\]
因而又含 $t^{m+1}e$. 这些矩阵多项式独立, 证明特征2下 $\mathfrak g_3$ 无限维. 所以本题的复数版本完整成立, 而原文省略的特征条件不可直接删除; 本题解不声称分类其他正特征下的这些生成代数.
\end{solution}

% MIT 18.712 Fall 2010, Homework 5--9. 中文题面与编者独立解答.
% 本文件只用主文件已定义的宏; 题源页码均为109页旧PDF的实际页码.
\originalsection{Homework 5: 扩张与代数表示}

以下解答均为编者独立推导, 并非官方答案的翻译. 本节的代数带单位, 表示保持单位. 涉及不可约表示的标量自同态以及射影空间分类时, 底域取代数闭域 $k$; 题目另有指定时服从题目条件. 扩张是一个正合列 $0\to V\to U\to W\to0$, 因而同时记录子表示和商表示的指定识别.

\begin{exercise}\label{hw:2.21}
\textbf{MIT 18.712 (2010), Homework 5, 原题 2.21; 旧PDF第30页.}
设 $A$ 为代数, $V,W$ 为其表示. 利用一个向量空间分裂将扩张 $U$ 写成 $V\oplus W$, 并令
\[
 \rho_U(a)=\begin{pmatrix}\rho_V(a)&f(a)\\0&\rho_W(a)\end{pmatrix},
 \qquad f:A\to\Hom_k(W,V)\text{ 线性}.
\]
(a) 求该矩阵定义表示的充要条件; 满足条件的映射组成空间 $Z^1(W,V)$.
(b) 对 $X:W\to V$, 令 $dX(a)=\rho_V(a)X-X\rho_W(a)$. 证明 $dX$ 为余循环, 且 $dX=0$ 当且仅当 $X$ 为表示同态. 由此说明余边界空间 $B^1(W,V)$ 同构于 $\Hom_k(W,V)/\Hom_A(W,V)$, 并定义 $\operatorname{Ext}^1(W,V)=Z^1(W,V)/B^1(W,V)$.
(c) 证明相差余边界的两个余循环给出同构的扩张; 反之, 若同构在子表示和商表示上均为恒等, 则两个余循环相差余边界.
(d) 若 $V,W$ 有限维不可约, 证明两个非零扩张类给出的中间表示同构, 当且仅当两类相差非零标量. 因而非平凡扩张的中间表示同构类由 $\mathbb P\operatorname{Ext}^1(W,V)$ 参数化; 所有扩张均平凡当且仅当此空间为零.
\end{exercise}
\begin{solution}
(a) 计算乘积的右上角得到必要且充分的条件
\[
 f(ab)=\rho_V(a)f(b)+f(a)\rho_W(b).                 \tag{H5.1}
\]
代入 $a=b=1$ 可得 $f(1)=2f(1)$, 所以 $f(1)=0$. 因而乘法条件也保证保持单位. 条件是关于 $f$ 的齐次线性条件, 所以余循环组成向量空间.

(b) 直接展开有
\[
\begin{aligned}
 dX(ab)&=\rho_V(a)\rho_V(b)X-X\rho_W(a)\rho_W(b)\\
 &=\rho_V(a)dX(b)+dX(a)\rho_W(b).
\end{aligned}
\]
故 $dX$ 是余循环. $dX=0$ 恰是 $X$ 与每个 $a$ 的作用交换. 线性映射 $d$ 的核为 $\Hom_A(W,V)$, 第一同构定理给出所要求的同构.

(c) 写 $T=\begin{pmatrix}1&X\\0&1\end{pmatrix}$. 条件 $T\rho_f(a)=\rho_{f'}(a)T$ 等价于 $f(a)-f'(a)=dX(a)$. 因此这样的 $T$ 给出扩张同构, 而任何在两端为恒等的扩张同构都必为此形式. 原题这一小问末尾把 $B^1(W,V)$ 写成 $B^1(V,W)$; 由映射的定义可见应为前者.

(d) 先注意非分裂扩张 $U$ 中的不可约子表示只能是指定的 $V$. 若另一个不可约子表示 $L$ 未包含在 $V$ 内, 则商映射 $L\to W$ 非零, 从而为同构; 其逆给出 $W\to U$ 的分裂, 矛盾. 若 $L\subset V$, 则 $L=V$. 因而任意中间表示同构 $U_f\to U_{f'}$ 都保持这个子表示, 并在子表示和商表示上分别为非零标量 $\alpha,\beta$. 它的矩阵为 $\begin{pmatrix}\alpha&X\\0&\beta\end{pmatrix}$, 交换条件给出
\[
 \alpha f-\beta f'=dX.
\]
在 $\operatorname{Ext}^1$ 中即为 $[f']=(\alpha/\beta)[f]$. 反过来, 此关系可选取 $X$ 实现, 所得矩阵可逆. 零类恰好对应分裂扩张. 非分裂中间表示不可能与 $V\oplus W$ 同构, 因为后者完全可约, 其中每个子表示都有补空间. 这证明射影空间分类及最后的等价关系. 代数闭域条件在将两端的自同构化为标量时使用; 一般域不能不加修改地沿用该射影空间结论.
\end{solution}

\begin{exercise}\label{hw:2.22}
\textbf{MIT 18.712 (2010), Homework 5, 原题 2.22; 旧PDF第30页.}
(a) 令 $A=\C[x_1,\ldots,x_n]$, $V_a,V_b$ 为一维表示, $x_i$ 分别作用为 $a_i,b_i$. 计算 $\operatorname{Ext}^1(V_a,V_b)$ 并分类 $A$ 的二维表示.
(b) 令复代数 $B$ 由 $x_1,\ldots,x_n$ 生成, 关系为所有 $x_ix_j=0$. 证明当 $n>1$ 时, $B$ 有无穷多个两两不同构的不可分解表示.
\end{exercise}
\begin{solution}
(a) 对扩张 $0\to V_b\to U\to V_a\to0$, 写
\[
 \rho(x_i)=\begin{pmatrix}b_i&t_i\\0&a_i\end{pmatrix}.
\]
生成元两两交换当且仅当 $(b_i-a_i)t_j=(b_j-a_j)t_i$. 如果 $a\ne b$, 选一个 $b_i-a_i\ne0$, 则所有 $t_j$ 为同一个标量乘 $b_j-a_j$, 正是余边界, 所以 $\operatorname{Ext}^1(V_a,V_b)=0$. 如果 $a=b$, 交换条件恒成立且余边界为零, 所以 $\operatorname{Ext}^1(V_a,V_a)\cong\C^n$.

一族交换的复矩阵有公共特征向量: 先取一个矩阵的非零特征子空间, 其余矩阵保持它, 再在该空间重复, 有限维数保证最终取得公共特征向量. 在二维情形其张成的一维空间给出上述扩张. 因而所有二维表示是 $V_a\oplus V_b$ (无序参数对 $\{a,b\}$), 或者
\[
 \rho(x_i)=a_i I+t_i\begin{pmatrix}0&1\\0&0\end{pmatrix},
 \qquad t\ne0.
\]
后一类由 $a\in\C^n$ 与 $[t]\in\mathbb P^{n-1}$ 唯一参数化. 它不可分解: 若分解为两条不变直线, 每个 $x_i$ 在该分解下都是对角矩阵; 但某个 $t_i\ne0$ 的矩阵有非平凡Jordan块. $t=0$ 则已经属于 $V_a\oplus V_a$.

(b) 令 $\rho(x_i)=t_iN$, 其中 $N=\begin{pmatrix}0&1\\0&0\end{pmatrix}$, $t\ne0$. 因为 $N^2=0$, 所有关系均成立. 同上这些表示不可分解, 而同构类由 $[t]\in\mathbb P^{n-1}$ 区分. 当 $n>1$ 时, $[1:s:0:\cdots:0]$, $s\in\C$, 已给出无穷多个不同的类.
\end{solution}

\begin{exercise}\label{hw:2.23}
\textbf{MIT 18.712 (2010), Homework 5, 原题 2.23; 旧PDF第30页.}
设 $Q$ 为有限且没有有向圈的箭图, $P_Q$ 为其路径代数. 求所有不可约表示及其间的 $\operatorname{Ext}^1$, 并分类总维数为2的表示.
\end{exercise}
\begin{solution}
路径代数以所有有向路径为基, 长度零的路径为顶点幂等元 $e_i$. 采用 $e_j\alpha e_i=\alpha$ 的约定, 其中 $\alpha:i\to j$; 路径的乘法对应线性映射的复合. 一个表示等价于给各顶点向量空间 $M_i=e_iM$, 给各箭头线性映射 $M_i\to M_j$. 总维数是 $\sum_i\dim M_i$.

正长度路径张成的理想 $J$ 是幂零理想, 因为无有向圈的路径长度小于顶点数. 若 $M$ 为简单模, $JM$ 是子模. 若 $JM=M$, 则反复乘 $J$ 得 $J^rM=M$, 与幂零性矛盾; 所以 $JM=0$. 因为 $P_Q/J\cong\prod_i k$, 不可约表示恰为顶点表示 $S_i$: 顶点 $i$ 放 $k$, 其余顶点放零, 箭头全为零.

扩张 $0\to S_i\to U\to S_j\to0$ 在各顶点均为向量空间分裂. 若 $i\ne j$, 只有从 $j$ 到 $i$ 的箭头能够有非零作用, 每条这样的箭头给一个任意标量; 其他箭头为零. 没有需要再施加的关系, 因为路径代数由这些箭头自由拼接生成. 固定两端的扩张同构在各顶点只能是恒等, 也没有能改变这些标量的跨顶点线性映射. 故
\[
 \operatorname{Ext}^1(S_j,S_i)\cong k^{r_{ji}},
 \qquad r_{ji}=\#\{\text{箭头 }j\to i\}.
\]
若 $i=j$, 顶点处只有二维空间且没有自环, 所有箭头作用均为零, 所以扩张分裂; 公式仍成立, 因为 $r_{ii}=0$.

总维数为2时, 或仅有一个顶点放二维空间, 得 $S_i\oplus S_i$; 或两个不同顶点各放一维空间. 后一种情形, 所有连向零空间的箭头均为零. 无有向圈保证这两个顶点之间不可能同时有两个方向的箭头. 全部箭头标量为零时得 $S_i\oplus S_j$; 若 $j\to i$ 有 $r_{ji}$ 条箭头且参数向量 $t\ne0$, 改变两顶点的基只把 $t$ 乘共同的非零标量. 因而非平凡类由 $\mathbb P^{r_{ji}-1}(k)$ 参数化. 至少一条箭头非零时表示不可分解, 因为可能的非零直和分解只能把两个顶点分开, 而这样会迫使全部箭头为零. 这也说明没有漏掉二维表示.
\end{solution}

\begin{exercise}\label{hw:2.24}
\textbf{MIT 18.712 (2010), Homework 5, 原题 2.24; 旧PDF第31页.}
设 $\rho:A\to\End(V)$. 一个形式变形是 $\widetilde\rho=\rho+t\rho_1+t^2\rho_2+\cdots$, 满足 $\widetilde\rho(ab)=\widetilde\rho(a)\widetilde\rho(b)$. 以形式可逆算子 $b(t)=1+tb_1+t^2b_2+\cdots$ 共轭得到的变形称为同构.
(a) 证明若 $\operatorname{Ext}^1(V,V)=0$, 所有形式变形均同构于 $\rho$.
(b) 逆命题是否成立? 考虑双数代数 $k[x]/(x^2)$.
\end{exercise}
\begin{solution}
(a) 假设已通过共轭消去次数小于 $r$ 的项, 则当前变形为 $\rho+t^r\eta+O(t^{r+1})$. 比较乘法公式中 $t^r$ 的系数, 得
\[
 \eta(ab)=\rho(a)\eta(b)+\eta(a)\rho(b).
\]
因此 $\eta$ 为余循环, 假设给出 $X_r$ 使 $\eta(a)=\rho(a)X_r-X_r\rho(a)$. 用 $1+t^rX_r$ 共轭, 新的 $t^r$ 项为 $\eta+X_r\rho-\rho X_r=0$, 更低次项不变. 从 $r=1$ 开始归纳. 这些共轭算子的有序乘积在形式幂级数意义下存在: 任意固定次数的系数只依赖有限多个因子. 其常数项为1, 所以可逆, 并把整个变形化为 $\rho$. 此处是形式逐系数论证, 不需要解析收敛.

(b) 不成立. 取一维表示 $V=k$, 令 $x$ 作用为零. 一阶扩张中 $x$ 可作用为 $\begin{pmatrix}0&t_0\\0&0\end{pmatrix}$, 任意 $t_0\in k$, 所以 $\operatorname{Ext}^1(V,V)\cong k\ne0$. 但在一维形式变形中, $x$ 必作用为 $a(t)\in tk[[t]]$, 且关系要求 $a(t)^2=0$. 幂级数环 $k[[t]]$ 是整环, 故 $a(t)=0$. 因而所有形式变形仍平凡. 一阶可变方向未必能够延伸为全部阶数的变形.
\end{solution}

\begin{exercise}\label{hw:2.25}
\textbf{MIT 18.712 (2010), Homework 5, 原题 2.25; 旧PDF第31页.}
设有限维复向量空间 $V$ 带对称双线性型 $(\ ,\ )$. 定义Clifford代数
\[
 \operatorname{Cl}(V)=T(V)/(v\otimes v-(v,v)1:v\in V),
\]
其中 $T(V)=\bigoplus_{r\ge0}V^{\otimes r}$. 等价关系为 $uv+vu=2(u,v)1$.
(i) 若型非退化, 证明 $\dim V=2n$ 时此代数为 $\operatorname{Mat}_{2^n}(\C)$, 有唯一的不可约表示, 维数 $2^n$; $\dim V=2n+1$ 时为两个这样的矩阵代数的直和, 有两个不可约表示.
(ii) 证明代数半单当且仅当型非退化. 型退化时求 $\operatorname{Cl}(V)/\operatorname{Rad}(\operatorname{Cl}(V))$.
\end{exercise}
\begin{solution}
先证明一个维数事实. 对任意基 $x_1,\ldots,x_d$, 用 $x_jx_i=-x_ix_j+2(x_i,x_j)$ 交换相邻乱序因子, 用 $x_i^2=(x_i,x_i)$ 消去重复, 可将任意字写成 $x_{i_1}\cdots x_{i_r}$, $i_1<\cdots<i_r$, 的线性组合, 所以维数至多 $2^d$.

这些单项式也线性无关. 在 $\bigwedge V$ 上令 $L_v$ 为左外乘 $v$, $\iota_v$ 为收缩:
\[
 \iota_v(w_1\wedge\cdots\wedge w_r)
 =\sum_{j=1}^r(-1)^{j-1}(v,w_j)
 w_1\wedge\cdots\widehat{w_j}\cdots\wedge w_r.
\]
直接展开得到 $\iota_vL_w+L_w\iota_v=(v,w)I$, 且两种同类算子分别反交换. 所以 $L_v+\iota_v$ 满足Clifford关系. 一个有序单项式作用于 $1\in\bigwedge^0V$ 时, 最高外次数项恰为 $x_{i_1}\wedge\cdots\wedge x_{i_r}$. 若存在单项式的线性关系, 取其中最大次数, 这些不同的外积基向量迫使所有最高次数系数为零; 向下归纳得关系平凡. 因而 $\dim\operatorname{Cl}(V)=2^d$, 对退化型也成立.

(i) 偶数维非退化复对称型可取双曲基 $a_i,b_i$, 使 $(a_i,a_j)=(b_i,b_j)=0$, $(a_i,b_j)=\delta_{ij}/2$. 这是将对角正交基每两个组合成各向同性向量得到的. 在 $S=\bigwedge\langle a_1,\ldots,a_n\rangle$ 上, 让 $a_i$ 左外乘, $b_i$ 按 $b_i(a_j)=\delta_{ij}$ 作外代数收缩. 它们满足 $a_ib_j+b_ja_i=\delta_{ij}$.

算子 $P_0=\prod_i b_ia_i$ 是投影到常数外积的算子: $b_ia_i$ 在未含 $a_i$ 的外积上为1, 在含 $a_i$ 的外积上为0, 这些投影互相交换. 对子集 $I,J$, 适当调整收缩的顺序和符号后, $a_IP_0b_J$ 将基向量 $a_J$ 送至 $a_I$, 将其余基向量送至零. 所以Clifford代数的像包含全部矩阵单位, 等于 $\End(S)$. 两边维数同为 $2^{2n}$, 故同构. 矩阵代数的唯一简单模是列向量空间 $S$, 从而得到不可约性与半单性.

奇数维再加一个与所有 $a_i,b_i$ 正交且 $c^2=1$ 的向量 $c$. 外代数上的次数奇偶算子 $\Gamma=(-1)^{\deg}$ 与每个 $a_i,b_i$ 反交换. 令 $c$ 分别作用为 $\Gamma$ 与 $-\Gamma$, 得到 $S_+,S_-$. 偶数部分已经作用为全部 $\End(S)$; 代数中也有 $\Gamma=\prod_i(1-2a_ib_i)$ 的对应元素. 因而 $c\Gamma$ 在两个表示上分别为 $+I,-I$, 证明两者不同构, 并且 $(1\pm c\Gamma)/2$ 分别投影到两个矩阵块. 两个表示合起来给出满射
\[
 \operatorname{Cl}(V)\longrightarrow\End(S_+)\oplus\End(S_-).
\]
两边维数均为 $2^{2n+1}$, 所以这是同构, 也证明没有遗漏其他不可约表示. $n=0$ 时同一论证给出 $\operatorname{Cl}(\C)\cong\C\oplus\C$.

(ii) 令 $R=\{r:(r,V)=0\}$ 为型的根基, 取补空间 $V=W\oplus R$. $W$ 上的型非退化. 理想 $J=(R)$ 的每个基单项式含至少一个 $R$ 中的基向量. 因为这些向量与所有生成元反交换且平方为零, $J^{\dim R+1}=0$. 商代数为 $\operatorname{Cl}(W)$, 由(i)半单.

幂零理想杀死每个简单模: 若 $JM\ne0$, 简单性迫使 $JM=M$, 与 $J$ 幂零矛盾. 所以 $J\subset\operatorname{Rad}$, 其中根基可定义为所有简单模零化理想的交. 反过来, 半单商 $\operatorname{Cl}(V)/J$ 的简单模的零化理想交为零, 故 $\operatorname{Rad}\subset J$. 因而
\[
 \operatorname{Rad}(\operatorname{Cl}(V))=(R),\qquad
 \operatorname{Cl}(V)/\operatorname{Rad}\cong\operatorname{Cl}(V/R).
\]
若 $R\ne0$, 前面的基证明保证 $R$ 在Clifford代数中的像非零, 所以根基非零, 代数不半单. 若 $R=0$, (i)已经证明半单.
\end{solution}

\begin{exercise}\label{hw:3.4}
\textbf{MIT 18.712 (2010), Homework 5, 原题 3.4; 旧PDF第34页.}
设 $|G|=p^n$, 底域 $k$ 的特征为 $p$. 证明 $G$ 的每个不可约 $k$ 表示均为平凡表示.
\end{exercise}
\begin{solution}
对 $|G|$ 归纳. 非平凡有限 $p$ 群的中心非平凡: 类方程中非中心共轭类的大小均被 $p$ 整除, 所以 $p\mid|Z(G)|$. 在中心中取一个阶为 $p$ 的元素 $z$. 在特征 $p$ 中,
\[
 (\rho(z)-I)^p=\rho(z)^p-I=0.
\]
一个幂零算子在非零空间上有非零核, 而 $z$ 中心保证该核为子表示. 不可约性使核等于整个表示, 所以 $z$ 作用平凡. 表示因此下降为 $G/\langle z\rangle$ 的不可约表示, 归纳得整个 $G$ 作用平凡. 一个平凡不可约表示只能是一维, 因为其任意子空间均不变. 这说明特征整除群阶时, Maschke定理的结论确实失效.
\end{solution}

\originalsection{Homework 6: 有限群的具体计算}

\begin{exercise}\label{hw:3.17}
\textbf{MIT 18.712 (2010), Homework 6, 原题 3.17; 旧PDF第43页.}
设 $G$ 为正 $N$ 边形的全部对称群, $|G|=2N$.
(a) 分别在 $N$ 为奇数和偶数时求所有复不可约表示.
(b) 令 $V$ 为平面标准实表示的复化, 求 $V\otimes V$ 的不可约分解.
\end{exercise}
\begin{solution}
写 $G=\langle r,s:r^N=s^2=1,\ srs=r^{-1}\rangle$, $\zeta=e^{2\pi i/N}$. 包括把正2边形视为线段的边界情形, 因此 $N\ge2$. 对 $k\in\Z/N\Z$, 定义二维表示
\[
 W_k(r)=\begin{pmatrix}\zeta^k&0\\0&\zeta^{-k}\end{pmatrix},
 \qquad W_k(s)=\begin{pmatrix}0&1\\1&0\end{pmatrix}.
\]
若 $2k\ne0\pmod N$, $r$ 的两条特征直线互异, $s$ 交换它们, 所以表示不可约. $W_k\cong W_{N-k}$, 且除这对指标外不同构, 由 $r$ 的特征值集合可知.

一维表示中 $r=r^{-1}$, 所以 $r$ 只能取允许的 $\pm1$, $s$ 任取 $\pm1$. 当 $N$ 为奇数时有 $1$ 与反射符号表示 $\varepsilon$ ($r\mapsto1,s\mapsto-1$), 以及 $W_k$, $1\le k\le(N-1)/2$. 其维数平方和为 $2+4(N-1)/2=2N$. 当 $N$ 为偶数时有四个一维表示 $\chi_{\alpha,\beta}$ ($r\mapsto\alpha,s\mapsto\beta$, $\alpha,\beta=\pm1$), 以及 $W_k$, $1\le k<N/2$; 平方和为 $4+4(N/2-1)=2N$. 因而清单完整.

(b) $V=W_1$. 用其基 $e_+,e_-$, 张成 $e_+\otimes e_+,e_-\otimes e_-$ 的子空间是 $W_2$; 其余两个基向量的和与差分别给 $1$ 与 $\varepsilon$. 所以
\[
 V\otimes V\cong1\oplus\varepsilon\oplus W_2.
\]
这里须把 $W_2$ 继续判断是否不可约. $N=3$ 时 $W_2\cong W_1$; $N=4$ 时 $r$ 在 $W_2$ 上为 $-I$, 所以 $W_2=\chi_{-1,1}\oplus\chi_{-1,-1}$. 对 $N\ge5$, $W_2$ 不可约 (指标可换成 $N-2$). 这避免在正方形情形把可分解的二维项误称不可约.

最后处理 $N=2$. 此时 $G\cong C_2\times C_2$, (a)的偶数情形已列出它的四个一维表示, 没有二维不可约表示. 标准平面表示中, $r=-I$, $s$ 在沿线段与垂直线段的两条直线上分别作用为 $1,-1$, 所以 $V=\chi_{-1,1}\oplus\chi_{-1,-1}$. 两个自张量均平凡, 两个交叉张量均为 $\varepsilon$ ($r\mapsto1,s\mapsto-1$), 故 $V\otimes V=2\cdot1\oplus2\cdot\varepsilon$. 也可在通式中使用 $W_2=W_0=1\oplus\varepsilon$ 得到同一结果.
\end{solution}

\begin{exercise}\label{hw:3.18}
\textbf{MIT 18.712 (2010), Homework 6, 原题 3.18; 旧PDF第43页.}
令 $G$ 为所有 $g(a,b,c)=\begin{pmatrix}1&a&c\\0&1&b\\0&0&1\end{pmatrix}$, $a,b,c\in\mathbb F_p$, 组成的Heisenberg群. 给定 $z^p=1$, 在 $\mathbb F_p$ 上复函数空间中要求
\[
 (\rho(g(1,0,0))f)(x)=f(x-1),\qquad
 (\rho(g(0,1,0))f)(x)=z^xf(x).
\]
(a) 证明该表示存在且唯一, 求所有 $g$ 的作用.
(b) 记此表示为 $R_z$, 证明其不可约当且仅当 $z\ne1$.
(c) 分类所有一维表示, 分解 $R_1$, 并说明其中每个出现的一维表示只出现一次.
(d) 利用维数平方和分类全部不可约表示.
\end{exercise}
\begin{solution}
群乘法为 $g(a,b,c)g(a',b',c')=g(a+a',b+b',c+c'+ab')$. 写 $A=g(1,0,0)$, $B=g(0,1,0)$, $C=g(0,0,1)$, 则 $AB=CBA$. 指定的移位 $T$ 与乘法 $M$ 满足 $TM=z^{-1}MT$, 所以 $C$ 必作用为 $z^{-1}I$. 因为 $g(a,b,c)=A^aB^bC^{c-ab}$, 唯一可能的公式是
\[
 (\rho(g(a,b,c))f)(x)=z^{bx-c}f(x-a).             \tag{H6.1}
\]
将两次作用复合, 指数成为 $(b+b')x-c-c'-ab'$, 正好符合群乘法, 证明存在性. $\mathbb F_p$ 指数可用任意整数代表, 因为 $z^p=1$.

若 $z\ne1$, 则 $z$ 是本原 $p$ 次单位根, $M$ 在点质量基 $\delta_x$ 上有互异特征值 $z^x$. 任意不变子空间由这些特征向量的一部分张成, 而 $T$ 循环置换所有 $\delta_x$, 所以非零子空间必为整个空间. 若 $z=1$, $B,C$ 平凡, $A$ 为循环移位, 故表示分解为其 $p$ 条特征直线而不可约.

一维表示杀死交换子 $C$, 所以因子化为 $G/\langle C\rangle\cong\mathbb F_p^2$. 固定 $\omega=e^{2\pi i/p}$, 全部一维表示为
\[
 L_{u,v}(g(a,b,c))=\omega^{ua+vb},\qquad u,v\in\mathbb F_p.
\]
函数 $f_u(x)=\omega^{-ux}$ 张成 $R_1$ 中的 $L_{u,0}$, 因而 $R_1=\bigoplus_uL_{u,0}$, 每项一次. 对各 $z\ne1$, 中心 $C$ 的标量 $z^{-1}$ 不同, 所以 $R_z$ 两两不同构. 已有 $p^2$ 个一维表示与 $p-1$ 个 $p$ 维表示, 维数平方和为 $p^2+(p-1)p^2=p^3=|G|$, 证明清单完整, 包括 $p=2$.
\end{solution}

\begin{exercise}\label{hw:3.19}
\textbf{MIT 18.712 (2010), Homework 6, 原题 3.19; 旧PDF第43页.}
设 $V$ 为有限维复向量空间. 证明 $\Sym^nV$ 与 $\bigwedge^mV$, $0\le m\le\dim V$, 在自然作用下是 $\GL(V)$ 的不可约表示.
\end{exercise}
\begin{solution}
按题目的不可约性要求, 所讨论的空间应非零; $V=0,n>0$ 时对称幂为零, 不称不可约. 其余情形中, 次数零的空间是一维平凡表示. 设 $d=\dim V>0$, 取基 $e_1,\ldots,e_d$. 对称幂的基为 $e^\alpha=e_1^{\alpha_1}\cdots e_d^{\alpha_d}$, $\sum\alpha_i=n$; 外幂的基为 $e_I=e_{i_1}\wedge\cdots\wedge e_{i_m}$.

取整数 $B>\max(n,1)$ 与实数 $t>1$, 令 $H e_i=t^{B^{i-1}}e_i$. 在对称幂中的特征值为 $t^{\sum\alpha_iB^{i-1}}$, 由 $B$ 进制表达唯一而互异; 外幂相同, 指数的数字为0或1. 因此一个 $H$ 不变子空间必由一部分基向量张成: 各特征投影都是 $H$ 的插值多项式, 可分别提取这些分量.

若对称幂的非零不变子空间包含 $e^\alpha$ 且 $\alpha_j>0$, 用 $1+E_{ij}$, $i\ne j$, 作用, 展开 $(e_j+e_i)^{\alpha_j}$, 再以特征投影提取, 得到 $e^{\alpha-e_j+e_i}$; 其系数为 $\alpha_j\ne0$. 重复这种搬移可从任意次数 $n$ 的单项式到达任意另一个, 所以子空间包含全基.

对外幂, 若 $j\in I,i\notin I$, 则 $(1+E_{ij})e_I=e_I\pm e_{I\setminus\{j\}\cup\{i\}}$. 特征投影取得后者. 任意两个 $m$ 元子集可逐次交换一个元素连接, 故同样包含全基. 这证明两种表示均不可约, 并明确使用了复数域的特征零性质.
\end{solution}

\begin{exercise}\label{hw:3.20}
\textbf{MIT 18.712 (2010), Homework 6, 原题 3.20; 旧PDF第43页.}
有限无重边图的邻接矩阵在相邻顶点处为1, 其余为0. 证明若该图的自同构群非交换, 邻接矩阵必有重特征值.
\end{exercise}
\begin{solution}
设邻接矩阵为 $A$. 图为无向图, 所以 $A$ 是实对称矩阵. 每个图自同构的置换矩阵 $P$ 满足 $PA=AP$. 若 $A$ 的所有特征值均单重, 它有由一维特征子空间组成的正交特征基; 每个 $P$ 保持各特征直线, 因而在该共同基下为对角矩阵. 于是全部 $P$ 两两交换. 自同构在顶点集上的作用忠实, 所以自同构群也交换, 与条件矛盾.
\end{solution}

\begin{exercise}\label{hw:3.21}
\textbf{MIT 18.712 (2010), Homework 6, 原题 3.21; 旧PDF第44页.}
正二十面体的旋转群为 $A_5$.
(a) 分解其在12个顶点上复函数空间的表示.
(b) 分别分解其在20个面及30条边上复函数空间的表示.
\end{exercise}
\begin{solution}
记 $A_5$ 的不可约表示为 $1,U,U',T,F$, 维数分别为 $1,3,3,4,5$. 令 $\phi=(1+\sqrt5)/2$, $\phi'=(1-\sqrt5)/2$. 所需特征标表为
\[
\begin{array}{c|rrrrr}
 &1&2A&3A&5A&5B\\ \hline
\text{类大小}&1&15&20&12&12\\
1&1&1&1&1&1\\
U&3&-1&0&\phi&\phi'\\
U'&3&-1&0&\phi'&\phi\\
T&4&0&1&-1&-1\\
F&5&1&-1&0&0
\end{array}                                                     \tag{H6.2}
\]
可直接得到这些表示与数值: $U$ 是空间旋转表示, 阶2、3、5旋转的迹是 $1+2\cos\theta$; 共轭的 $\sqrt5$ 取值给 $U'$. $T$ 是五点置换表示去掉常数, $F$ 是 $\Sym^2U$ 去掉标量型. 对后三者及 $U,U'$ 用表中的类大小计算, 行范数均为1, 行间内积为0, 且维数平方和 $1+9+9+16+25=60$, 从而这是完整表. $U'$ 的存在也可由 $A_5$ 的外自同构取得: 以奇置换共轭, 交换两个5循环类.

置换特征标等于固定对象数. 顶点、面、边的特征标分别是
\[
 (12,0,0,2,2),\quad (20,0,2,0,0),\quad (30,2,0,0,0).
\]
因为阶5轴穿过一对相对顶点, 阶3轴穿过一对相对面的中心, 阶2轴穿过一对相对边的中点; 其他非恒等旋转不固定相应对象. 将这三行与(H6.2)各行取内积, 得到
\[
\begin{aligned}
 \operatorname{Fun}(\text{顶点})&=1\oplus U\oplus U'\oplus F,\\
 \operatorname{Fun}(\text{面})&=1\oplus U\oplus U'\oplus2T\oplus F,\\
 \operatorname{Fun}(\text{边})&=1\oplus U\oplus U'\oplus2T\oplus3F.
\end{aligned}
\]
例如顶点中 $U$ 的重数为 $(36+24(\phi+\phi'))/60=1$, 边中 $F$ 的重数为 $(150+30)/60=3$. 各行的维数分别为12、20、30, 与原空间一致.
\end{solution}

\begin{exercise}\label{hw:3.22}
\textbf{MIT 18.712 (2010), Homework 6, 原题 3.22; 旧PDF第44页.}
设 $G$ 为有限域 $\mathbb F_q$ 上所有仿射变换 $x\mapsto ax+b$, $a\in\mathbb F_q^\times,b\in\mathbb F_q$. 求所有复不可约表示、特征标和不可约表示的全部张量积.
\end{exercise}
\begin{solution}
群阶为 $q(q-1)$. 将乘法群 $\mathbb F_q^\times$ 的 $q-1$ 个一维特征标 $\lambda$ 拉回, 得一维表示 $L_\lambda(a,b)=\lambda(a)$. 在 $\mathbb F_q$ 的置换表示中去掉常数, 得 $q-1$ 维子表示 $W$ (函数值总和为零). 恒等仿射变换固定 $q$ 点, 非零平移不固定点, 而 $a\ne1$ 的仿射变换有唯一固定点 $b/(1-a)$. 所以
\[
 \chi_W(a,b)=\begin{cases}q-1,&a=1,b=0,\\-1,&a=1,b\ne0,\\0,&a\ne1.\end{cases}
\]
其范数为 $((q-1)^2+(q-1))/(q(q-1))=1$, 因而不可约. 已列表示的平方和为 $(q-1)+(q-1)^2=q(q-1)$, 清单完整. $q=2$ 时 $W$ 本身是一维, 仍是与平凡表示不同的另一不可约表示.

一维表示间显然有 $L_\lambda\otimes L_\mu=L_{\lambda\mu}$. 由于 $\lambda(1)=1$ 且 $\chi_W$ 仅在 $a=1$ 处非零, $L_\lambda\otimes W\cong W$. 最后
\[
 \inner{\chi_W^2}{\lambda}=1,\qquad
 \inner{\chi_W^2}{\chi_W}
 =\frac{(q-1)^3-(q-1)}{q(q-1)}=q-2.
\]
因此 $W\otimes W\cong\bigoplus_\lambda L_\lambda\oplus(q-2)W$. 维数为 $(q-1)+(q-2)(q-1)=(q-1)^2$, 也检查了这一分解.
\end{solution}
\originalsection{Homework 7: 实表示与张量积}

\begin{exercise}\label{hw:3.23}
\textbf{MIT 18.712 (2010), Homework 7, 原题 3.23; 旧PDF第44页.}
令 $G=\operatorname{SU}(2)$, 将标准空间 $V=\C^2$ 看成四维实表示.
(a) 证明 $V$ 实不可约.
(b) 证明交换代数 $H=\End_{\R[G]}V$ 为四维实除代数.
(c) 找基 $1,i,j,k$, 满足 $i^2=j^2=k^2=-1$, $ij=-ji=k$, $jk=-kj=i$, $ki=-ik=j$, 从而四元数群 $Q_8$ 是 $H^\times$ 的子群.
(d) 对 $q=a+bi+cj+dk$ 定义 $\bar q=a-bi-cj-dk$ 与 $\|q\|^2=q\bar q$. 证明 $\overline{q_1q_2}=\bar q_2\bar q_1$ 及范数乘法性.
(e) 证明单位四元数群同构于 $\operatorname{SU}(2)$.
(f) 单位四元数在纯虚四元数空间上以 $x\mapsto qxq^{-1}$ 作用. 证明由此得到的 $h:\operatorname{SU}(2)\to\operatorname{SO}(3)$ 满射, 核为 $\{1,-1\}$.
\end{exercise}
\begin{solution}
(a) $\operatorname{SU}(2)$ 在 $\C^2$ 的单位球面上传递: 任意单位向量 $(z_1,z_2)$ 都是某个矩阵 $\begin{pmatrix}z_1&-\bar z_2\\z_2&\bar z_1\end{pmatrix}$ 的第一列. 非零实不变子空间含一个非零向量, 也就包含同长度的整个球面, 其线性张成是整个实空间. 故实不可约.

(b),(c) 令 $i$ 为乘以复数 $i$, 令 $j(z_1,z_2)=(-\bar z_2,\bar z_1)$, $k=ij$. 直接计算 $j^2=-1$, $ji=-ij$, 而每个 $\operatorname{SU}(2)$ 矩阵都与 $j$ 交换. 它们产生四元数乘法关系.

任意实线性算子 $T$ 唯一分解为复线性部分 $(T-iTi)/2$ 与复反线性部分 $(T+iTi)/2$. 若 $T$ 与群作用交换, 两部分也交换. 标准复表示不可约 (由(a)即知), Schur引理使复线性部分为复标量. 复反线性部分与 $j^{-1}$ 复合为复线性群同态, 所以为复标量乘 $j$. 因而 $H=\C\oplus\C j$ 为四维实代数, 基为 $1,i,j,k$. 若 $T\ne0$, 实Schur引理给出其核为零且像为全空间, 所以可逆, 逆仍与群作用交换. 八个元素 $\pm1,\pm i,\pm j,\pm k$ 构成 $Q_8$.

(d) 基元素的关系逐项给出共轭反转乘法次序, 双线性推广得 $\overline{q_1q_2}=\bar q_2\bar q_1$. 同样展开有 $q\bar q=\bar qq=a^2+b^2+c^2+d^2$. 因此
\[
 (q_1q_2)\overline{q_1q_2}
 =q_1q_2\bar q_2\bar q_1=\|q_2\|^2\|q_1\|^2.
\]
取非负平方根得到范数乘法性, 并得非零四元数的逆为 $\bar q/\|q\|^2$.

(e) 此处的抽象四元数代数另有复矩阵实现
\[
 a+bi+cj+dk\longmapsto
 \begin{pmatrix}a+bi&c+di\\-c+di&a-bi\end{pmatrix}.
\]
三个虚单位的像满足相同乘法关系, 所以此映射保持乘法且单射. 矩阵与其共轭转置之积为 $\|q\|^2I$, 行列式为 $\|q\|^2$. 范数1的像恰为所有 $\operatorname{SU}(2)$ 矩阵, 因为后者均有 $\begin{pmatrix}\alpha&\beta\\-\bar\beta&\bar\alpha\end{pmatrix}$, $|\alpha|^2+|\beta|^2=1$, 的形式. 这证明同构, 也将群的底层空间识别为三维球面 $S^3$.

(f) 将纯虚四元数看作 $\R^3$, 有 $uv=-u\cdot v+u\times v$. 共轭作用保持纯虚性和范数. 若 $u$ 为单位纯虚四元数, 则 $q=\cos(\theta/2)+u\sin(\theta/2)$ 的共轭作用为
\[
 x\longmapsto x\cos\theta+(u\times x)\sin\theta
              +u(u\cdot x)(1-\cos\theta).
\]
这由展开 $qx\bar q$ 得到, 是绕 $u$ 轴旋转 $\theta$ 的公式. 每个三维旋转都有一个固定轴, 并在正交平面上是二维旋转, 所以上式给出满射. 若 $q$ 在核中, 它与 $i,j,k$ 都交换; 由乘法关系比较系数可得 $q$ 为实数. 再用范数1得到 $q=\pm1$.
\end{solution}

\begin{exercise}\label{hw:3.24}
\textbf{MIT 18.712 (2010), Homework 7, 原题 3.24; 旧PDF第44--45页.}
(a) 推导 $\operatorname{SO}(3)$ 的有限子群分类: 循环群、阶 $2n$ 的二面体旋转群 ($n\ge2$)、正四面体旋转群 $A_4$、立方体/正八面体旋转群 $S_4$、正十二面体/正二十面体旋转群 $A_5$.
(b) 利用 $\operatorname{SU}(2)\to\operatorname{SO}(3)$ 分类 $\operatorname{SU}(2)$ 的有限子群.
\end{exercise}
\begin{solution}
(a) 设非平凡有限子群 $G$ 的阶为 $n$. 非恒等旋转在单位球面上恰固定轴的两个端点, 称为极点. 某极点 $P$ 的稳定子是绕该轴的有限旋转群, 因而循环, 设阶为 $m_P\ge2$. 若极点轨道代表为 $P_1,\ldots,P_k$, 对有序对 $(g,P)$, $g\ne1,gP=P$, 计数, 得
\[
 2(n-1)=\sum_{i=1}^k\frac n{m_i}(m_i-1),\qquad
 2\left(1-\frac1n\right)=\sum_{i=1}^k\left(1-\frac1{m_i}\right). \tag{H7.1}
\]
每项至少 $1/2$, 而总和小于2, 故 $k\le3$. $k=1$ 不可能: 左边至少1, 右边小于1. $k=2$ 时, $1/m_1+1/m_2=2/n$; 因为 $m_i\le n$, 每项至少 $1/n$, 故 $m_1=m_2=n$. 全群固定同一轴, 所以为循环群.

$k=3$ 时按 $m_1\le m_2\le m_3$ 排序, 有 $1/m_1+1/m_2+1/m_3=1+2/n>1$. 若 $m_1\ge3$, 总和至多1, 所以 $m_1=2$. 若 $m_2\ge4$, 总和仍至多1. 剩下的可能恰为
\[
\begin{array}{c|c}
(m_1,m_2,m_3)&n\\\hline
(2,2,m)&2m\\
(2,3,3)&12\\
(2,3,4)&24\\
(2,3,5)&60
\end{array}                                                     \tag{H7.2}
\]
还须由这张数值表识别实际旋转群, 不能仅靠群阶判定.

对 $(2,2,m)$ 且 $m>2$, 阶 $m$ 的极点轨道有两个点. 它们必互为相反点, 因为稳定子的非恒等旋转仅固定这两个点. 群保持该轴, 绕轴旋转组成阶 $m$ 循环子群, 其余元素交换轴的两端, 必为垂直轴的半周旋转. 选一条这样的轴, 得生成元 $r,s$ 满足 $r^m=s^2=1,srs=r^{-1}$, 即二面体旋转群. $m=2$ 时群阶4, 三个非恒等元素均为半周旋转. 任意两个不同的轴必须垂直, 因为两次半周旋转的复合仍为半周旋转; 可从旋转复合的角度公式或四元数乘法验证. 因而为三条互相垂直轴的半周旋转群 $D_2$.

后三种情形中, 空间表示实不可约. 若有不变直线, 群保持一对相反极点; 保持轴方向的子群为循环, 交换两端的元素为半周旋转, 故群只能为上述循环或二面体群. 若有不变平面, 它的正交补给出不变直线, 也不可能. 因而对任意轨道 $\mathcal O$, 对称算子 $\sum_{P\in\mathcal O}PP^{\mathsf T}$ 的各特征空间都不变, 它必为 $|\mathcal O|I/3$. 同时 $\sum P=0$, 因为此和是固定向量.

对 $(2,3,3)$, 取阶3极点的四点轨道. 一个极点的阶3稳定子循环置换另外三点 (若其相反点也在轨道, 剩下两点不可能被非恒等三循环作用). 所以该点与另外三点的内积相同. 用四点和为零得内积 $-1/3$, 即正四面体顶点. 群嵌入四面体旋转群且两者阶均为12, 所以相等, 对四个顶点的作用给出 $A_4$.

对 $(2,3,4)$, 取阶4极点的六点轨道. 稳定子的非轴点轨道长为4, 所以除所选极点外的五点中必含其相反点, 其余四点是绕轴的正方形. 整个轨道有三对相反点, 故该正方形的纬度为零. 这是正八面体的六个顶点. 群等于其阶24旋转群; 对立方体四条体对角线的忠实作用将它识别为 $S_4$.

对 $(2,3,5)$, 取阶5极点的十二点轨道. 同理它含所选点 $P$ 及 $-P$, 其余十点组成两个正五边形, 位于相反纬度 $a,-a$. 用前述对称算子在 $P$ 方向的值, 得 $2+10a^2=12/3=4$, 所以 $a^2=1/5$. 两个五边形由取相反点对应, 方位角相差 $\pi$, 即模 $2\pi/5$ 相差 $\pi/5$. 这些是正二十面体的标准顶点: 每个点恰与内积 $1/\sqrt5$ 的五点相连, 共30条等长边、20个等边三角形面. 因而群等于其阶60旋转群. 为识别抽象群, 先计算共轭类. 阶3与阶5旋转的中心化子只能是同轴旋转, 阶分别为3与5, 因而类大小为20与12. 阶2轴的无序极点对稳定子阶为4, 正是该半周旋转的中心化子, 故15个阶2元素组成一个共轭类. 每个此中心化子都是四元Klein群: 其中不存在阶4元素, 因为所有元素的阶由极点稳定子只可能为1、2、3、5. 每个阶2元素只属于一个这样的四元群, 因为四元群交换而其阶已等于中心化子阶. 所以15个阶2元素分成五个四元群的三非恒等元素集. 共轭作用传递地置换这五个群, 给出 $G\to S_5$. 其像的阶被5整除, 故核阶至多12. 但任何非平凡正规子群都包含一个完整非恒等共轭类, 这些类大小分别为15、20、12、12, 连同单位元使子群阶至少13. 因而核平凡, 像为 $S_5$ 的指数2子群 $A_5$. 指数2子群必含全部偶置换: 所有换位在商群中有同一像, 所以偶数个换位的乘积落入核, 而 $A_5$ 本身已有阶60. 这完成几何群的识别. 各类列出的标准旋转群均存在, 所以分类无缺项.

(b) 设 $H\subset\operatorname{SU}(2)$ 有限, $G=h(H)$. 若 $-1\in H$, 则 $H=h^{-1}(G)$: 每个 $g\in G$ 的两个提升都在 $H$, 而总共只有这两个. 因而得到循环群的二重提升 (阶 $2n$ 的循环群), 二面体群的二重提升 (阶 $4n$ 的二元二面体群), 以及二元四面体、二元八面体、二元二十面体群, 阶分别为24、48、120. 二元二面体群可写为
\[
 \langle x,y:x^{2n}=1,\ y^2=x^n,\ yxy^{-1}=x^{-1}\rangle,
 \qquad n\ge2.
\]
若 $-1\notin H$, 映射 $h|_H$ 单射. $\operatorname{SU}(2)$ 中唯一的阶2元素为 $-1$, 因为单位四元数 $q^2=1$ 迫使 $q=\pm1$. Cauchy定理使 $|H|$ 为奇数, 所以 $G$ 的上述分类只能给奇数阶循环群. 它的全提升是阶 $2|G|$ 的循环群, 其中唯一的奇数阶子群就是 $H$. 合并两种情形, 完整名单是任意阶的循环群, 阶 $4n$ ($n\ge2$) 的二元二面体群, 及三种二元多面体群. 单位子群也包括在循环群中.
\end{solution}

\begin{exercise}\label{hw:3.25}
\textbf{MIT 18.712 (2010), Homework 7, 原题 3.25; 旧PDF第45页.}
求原题3.18中的Heisenberg群所有复不可约表示的特征标及张量积.
\end{exercise}
\begin{solution}
沿用习题\ref{hw:3.18}的 $L_{u,v}$ 与 $R_z$, $z\ne1,z^p=1$. 一维特征标是 $\omega^{ua+vb}$. 在公式(H6.1)的点质量基中, $a\ne0$ 时作用移位, 迹为零; $a=0,b\ne0$ 时迹为 $z^{-c}\sum_xz^{bx}=0$. 所以
\[
 \chi_{R_z}(g(a,b,c))=
 \begin{cases}p z^{-c},&a=b=0,\\0,&(a,b)\ne(0,0).\end{cases}
\]
中心指数的负号来自原题移位 $f(x-1)$ 的约定, 必须与3.18一致.

一维表示的张量积为 $L_{u,v}\otimes L_{u',v'}=L_{u+u',v+v'}$. 一维表示在中心上恒为1, 所以 $L_{u,v}\otimes R_z\cong R_z$. 若 $zw\ne1$, $R_z\otimes R_w$ 的特征标在中心为 $p^2(zw)^{-c}$, 在其他元素处为0, 即 $p\chi_{R_{zw}}$, 因而
\[
 R_z\otimes R_w\cong pR_{zw}\quad(zw\ne1).
\]
若 $zw=1$, 对所有一维表示之和, 在 $(a,b)=(0,0)$ 时特征标为 $p^2$, 否则为0 (分别对 $u,v$ 求单位根和). 故
\[
 R_z\otimes R_{z^{-1}}\cong
 \bigoplus_{u,v\in\mathbb F_p}L_{u,v}.
\]
两种结果的维数均为 $p^2$; 对 $p=2$ 只有第二种非线性张量积情形.
\end{solution}

\begin{exercise}\label{hw:3.26}
\textbf{MIT 18.712 (2010), Homework 7, 原题 3.26; 旧PDF第45页.}
若有限群 $G$ 的复表示 $V$ 忠实, 证明每个不可约表示都包含在某个 $\Sym^nV$ 中, 从而包含在某个 $V^{\otimes n}$ 中.
\end{exercise}
\begin{solution}
对每个 $g\ne1$, 忠实性使 $V^*$ 上 $g$ 的固定空间为真子空间. 复向量空间不是有限多个真线性子空间的并: 可分别选一个非零线性方程在各子空间上消失, 它们的乘积是非零多项式, 不可能处处为零. 因而可取 $u\in V^*$ 的稳定子为1, 其轨道含 $|G|$ 个不同点.

将 $\Sym V$ 看作 $V^*$ 上的多项式环, 定义 $\Phi(f)(g)=f(gu)$. 群作用满足 $(hf)(\xi)=f(h^{-1}\xi)$, 所以 $\Phi(hf)(g)=\Phi(f)(h^{-1}g)$, 即 $\Phi$ 是到左正则函数表示的同态. 它满射: 对轨道中的一个点 $u_i$, 每个其他点 $u_j$ 可取仿射线性函数 $\ell_{ij}$, 使 $\ell_{ij}(u_i)=1,\ell_{ij}(u_j)=0$. 乘积 $\prod_{j\ne i}\ell_{ij}$ 在轨道上给出点质量函数, 这些函数构成目标空间的一组基.

这些有限多个插值多项式有某个最大次数 $D$, 因而 $\bigoplus_{n=0}^D\Sym^nV$ 已满射到正则表示. Maschke定理使该满射分裂, 正则表示又包含每个不可约表示. 对一个不可约表示 $L$ 的嵌入, 投影到某个次数分量必须非零, 所以把 $L$ 嵌入对应的 $\Sym^nV$. 最后复数域中对称化算子 $\frac1{n!}\sum_{\sigma\in S_n}\sigma$ 将对称幂识别为 $V^{\otimes n}$ 的子空间. 平凡表示可取 $n=0$.
\end{solution}

\begin{exercise}\label{hw:3.27}
\textbf{MIT 18.712 (2010), Homework 7, 原题 3.27; 旧PDF第45--46页.}
弹性论中令 $V=\R^3$. 小变形的变形张量为对称算子 $d\in\Sym^2V$, 应力张量为 $S\in\End(V)$. 具体地, 对局部变形的微分作极分解 $A=DO$, $D=1+d$; $O$ 为正交旋转部分. 应力定义为通过面积为 $\|v\|$、法向为 $v$ 的小圆盘的力 $F_v=Sv$. 各向同性的线性弹性律是 $\operatorname{SO}(3)$ 等变线性映射 $f:\Sym^2V\to\End(V)$, $S=f(d)$.
(a) 证明 $\End(V)=\R\oplus V\oplus W$, $\Sym^2V=\R\oplus W$, 其中 $W$ 为五维表示.
(b) 证明 $V,W$ 复化后仍不可约. 由此证明 $S$ 总是对称, 且对 $x\in\R,y\in W$ 有 $f(x+y)=Kx+\mu y$, $K,\mu\in\R$.
\end{exercise}
\begin{solution}
(a) 使用欧氏内积识别 $\End(V)$ 与 $V\otimes V$. 一个矩阵唯一分解为标量矩阵、反对称矩阵及对称无迹矩阵. 这些空间在旋转共轭下不变; 反对称矩阵由 $v\mapsto(x\mapsto v\times x)$ 与标准三维表示等变同构. 对称无迹部分 $W$ 维数 $6-1=5$. 对称矩阵恰为标量部分与 $W$ 的和.

(b) 为避免仅用几何直觉断言复不可约, 在多项式模型中检查. 写 $u=x+iy,v=x-iy,w=z$, 则欧氏二次型为 $Q=uv+w^2$. 复旋转李代数的下列微分算子均保持 $Q$:
\[
 H=2(u\partial_u-v\partial_v),\quad
 E=u\partial_w-2w\partial_v,\quad
 F=2w\partial_u-v\partial_w.
\]
直接算出 $[H,E]=2E,[H,F]=-2F,[E,F]=H$. 这三个算子张成保持 $Q$ 的三维李代数, 因而群不变子空间对它们均不变 (沿实旋转的一参数子群微分, 再复线性组合). 在一次多项式空间中, $u,w,v$ 的 $H$ 权分别为 $2,0,-2$, $E,F$ 在相邻权之间非零作用, 所以任意非零不变子空间通过特征投影及升降作用包含全部三条权线.

在二次多项式中, 无迹对称张量对应调和二次式, 即 $4\partial_u\partial_v+\partial_w^2$ 的核. 从 $u^2$ 连续用 $F$ 得
\[
 u^2,\quad4uw,\quad8w^2-4uv,\quad-24wv,\quad24v^2.
\]
这五项均调和, 其 $H$ 权为 $4,2,0,-2,-4$, 因而构成 $W_\C$ 的基. $F$ 连接相邻基向量, 而 $E$ 反向连接且系数非零, 例如由 $[E,F^r]u^2=r(5-r)F^{r-1}u^2$ 可直接验证. 同样的特征投影与升降论证证明 $W_\C$ 不可约.

这三个不可约分量的维数 $1,3,5$ 不同, 因而两两不同构. 对 $f$ 的复化用Schur引理, 从 $\R\oplus W$ 到 $V$ 的分量为零, 在标量空间与 $W$ 上分别为标量 $K,\mu$. $f$ 保持实空间, 所以两标量均为实数, 其像落在对称矩阵空间. 若 $d=(\tr d/3)I+d_0$, 则具体公式是
\[
 S=K\frac{\tr d}{3}I+\mu d_0.
\]
这里 $K,\mu$ 的数值约定跟随原题对两个分量的参数化; 弹性稳定性所要求的正性是额外的物理假设, 不是仅由等变性推出.
\end{solution}

\begin{exercise}\label{hw:4.1}
\textbf{MIT 18.712 (2010), Homework 7, 原题 4.1; 旧PDF第47页.}
设 $V$ 为有限群 $G$ 的复不可约表示. 非自对偶者称复型; 有非退化不变对称双线性型者称实型; 有非退化不变反对称双线性型者称四元数型.
(a) 将 $V$ 看成实表示, 证明其交换代数分别为 $\C$, $\operatorname{Mat}_2(\R)$, $\mathbb H$.
(b) 证明 $V$ 为实型当且仅当存在实表示 $V_\R$ 使 $V\cong V_\R\otimes_\R\C$.
\end{exercise}
\begin{solution}
取平均得到 $G$ 不变的正定Hermitian型 $\langle\ ,\ \rangle$, 约定第一变量复线性. 将任意实线性群自同态分解为复线性与复反线性部分, 前一部分由Schur引理为 $\C I$. 一个反线性群同态等价于 $\bar V\to V$ 的复线性同态. 酉型又给出 $\bar V\cong V^*$, 所以复型时后一部分为零, 交换代数恰为 $\C$.

自对偶时, 一个非零不变双线性型 $B$ 对应唯一反线性可逆算子 $j$, 满足 $B(v,w)=\langle v,jw\rangle$. 不变性使 $j$ 与群作用交换. 因而 $j^2=\lambda I$, $\lambda\ne0$. 比较 $j j^2$ 与 $j^2j$, 利用反线性得 $\lambda=\bar\lambda$. 由Hermitian型的唯一性, 存在 $c>0$ 使 $\langle jv,jw\rangle=c\langle w,v\rangle$. 若 $B(w,v)=\epsilon B(v,w)$, $\epsilon=\pm1$, 则
\[
 c\langle v,v\rangle=B(jv,v)
 =\epsilon B(v,jv)=\epsilon\lambda\langle v,v\rangle.
\]
故 $\lambda$ 的符号为 $\epsilon$. 将 $j$ 乘一个正实数, 可正规化为实型 $j^2=1$, 四元数型 $j^2=-1$. 全部反线性同态是 $\C j$, 因为与 $j^{-1}$ 复合后是复标量. 于是交换代数为 $\C\oplus\C j$, 且 $ji=-ij$.

若 $j^2=-1$, 这正是四元数代数. 若 $j^2=1$, 将 $i,j$ 分别送至 $\begin{pmatrix}0&-1\\1&0\end{pmatrix}$ 与 $\begin{pmatrix}1&0\\0&-1\end{pmatrix}$, 它们满足相同关系且连同乘积与单位构成四维矩阵基, 所以交换代数同构于 $\operatorname{Mat}_2(\R)$. 注意此处研究的是复不可约表示的实化, 实型时它分解为两个同构的实表示, 因而交换代数不是仅有 $\R$.

(b) 实型时取 $V_\R=\{v:jv=v\}$. 对任意 $v$, 唯一分解为
\[
 v=\frac{v+jv}{2}+i\frac{v-jv}{2i},
\]
两项的实向量部分均在 $V_\R$. 所以 $V=V_\R\oplus iV_\R$, 且群保持 $V_\R$, 得到所需复化同构. 反之, 若 $V$ 是实表示的复化, 对实表示平均一个正定实内积, 再作复双线性延拓, 得非退化不变对称型. 因而 $V$ 为实型.
\end{solution}
\originalsection{Homework 8: 诱导与对称群}

\begin{exercise}\label{hw:4.15}
\textbf{MIT 18.712 (2010), Homework 8, 原题 4.15; 旧PDF第51页.}
(a) 证明对任意有限群 $G$, 存在 $\mathbb Q$ 的有限Galois扩张 $K\subset\C$, 使 $G$ 的每个有限维复表示均可选基令所有群元素矩阵的系数落在 $K$ 中.
(b) 若 $V$ 为维数大于1的复不可约表示, 证明存在 $g\in G$ 使 $\chi_V(g)=0$.
\end{exercise}
\begin{solution}
(a) 记 $\overline{\mathbb Q}$ 为复数中的代数数域. Maschke定理在此特征零域成立, 且该域代数闭, 所以群代数分块为
\[
 \overline{\mathbb Q}[G]\cong\bigoplus_i\operatorname{Mat}_{d_i}(\overline{\mathbb Q}).
\]
给出这一分块的依据. 令 $A=\overline{\mathbb Q}[G]$, 每个简单模 $S_i$ 都是 $A$ 的商, 因为任意非零向量生成 $S_i$; 正则模完全可约, 所以所有简单模都在其中出现. 同构 $\Hom_A(A,S_i)\cong S_i$, $f\mapsto f(1)$, 结合Schur引理说明 $S_i$ 在正则模中的重数为 $d_i=\dim S_i$. 因而 $\dim A=\sum_i d_i^2$. 作用映射 $A\to\bigoplus_i\End(S_i)$ 是单射: 若一个元素在所有简单模上为零, 它在正则模上也为零, 作用于单位元即知该元素为零. 两边维数相同, 故映射满射, 得到上面的分块, 无需预先调用一般半单代数分类.

将底域扩张到 $\C$, 上式变为 $\C[G]\cong\bigoplus_i\operatorname{Mat}_{d_i}(\C)$, 因而每个复不可约表示由一个 $\overline{\mathbb Q}$ 上的简单模扩域得到, 没有新增或合并简单模. 对有限多个简单模各选基, 所有有限多个群元素的全部矩阵系数是有限多个代数数. 它们生成一个有限数域, 取其正规闭包 $K$ 即为有限Galois扩张. 任意有限维表示是这些不可约表示的直和, 选分块基后仍只用 $K$ 中的矩阵系数.

(b) 给出旧讲义备注中的圆分域论证, 它还省去(a)的调用. 令 $N=|G|$, $d=\dim V>1$, 反设所有特征标值非零. 因为每个 $g$ 的特征值是 $N$ 次单位根, 所有 $\chi(g)$ 都是 $\mathbb Q(\zeta)$ 中的代数整数, $\zeta=e^{2\pi i/N}$. 酉化给 $\chi(g^{-1})=\overline{\chi(g)}$. 所以
\[
 \alpha=\prod_{g\ne1}\chi(g)\chi(g^{-1})
       =\prod_{g\ne1}|\chi(g)|^2
\]
是正的代数整数. 特征标正交给 $\sum_g|\chi(g)|^2=N$, 故非恒等项的算术平均为 $(N-d^2)/(N-1)<1$. 算术--几何平均不等式给 $0<\alpha<1$.

圆分域的每个Galois自同构形如 $\zeta\mapsto\zeta^j$, $(j,N)=1$, 并把 $\chi(g)$ 变为 $\chi(g^j)$. 幂映射 $g\mapsto g^j$ 是集合上的双射: 取 $\ell$ 满足 $j\ell\equiv1\pmod N$, 则 $g\mapsto g^\ell$ 为其逆; 此处不要求该映射保持群乘法. 它也双射非恒等元素集. 因而每个自同构固定 $\alpha$, 所以 $\alpha\in\mathbb Q$. 有理代数整数为整数 (将既约分数代入首一整系数方程, 分母必须为1), 与 $0<\alpha<1$ 矛盾. 因此存在零特征标值. 本题使用有限代数扩张的正规闭包及圆分域Galois群这两个基础数论事实, 未把它们扩展为一章数论课程.
\end{solution}

\begin{exercise}\label{hw:4.31}
\textbf{MIT 18.712 (2010), Homework 8, 原题 4.31; 旧PDF第55页.}
设 $K\subset H\subset G$ 为群, $V$ 为 $K$ 的表示. 证明
\[
 \Ind_H^G(\Ind_K^HV)\cong\Ind_K^GV.
\]
\end{exercise}
\begin{solution}
在群代数张量模型中, 左边为 $\C[G]\otimes_{\C[H]}(\C[H]\otimes_{\C[K]}V)$. 定义
\[
 F(g\otimes(h\otimes v))=gh\otimes v,\qquad
 J(g\otimes v)=g\otimes(1\otimes v).
\]
对 $h_0\in H$, 两个平衡表达式 $gh_0\otimes(h\otimes v)$ 与 $g\otimes(h_0h\otimes v)$ 的像相同; 对 $k\in K$, $h k\otimes v=h\otimes kv$ 的像也相同. 所以 $F$ 良定义. 对 $J$, $gk\otimes v$ 的像为 $g\otimes(k\otimes v)=g\otimes(1\otimes kv)$, 也良定义. 两映射显然 $G$ 等变, 并且 $FJ=1$, $JF(g\otimes(h\otimes v))=gh\otimes(1\otimes v)=g\otimes(h\otimes v)$. 因而互逆. 论证对一般底域及一般群同样成立.
\end{solution}

\begin{exercise}\label{hw:4.34}
\textbf{MIT 18.712 (2010), Homework 8, 原题 4.34; 旧PDF第57--58页.}
分别对 $A_5$ 的子群 (a) $C_2$, (b) $C_3$, (c) $C_5$, (d) $A_4$, (e) $C_2\times C_2$, 从其每一个不可约表示诱导到 $A_5$, 求全部不可约分解.
\end{exercise}
\begin{solution}
沿用(H6.2)的 $1,U,U',T,F$. Frobenius互反使诱导的重数等于对应子群上的限制重数. 下面各行列出诱导表示中 $(1,U,U',T,F)$ 的重数; 一行完全确定其直和分解.
\[
\begin{array}{c|c|rrrrr}
\text{子群}&\text{其不可约表示}&1&U&U'&T&F\\\hline
C_2&1&1&1&1&2&3\\
&-&0&2&2&2&2\\\hline
C_3&1&1&1&1&2&1\\
&\zeta,\zeta^2&0&1&1&1&2\\\hline
C_5&1&1&1&1&0&1\\
&\xi,\xi^{-1}&0&1&0&1&1\\
&\xi^2,\xi^{-2}&0&0&1&1&1\\\hline
A_4&1&1&0&0&1&0\\
&\eta,\bar\eta&0&0&0&0&1\\
&D&0&1&1&1&1\\\hline
C_2\times C_2&1&1&0&0&1&2\\
&\text{任一非平凡一维表示}&0&1&1&1&1
\end{array}                                                     \tag{H8.1}
\]
一格中并列的子群表示分别有相同分解, 并非先把它们相加. 对 $C_5$, 选择生成元在类 $5A$, 令 $\xi=e^{2\pi i/5}$; 改选 $5B$ 生成元会交换 $U,U'$ 对应的两行. 对 $A_4$, $D$ 是三维不可约表示, $\eta,\bar\eta$ 为通过 $A_4/(C_2\times C_2)\cong C_3$ 的两个非平凡一维表示.

逐项核算如下. $C_2$ 的正、负重数为 $(d\pm\chi(2A))/2$, 给前两行. $C_3$ 的平凡重数为 $(d+2\chi(3A))/3$, 另外两者的重数为 $(d-\chi(3A))/3$. $C_5$ 上, $U$ 的特征值为 $1,\xi,\xi^{-1}$, $U'$ 为 $1,\xi^2,\xi^{-2}$; $T$ 是四个非平凡一维特征标的和, $F$ 是全部五个一维特征标的和, 得到三行. 这些特征值也可将(H6.2)在四个非恒等幂上的值作离散Fourier变换得到.

在 $A_4$ 上, $U,U'$ 限制均为 $D$, $T$ 限制为 $1\oplus D$, $F$ 限制为 $\eta\oplus\bar\eta\oplus D$. 检查恒等、三个双换位、两个各含四个元素的3循环类的值分别相等即可确认, 给出相应三行. 最后 $C_2\times C_2$ 的三个非恒等元素都在 $2A$, 所以平凡重数为 $(d+3\chi(2A))/4$, 每个非平凡重数为 $(d-\chi(2A))/4$. 各行维数依次为指数乘子群表示维数: $30,20,12,5$ (或 $15$), $15$, 与诱导的维数公式相符.
\end{solution}

\textbf{对称群题目的前置工具.}
下面五道关于 $S_n$ 的指定作业题需要一般Young模. 这里补出所用定义和关键工具, 仅服务这些题解; 不把整套Young表、钩长公式与Schur--Weyl表示分类称为正文已全面覆盖的内容.

一个分拆 $\lambda=(\lambda_1\ge\cdots\ge\lambda_r>0)$ 满足 $\sum\lambda_i=n$. Young图第 $j$ 行有 $\lambda_j$ 个格子, 格子位置用 (列 $i$, 行 $j$) 记录, 与原题 $i-j$ 的内容约定一致. 固定给格子填入 $1,\ldots,n$ 的表, 令 $P_\lambda$ 为行置换群, $Q_\lambda$ 为列置换群, 并令
\[
 a_\lambda=\frac1{|P_\lambda|}\sum_{p\in P_\lambda}p,
 \quad b_\lambda=\frac1{|Q_\lambda|}\sum_{q\in Q_\lambda}\sgn(q)q,
 \quad c_\lambda=a_\lambda b_\lambda,
 \quad V_\lambda=\C[S_n]c_\lambda.
\]
左右乘法次序遵循旧讲义, 不任意交换 $a,b$.

\begin{lemma}\label{hw:young-tools}
各 $V_\lambda$ 是两两不同构的全部不可约表示, 且 $c_\lambda^2=\kappa_\lambda c_\lambda$, $\kappa_\lambda>0$. 它们可在 $\mathbb Q$ 上实现.
\end{lemma}
\begin{proof}
先证 $a_\lambda x b_\lambda\in\C c_\lambda$ 对所有 $x\in\C[S_n]$ 成立. 只需 $x=g$. 若有两个数字同处原表的一行, 又同处 $g$ 变换后的表的一列, 相应换位 $t$ 在 $P_\lambda$ 中, $g^{-1}tg\in Q_\lambda$. 因此 $a_\lambda g b_\lambda=a_\lambda t g b_\lambda=-a_\lambda g b_\lambda$, 得零. 若不存在这样的两个数字, 原表第一行的 $\lambda_1$ 个数字在变换后表中分占全部 $\lambda_1$ 列. 可分别在各列内置换, 将它们送到第一行, 再在第一行内排列到原来的位置. 删除已经对齐的第一行, 重复此操作. 最终得到 $g\in P_\lambda Q_\lambda$, 从而 $a_\lambda g b_\lambda=\pm c_\lambda$.

取 $x=b_\lambda a_\lambda$, 得 $c_\lambda^2=\kappa_\lambda c_\lambda$. 因为 $P_\lambda\cap Q_\lambda=\{1\}$, $c_\lambda$ 的单位元系数为 $1/(|P_\lambda||Q_\lambda|)>0$. 群代数左正则作用的迹为群阶乘单位元系数, 因而 $c_\lambda$ 不是平方为零的算子, $\kappa_\lambda\ne0$. 在群元素正交的内积下 $a,b$ 均为正交投影, $aba$ 是非零正半定算子且 $(aba)^2=\kappa_\lambda aba$. 取它的非零特征值, 即得 $\kappa_\lambda>0$.

所以 $e_\lambda=c_\lambda/\kappa_\lambda$ 为幂等元, 上述公式还给 $e_\lambda\C[S_n]e_\lambda=\C e_\lambda$. 在半单代数的矩阵分块中, 一个幂等元在各块的秩为 $r_i$, 其角代数维数为 $\sum r_i^2$. 此处为1, 故恰一个块有秩1, 左理想 $\C[S_n]e_\lambda$ 不可约.

不同分拆按字典序比较. 若 $\lambda>\mu$, 则 $a_\lambda x b_\mu=0$: 在两图前面的等长行可像上述过程逐行对齐并删除, 在第一个长度不同的行, $\lambda$ 的该行比 $\mu$ 剩余图的列数多, 必有两个数字同处一行又同处一列, 给出同样的换位消去. 对此等式取群代数的共轭转置, 得 $b_\mu\C[S_n]a_\lambda=0$, 所以 $e_\mu\C[S_n]e_\lambda=0$. 两个对应的简单左理想因此不可能同构. 分拆数等于 $S_n$ 的共轭类数 (由循环长度分类), 故这些不可约表示已全部列出.

最后 $c_\lambda$ 的系数均有理, 在有理群代数中取其左理想的一组有理基. 左乘群元素的矩阵系数有理, 复化后就是 $V_\lambda$, 证明有理实现.
\end{proof}

为证明分支规则, 还需要一个可计算的特征标模型. 取变量数 $N\ge n$, 令 $p_r(x)=\sum_i x_i^r$. 对循环型 $\mu=(1^{m_1}2^{m_2}\cdots)$, 令 $p_\mu=\prod_rp_r^{m_r}$, $z_\mu=\prod_r r^{m_r}m_r!$. 类大小为 $n!/z_\mu$. 定义特征标映射
\[
 \operatorname{ch}(f)=\sum_{\mu\vdash n}\frac{f(\mu)}{z_\mu}p_\mu,
 \qquad \langle p_\mu,p_\nu\rangle=\delta_{\mu\nu}z_\mu.
\]
这是等距映射, 由类大小与特征标内积公式直接得到. 令 $h_r$ 为次数 $r$ 的完全对称多项式, 即所有非降指标单项式之和, $h_0=1,h_r=0$ ($r<0$). 生成函数
\[
 \sum_{r\ge0}h_r t^r=\prod_i(1-x_it)^{-1}
  =\exp\left(\sum_{r\ge1}\frac{p_r t^r}{r}\right)
\]
说明 $\operatorname{ch}(1_{S_r})=h_r$. 从 $S_a\times S_b$ 的外张量积诱导到 $S_{a+b}$ 时, $\operatorname{ch}$ 将诱导变成乘法: 在给定循环型中选择属于前 $a$ 个点的循环, 对每种长度选出 $k_r$ 个, 产生二项系数 $\binom{m_r}{k_r}$; 与 $z_\mu$ 的阶乘因子相除恰为两个相应 $z$ 因子的倒数. 这逐循环验证了该乘法公式.

\begin{lemma}\label{hw:schur-tools}
令 $\delta=(N-1,N-2,\ldots,0)$, $A_\ell(x)=\det(x_i^{\ell_j})$, 则
\[
 s_\lambda(x)=\frac{A_{\lambda+\delta}(x)}{A_\delta(x)}
 =\det(h_{\lambda_i-i+j}(x)),\qquad
 \operatorname{ch}(\chi_{V_\lambda})=s_\lambda.
\]
并有 $p_1s_\mu=\sum_{\lambda\in A(\mu)}s_\lambda$, 其中 $A(\mu)$ 为加一个格子所得Young图的集合.
\end{lemma}
\begin{proof}
分子交错, 被Vandermonde多项式 $A_\delta=\prod_{i<j}(x_i-x_j)$ 整除, 商为对称多项式. 两套各 $N$ 个变量满足Cauchy恒等式
\[
 \sum_\lambda s_\lambda(x)s_\lambda(y)
 =\prod_{i,j}(1-x_iy_j)^{-1}
 =\sum_\mu\frac{p_\mu(x)p_\mu(y)}{z_\mu}.       \tag{H8.2}
\]
第一等式可纯用行列式证明: 对矩阵 $((1-x_iy_j)^{-1})$, 将每项展开为 $\sum_{r\ge0}x_i^ry_j^r$, 用Cauchy--Binet公式, 行列式变成对严格递减非负指数 $\ell$ 的 $A_\ell(x)A_\ell(y)$ 之和. 另一方面Cauchy行列式公式给其值为 $A_\delta(x)A_\delta(y)/\prod_{i,j}(1-x_iy_j)$. 该公式可将两边乘分母验证: 分子分别对 $x,y$ 交错, 每个变量次数至多 $N-1$, 因而必为两个Vandermonde的常数倍, 比较最高交错项得常数1. 两边相除并写 $\ell=\lambda+\delta$, 即得第一等式. 第二等式是对数展开后逐项指数化.

由(H8.2), $s_\lambda$ 在上述 $p$ 内积下构成正交归一基. 再将(H8.2)乘 $A_\delta(y)$, 提取 $y^{\lambda+\delta}$ 系数. 右侧的乘积是 $\prod_i\sum_{r\ge0}h_r(x)y_i^r$, 与 $A_\delta(y)$ 相乘并展开行列式, 系数恰为 $\det(h_{\lambda_i-i+j}(x))$; 左侧只有 $A_{\lambda+\delta}(y)$ 的恒等排列贡献该严格递减指数, 给 $s_\lambda(x)$. 因而得到行列式公式.

每个 $h_r$ 对应一个平凡特征标, 乘积对应外张量诱导. 所以该行列式是某个整系数虚特征标的像. 范数为1使该虚特征标只能是一个不可约特征标的正号或负号. 其在单位元的值为
\[
 n!\det\left(\frac1{(\lambda_i-i+j)!}\right)
 =\frac{n!\prod_{i<j}(\ell_i-\ell_j)}{\prod_i\ell_i!}>0,
 \qquad \ell_i=\lambda_i+N-i.                   \tag{H8.3}
\]
负阶乘的倒数按零解释. 为验证第二等式, 每行乘 $\ell_i!$, 各列成为次数 $N-j$、首项系数1的下降阶乘多项式; 列消元化为Vandermonde行列式. 第一等式来自提取 $p_1^n$ 的系数, 因为令其他 $p_r$ 为零后 $h_r=p_1^r/r!$. 正性使它为实际不可约特征标.

还需说明其标签与 $V_\lambda$ 一致. 在字典序下 $s_\lambda$ 的最高单项式为 $x^\lambda$, 系数1: 比较上述分子、分母的最高单项式即可. 从(H8.2)提取 $y^\mu$ 系数, 得
\[
 h_{\mu_1}\cdots h_{\mu_N}
 =\sum_\lambda [y^\mu]s_\lambda(y)\,s_\lambda(x).
\]
因而行置换诱导模 $M_\mu=\C[S_n]a_\mu$ 的分解只含标签 $\lambda\ge\mu$, 标签 $\mu$ 的重数为1. 对Young构造, 上一引理的消去论证给 $a_\mu\C[S_n]c_\lambda=0$ ($\mu>\lambda$), 且 $a_\lambda\C[S_n]c_\lambda=\C c_\lambda$. 这两个空间的维数就是 $\Hom(M_\mu,V_\lambda)$ 的维数 (同态由幂等元 $a_\mu$ 的像决定). 因而Young构造也有同样的三角分解及对角重数1. 从最大分拆开始向下归纳, 两种不可约标签逐个一致, 证明特征标公式.

最后 $p_1A_\ell=\sum_j A_{\ell+e_j}$. 这是将行列式按排列展开后, 把额外的 $x_i$ 因子归到其对应的指数列得到的. 若 $\ell_j+1=\ell_{j-1}$, 该项两列相同而为零; 其余项恰对应在 $\mu$ 的可加位置增加一格. 除以 $A_\delta$, 得所需加格公式.
\end{proof}

\begin{exercise}\label{hw:4.39}
\textbf{MIT 18.712 (2010), Homework 8, 原题 4.39; 旧PDF第59页.}
求对称群 $S_n$ 全部不可约表示维数的和.
\end{exercise}
\begin{solution}
由引理\ref{hw:young-tools}, 每个不可约表示有实矩阵实现. 在其实空间上平均正定内积, 复双线性延拓后是非退化不变对称型, 所以全部为实型.

说明所用Frobenius--Schur计数. 对不可约 $L$, 投影 $P=|G|^{-1}\sum_g\rho(g)\otimes\rho(g)$ 投到 $(L\otimes L)^G$. 翻转 $\tau(v\otimes w)=w\otimes v$ 的迹满足 $\tr(\tau(A\otimes A))=\tr(A^2)$, 故 $\tr(\tau P)=|G|^{-1}\sum_g\chi_L(g^2)$. 对实型表示, 非退化对称不变型的逆对应一个对称不变张量, Schur引理又使不变张量空间为一维, 故此迹为1. 令 $r(g)=\#\{h:h^2=g\}$. 则 $\langle r,\chi_L\rangle=|G|^{-1}\sum_h\overline{\chi_L(h^2)}=1$ (这里特征标实值). 将类函数 $r$ 在全部不可约特征标基下展开, 取单位元的值, 得
\[
 \sum_{\lambda\vdash n}\dim V_\lambda
 =\#\{\sigma\in S_n:\sigma^2=1\}
 =\sum_{k=0}^{\lfloor n/2\rfloor}\frac{n!}{2^k k!(n-2k)!}.
\]
最后一式是先选 $2k$ 个被移动的点, 再将它们分成 $k$ 个无序对子形成互不相交的换位; 其余点固定. 此处计数包含恒等置换, 对应 $k=0$.
\end{solution}

\begin{exercise}\label{hw:4.50}
\textbf{MIT 18.712 (2010), Homework 8, 原题 4.50; 旧PDF第63页.}
对Young图 $\mu$, 令 $A(\mu)$ 为加一格所得图, $R(\mu)$ 为删一格所得图.
(a) 若 $|\mu|=n$, 证明 $\Res_{S_{n-1}}^{S_n}V_\mu=\bigoplus_{\lambda\in R(\mu)}V_\lambda$.
(b) 若 $|\mu|=n-1$, 证明 $\Ind_{S_{n-1}}^{S_n}V_\mu=\bigoplus_{\lambda\in A(\mu)}V_\lambda$.
\end{exercise}
\begin{solution}
先证(b). 把 $S_{n-1}$ 看作 $S_{n-1}\times S_1$, 其 $S_1$ 因子平凡. 特征标映射的诱导乘法公式及引理\ref{hw:schur-tools}给
\[
 \operatorname{ch}\left(\chi_{\Ind V_\mu}\right)
 =s_\mu h_1=s_\mu p_1=\sum_{\lambda\in A(\mu)}s_\lambda.
\]
映射单射, 特征标决定完全可约表示, 故得到(b), 每项重数1. 对(a), Frobenius互反给
\[
 [\Res V_\mu:V_\lambda]=[V_\mu:\Ind V_\lambda]
 =\begin{cases}1,&\mu\in A(\lambda),\\0,&\text{其他}.
 \end{cases}
\]
加格与删格互逆, 于是正好得到(a). 前置证明中的行列式与加格公式补齐了分支规则的依据, 不仅引用其名称.
\end{solution}

\begin{exercise}\label{hw:4.51}
\textbf{MIT 18.712 (2010), Homework 8, 原题 4.51; 旧PDF第63页.}
Young图的内容为 $c(\lambda)=\sum_j\sum_{i=1}^{\lambda_j}(i-j)$. 令 $C=\sum_{1\le i<j\le n}(ij)\in\C[S_n]$. 证明 $C$ 在Specht模 $V_\lambda$ 上作用为标量 $c(\lambda)$.
\end{exercise}
\begin{solution}
换位之和 $C$ 是中心元, 因而由Schur引理在 $V_\lambda$ 上为某标量 $\gamma$, 特别有 $C a_\lambda b_\lambda=\gamma a_\lambda b_\lambda$. 比较单位元的系数: $a_\lambda b_\lambda$ 的该系数为 $1/(|P_\lambda||Q_\lambda|)$; $C a_\lambda b_\lambda$ 的该系数为所有换位在 $a_\lambda b_\lambda$ 中的系数之和.

一个换位 $t$ 若可写成 $pq$, $p\in P_\lambda,q\in Q_\lambda$, 则它必在行群或列群中. 确实, 对被 $t$ 固定的数字 $a$, $q(a)$ 同时位于 $a$ 的列和 $a$ 的行 (因为 $p(q(a))=a$), 行列交集只有该格, 所以 $q(a)=a$ 且 $p(a)=a$. 因而 $p,q$ 均只可能作用于被交换的那两个数字, 至多一个非平凡, 因为两不同格不可能同时同行又同列. 反之同行换位的系数为正 $1/(|P||Q|)$, 同列换位为负同一值, 其他为零. 故
\[
 \gamma=\sum_j\binom{\lambda_j}{2}-\sum_i\binom{\lambda'_i}{2}
 =\sum_{(i,j)\in\lambda}\bigl((i-1)-(j-1)\bigr)=c(\lambda).
\]
这里 $\lambda'_i$ 是第 $i$ 列的高度. 最后的逐格计数将行内左侧格数与列内上方格数相减, 与原题内容符号完全一致.
\end{solution}
\originalsection{Homework 9: 内容与矩阵不变量}

\begin{exercise}\label{hw:4.52}
\textbf{MIT 18.712 (2010), Homework 9, 原题 4.52; 旧PDF第63页.}
(a) 对 $S_n$ 的任意有限维复表示 $V$, 证明 $E=(12)+(13)+\cdots+(1n)$ 可对角化, 特征值为整数且介于 $1-n$ 与 $n-1$ 之间.
(b) 证明 $E$ 在 $V_\lambda$ 上为标量当且仅当 $\lambda$ 是矩形Young图, 并计算标量.
\end{exercise}
\begin{solution}
令 $C_n$ 为全部换位之和, $C_{n-1}$ 为仅交换数字 $2,\ldots,n$ 的换位之和, 则 $E=C_n-C_{n-1}$. 对不可约 $V_\lambda$, $C_n$ 作用为 $c(\lambda)$; 在限制分解的 $V_\mu$, $\mu\in R(\lambda)$, 上, $C_{n-1}$ 作用为 $c(\mu)$. 因而 $E$ 在这一分量的标量为删去那一格的内容 $c(\lambda)-c(\mu)=i-j$. 完全可约性与分支规则给出直和分解, 故 $E$ 可对角化. 每个格子至少在第1行和第1列, 至多在第 $n$ 行或第 $n$ 列, 因而 $1-n\le i-j\le n-1$. 任意表示是不可约表示的直和, 结论随之成立.

(b) 可删格位于满足 $\lambda_j>\lambda_{j+1}$ 的行末, 其内容是 $\lambda_j-j$. 对两条这样的行 $j<k$, 有 $\lambda_j-j>\lambda_k-k$, 所以不同可删格的内容互异. 因而 $E$ 为标量恰好等价于只有一个可删格, 即所有非零行长相等, Young图为矩形. 若矩形有 $a$ 列、$b$ 行, 唯一可删格为 $(a,b)$, 标量为 $a-b$. $n=1$ 时 $E=0$, 对一格矩形公式仍成立.
\end{solution}

\begin{exercise}\label{hw:4.68}
\textbf{MIT 18.712 (2010), Homework 9, 原题 4.68; 旧PDF第68页.}
(a) 证明 $V'_\lambda=\C[S_n]b_\lambda a_\lambda$ 与 $V_\lambda=\C[S_n]a_\lambda b_\lambda$ 同构.
(b) 令代数自同构 $\varphi(s)=\sgn(s)s$. 说明它将一个表示扭曲为与符号表示 $\C_-$ 的张量积, 且 $\varphi(\C[S_n]a)=\C[S_n]\varphi(a)$. 由(a)证明 $V_\lambda\otimes\C_-\cong V_{\lambda'}$, 其中 $\lambda'$ 为转置Young图.
\end{exercise}
\begin{solution}
(a) 写 $a=a_\lambda,b=b_\lambda,\kappa=\kappa_\lambda>0$. 由引理\ref{hw:young-tools}, $(ab)^2=\kappa ab$. 取共轭转置, 因 $a,b$ 自伴且 $\kappa$ 实, 同样有 $(ba)^2=\kappa ba$. 定义 $f(x)=xa$, $g(y)=yb$. 若 $x=uab$, 则 $xa=(ua)ba\in V'$, 若 $y=uba$, 则 $yb=(ub)ab\in V$, 所以良定义, 且两者均与左群作用交换. 对这两种元素分别有
\[
 gf(x)=xab=\kappa x,\qquad fg(y)=yba=\kappa y.
\]
因为 $\kappa\ne0$, $f$ 可逆, 逆为 $\kappa^{-1}g$, 证明同构.

(b) $\varphi$ 保持乘法且自身互逆. 对任意左理想 $L$, 限制映射 $\varphi:L\to\varphi(L)$ 满足 $\varphi(sx)=\sgn(s)s\varphi(x)$, 因而是符号扭曲后的表示同构. 对一般表示同样用 $\rho\circ\varphi$ 定义扭曲, 在群元素上作用即 $\sgn(s)\rho(s)$. 自同构的满射性还给 $\varphi(\C[S_n]a)=\C[S_n]\varphi(a)$.

将同一张填数表转置, 行群与列群交换. 因而 $\varphi(a_\lambda)=b_{\lambda'}$, $\varphi(b_\lambda)=a_{\lambda'}$, 连同归一化因子均恰好对应. 所以
\[
 V_\lambda\otimes\C_-\cong
 \C[S_n]b_{\lambda'}a_{\lambda'}\cong V_{\lambda'},
\]
最后一步使用(a), 不把两个Young算子的顺序直接视为相同.
\end{solution}

\begin{exercise}\label{hw:4.69}
\textbf{MIT 18.712 (2010), Homework 9, 原题 4.69; 旧PDF第68页.}
令 $R_{k,N}$ 为 $k$ 个复 $N\times N$ 矩阵 $X_1,\ldots,X_k$ 的多项式函数中, 对同时共轭 $X_i\mapsto gX_ig^{-1}$ 不变者组成的代数. 对任意有限字 $w$ 定义 $T_w=\tr(w(X_1,\ldots,X_k))$. 证明 $R_{k,N}$ 由全部 $T_w$ 生成.
\end{exercise}
\begin{solution}
先证明本题所需的Schur--Weyl交换代数结论, 不需要完整的多项式表示分类. 令 $U=\C^N$, $d\ge0$. 置换群 $S_d$ 在 $U^{\otimes d}$ 上交换张量因子, $\GL(U)$ 以 $g^{\otimes d}$ 作用. 通过 $\End(U^{\otimes d})\cong\End(U)^{\otimes d}$, 与全部置换交换的算子恰为对张量因子置换不变的张量, 即 $\Sym^d(\End U)$.

这些对称张量由 $g^{\otimes d}$, $g\in\GL(U)$, 线性张成. 证明如下: 一个线性泛函若在所有可逆 $g$ 的纯幂张量上为零, 则多项式 $P(X)=\ell(X^{\otimes d})$ 在所有可逆矩阵处为零. 对任意矩阵 $X$, $X+tI$ 仅在有限多个 $t$ 处不可逆, 故 $P(X+tI)$ 这一一元多项式恒为零, 特别 $P(X)=0$. 对纯幂作极化, 即在 $P(\sum t_iX_i)$ 中提取 $t_1\cdots t_d$ 系数, 知 $\ell$ 在所有对称张量上为零. 有限维线性代数给出所需张成性.

因此 $S_d$ 的交换代数正是 $g^{\otimes d}$ 的线性张成代数. 半单有限群表示的双交换代数是原群代数的像: 将空间写为 $\bigoplus_\alpha L_\alpha\otimes M_\alpha$, 原群代数像为 $\bigoplus_\alpha\End(L_\alpha)\otimes I$, 其交换代数为 $\bigoplus_\alpha I\otimes\End(M_\alpha)$, 再取交换代数恢复原像. 故
\[
 \End_{\GL(U)}(U^{\otimes d})
 =\text{张量因子置换算子的线性张成}.              \tag{H9.1}
\]
论证对 $N<d$ 同样成立; 置换算子此时可能线性相关, 但生成结论不变.

现在令 $f$ 为不变多项式. 按各矩阵的次数分解为多齐次部分, 共轭保持次数, 所以每部分仍不变. 对多次数 $(d_1,\ldots,d_k)$, $d=\sum d_i$, 对每组变量作极化, 得一个 $d$ 重线性函数, 在属于同一 $X_i$ 的 $d_i$ 个位置内对称, 且仍对同时共轭不变. 在复数域中, 极化与令同组各变量相等互逆 (包含相应阶乘归一化), 所以只需描述全部不变多线性函数.

迹配对 $\End(U)\times\End(U)\to\C$, $(A,B)\mapsto\tr(AB)$, 非退化且共轭不变. 因而 $d$ 重线性函数唯一写成
\[
 (Y_1,\ldots,Y_d)\longmapsto
 \tr_{U^{\otimes d}}\bigl(A(Y_1\otimes\cdots\otimes Y_d)\bigr).
\]
它不变当且仅当 $A$ 与 $g^{\otimes d}$ 交换. 由(H9.1), $A$ 是置换算子 $P_\sigma$ 的线性组合. 对单个置换, 在张量基中展开迹, 每个循环给一串首尾相接的矩阵指标, 因而
\[
 \tr\bigl(P_\sigma(Y_1\otimes\cdots\otimes Y_d)\bigr)
 =\prod_{(i_1\cdots i_r)\text{ 为 }\sigma\text{ 的循环}}
       \tr(Y_{i_r}\cdots Y_{i_1}).                 \tag{H9.2}
\]
这里取 $P_\sigma(v_1\otimes\cdots\otimes v_d)
 =v_{\sigma^{-1}(1)}\otimes\cdots\otimes v_{\sigma^{-1}(d)}$; 换一种置换约定只会反转循环中矩阵的排列, 仍属于某个字的迹. 例如固定点给 $\tr(Y_i)$, 二循环给 $\tr(Y_jY_i)$, 三循环给 $\tr(Y_kY_jY_i)$, 可直接复核指标收缩.

将同一组的 $Y$ 全部代回 $X_i$, (H9.2)每一项就是若干 $T_w$ 的乘积. 各多齐次部分因而均为这些迹的多项式, 相加得到 $f$ 也在它们生成的代数中. 反过来每个 $T_w$ 显然对同时共轭不变, 因为共轭在字的乘积中逐项抵消且迹不变. 零次数部分是常数, 已包含在生成代数中. 这完成两个方向的证明.
\end{solution}

% 题面依据2011-02-01公开旧讲义，PDF78--81、96--97。
% 未见公开官方解答；本文件解答为编者独立编写，印刷问题显式登记。
\originalsection{Homework 10：箭图与 McKay 图}
本次作业的原题是旧讲义的5.1--5.5，其中5.5的末两问在原PDF第81页。题面译自公开旧讲义，以下解答由AI辅助独立编写，并非MIT官方答案。为解答这些题，本节补充箭图、Dynkin图和根格的必要知识；这些局部补充不表示本书已覆盖旧讲义的箭图、李代数或一般有限维代数全部专题。以下默认箭图表示有限维，底域代数闭；涉及特征标和有限旋转群时取底域为 $\C$。

\begin{exercise}\label{hw:5.1}
\textbf{MIT 18.712（2010）Homework 10 原题5.1；旧PDF第78页。}
设 $k$ 是域，$k(y_1,\ldots,y_m)$ 是 $m$ 个变量上的有理函数域。若有单射的 $k$-代数同态
\[
 f:k[x_1,\ldots,x_n]\longrightarrow k(y_1,\ldots,y_m),
\]
证明 $m\ge n$。原题提示：比较 $x_i$ 中次数为 $N$ 的多项式空间及其像随 $N\to\infty$ 的维数增长。再推得，若有 $k$-域嵌入 $k(x_1,\ldots,x_n)\to k(y_1,\ldots,y_m)$，仍有 $m\ge n$。
\end{exercise}
\begin{solution}
为得到足够强的增长率，使用\textbf{总次数至多} $N$ 的空间 $W_N$，而非只取齐次次数恰为 $N$ 的空间。其单项式基给出
\[
 \dim_k W_N=\binom{N+n}{n}.
\]
把所有 $f(x_i)$ 写成共同分母的形式 $p_i/q$，其中 $p_i,q\in k[y_1,\ldots,y_m]$，$q\ne0$。取 $D\ge1$，使这些多项式的总次数都不超过 $D$。若 $x^\alpha$ 的次数 $|\alpha|\le N$，则
\[
 q^N f(x^\alpha)=q^{N-|\alpha|}\prod_i p_i^{\alpha_i}
\]
是次数至多 $DN$ 的多项式。因此，$f$ 的单射性和乘以 $q^N$ 的单射性共同给出
\[
 \binom{N+n}{n}\le\binom{DN+m}{m}\qquad(N\ge0).
\]
左侧关于 $N$ 的次数为 $n$，右侧为 $m$，且最高次系数都为正。若 $n>m$，当 $N$ 足够大时不等式不可能成立，故 $n\le m$。域嵌入限制到多项式子代数仍然单射，第二个结论随即成立。
\end{solution}

\begin{exercise}\label{hw:5.2}
\textbf{MIT 18.712（2010）Homework 10 原题5.2；旧PDF第78页。}
设 $k$ 代数闭，$G=\GL_n(k)$，$V$ 是 $G$ 的多项式表示。证明：若 $G$ 在 $V$ 上只有有限个轨道，则 $\dim V\le n^2$。
\begin{enumerate}[label=(\alph*)]
\item 取 $V$ 的线性坐标 $x_1,\ldots,x_N$。称子集 $X\subset V$ 为Zariski稠密，如果在 $X$ 上为零的多项式只能是零多项式。证明有限个轨道中至少有一个Zariski稠密轨道。
\item 据此构造域嵌入 $k(x_1,\ldots,x_N)\to k(g_{pq})$，再用原题5.1。
\item 推广到 $G=\GL_{n_1}(k)\times\cdots\times\GL_{n_s}(k)$。
\end{enumerate}
\end{exercise}
\begin{solution}
\textbf{(a)} 若所有轨道 $O_1,\ldots,O_t$ 都不稠密，可为每个轨道取非零多项式 $F_j$，使 $F_j|_{O_j}=0$。则非零多项式 $\prod_jF_j$ 在整个 $V=k^N$ 上为零。代数闭域是无限域，而无限域上在所有点取零值的多项式必为零：对变量数归纳，把它当作最后一个变量的一元多项式，并用一元非零多项式的根数有限。这一矛盾证明了结论。

\textbf{(b)} 取稠密轨道中的点 $v$。轨道映射 $g\mapsto gv$ 把坐标函数 $x_i$ 拉回为矩阵坐标 $g_{pq}$ 的多项式，因而定义同态
\[
 k[x_1,\ldots,x_N]\longrightarrow k(g_{pq}),\qquad F\longmapsto[g\mapsto F(gv)].
\]
若像为零，$F$ 在稠密轨道上消失，故 $F=0$；同态单射，并能唯一延拓到分式域。目标有 $n^2$ 个变量，原题5.1给出 $N\le n^2$。

\textbf{(c)} 相同证明给出
\[
 \boxed{\dim V\le\sum_{j=1}^s n_j^2.}
\]
下一题需要稍宽的版本：作用矩阵的系数可以是矩阵坐标及行列式的逆所组成的有理函数。此时轨道映射的拉回仍落在同一有理函数域，以上单射和变量数论证完全保留。因此结论也适用于这种有理表示，而不必强求所有系数都是多项式。
\end{solution}

\begin{exercise}\label{hw:5.3}
\textbf{MIT 18.712（2010）Homework 10 原题5.3；旧PDF第78--80页。}
设 $\Gamma$ 是有限连通无自环图，允许重边，顶点为 $1,\ldots,N$。令 $r_{ij}$ 为 $i,j$ 间的边数，$R_\Gamma=(r_{ij})$，$A_\Gamma=2I-R_\Gamma$。若 $A_\Gamma$ 定义的实二次型正定，称 $\Gamma$ 为Dynkin图。证明这类图恰为 $A_n$（$n\ge1$）、$D_n$（$n\ge4$）、$E_6,E_7,E_8$。

为准确说明原题的图形，$A_n$ 是 $n$ 个顶点的链；$D_n$ 是从一个三价顶点伸出长度为 $n-3,1,1$ 的三条臂；$E_6,E_7,E_8$ 的三臂长度分别为 $(2,2,1),(3,2,1),(4,2,1)$。这里臂长是边数，中心不重复计入。下图逐一重画旧讲义第78--79页的有限型图：
\begin{center}
\begin{tikzpicture}[scale=.64,every node/.style={circle,draw,fill=white,inner sep=1.4pt}]
\node[draw=none] at (-1,0) {$A_n$};
\foreach \i in {0,1,4,5}{\node (a\i) at (\i,0) {};}
\draw (a0)--(a1);\draw[dashed] (a1)--(a4);\draw (a4)--(a5);
\node[draw=none] at (-1,-1.5) {$D_n$};
\foreach \i in {0,1,4,5}{\node (d\i) at (\i,-1.5) {};}
\node (db) at (4,-2.2) {};\draw (d0)--(d1);\draw[dashed] (d1)--(d4);\draw (d4)--(d5);\draw (d4)--(db);
\node[draw=none] at (-1,-3.5) {$E_6$};
\foreach \i in {0,...,4}{\node (s\i) at (\i,-3.5) {};}
\node (sb) at (2,-4.2) {};\foreach \i [evaluate=\i as \j using int(\i+1)] in {0,...,3}{\draw (s\i)--(s\j);}\draw (s2)--(sb);
\node[draw=none] at (-1,-5.5) {$E_7$};
\foreach \i in {0,...,5}{\node (t\i) at (\i,-5.5) {};}
\node (tb) at (3,-6.2) {};\foreach \i [evaluate=\i as \j using int(\i+1)] in {0,...,4}{\draw (t\i)--(t\j);}\draw (t3)--(tb);
\node[draw=none] at (-1,-7.5) {$E_8$};
\foreach \i in {0,...,6}{\node (u\i) at (\i,-7.5) {};}
\node (ub) at (4,-8.2) {};\foreach \i [evaluate=\i as \j using int(\i+1)] in {0,...,5}{\draw (u\i)--(u\j);}\draw (u4)--(ub);
\end{tikzpicture}
\end{center}
\begin{enumerate}[label=(\alph*)]
\item 用按行展开建立递推，计算 $A_N,D_N$ 的Cartan矩阵行列式，并用Sylvester判据证明正定。
\item 计算 $E_6,E_7,E_8$ 的行列式并证明正定。
\item 证明Dynkin图没有圈。原题提示：圈图的Cartan矩阵各行之和为零。
\item 证明其顶点度数至多三，且三价顶点至多一个。原题给出的禁图为四臂星形，以及两个三价顶点由一条链连接、各自另接两个叶子的图。
\item 证明三臂长度分别为 $(2,2,2),(3,3,1),(5,2,1)$ 的图的Cartan矩阵行列式为零。这是原题第79--80页的三个带数字禁图。
\item 由以上结果完成有限型分类。
\item 分类仿射Dynkin图。原题把它定义为Cartan矩阵半正定的连通无自环图，并要求它们恰为前面的禁图。
\end{enumerate}
\end{exercise}
\begin{remark}
原题(g)的定义少了“非正定”这一条件：正定矩阵本来也是半正定矩阵，照字面会把有限型全部包含进去。以下将“仿射”按半正定而非正定解释；若严格沿用原题字面，则答案应是有限型与仿射型的并集。允许重边时还须补入两个顶点之间的双边图 $\widetilde A_1$，它是长度二的圈；旧图中没有单独画出这个边界情形。
\end{remark}
\begin{solution}
先记两个可重复使用的计算。若树 $T$ 中的叶子 $v$ 邻接 $u$，沿叶子所在行、列展开，有
\[
 \Delta(T)=2\Delta(T-v)-\Delta(T-\{u,v\}),\qquad
 \Delta(T)=\det A_T.
\]
删掉顶点后若图不连通，行列式就是各连通分支行列式的乘积；空图的行列式取一。

\textbf{(a)} 令 $d_n=\Delta(A_n)$，则 $d_0=1,d_1=2$，且 $d_n=2d_{n-1}-d_{n-2}$，归纳得到 $d_n=n+1$。链按顺序编号时，各顺序主子式为 $2,3,\ldots,n+1$，故由Sylvester判据正定。对 $D_n$，删去其中一个短臂叶子，剩下 $A_{n-1}$；再删去中心，剩下 $A_{n-3}$ 与一个孤立点。因此
\[
 \Delta(D_n)=2n-2(n-2)=4.
\]
先编号长臂、中心和一个短臂，最后编号另一个短臂，则前 $n-1$ 个顺序主子式来自链，最后一个为四，全部为正。

\textbf{(b)} 将三臂图记为 $T(a,b,c)$。一条长度为 $a$ 的臂对应链矩阵 $A_{A_a}$，其靠中心端的逆矩阵对角元为
\[
 (A_{A_a}^{-1})_{11}=\frac{\Delta(A_{a-1})}{\Delta(A_a)}=\frac a{a+1}.
\]
先消去三条臂，中心的Schur补为
\[
 2-\frac a{a+1}-\frac b{b+1}-\frac c{c+1}
 =\frac1{a+1}+\frac1{b+1}+\frac1{c+1}-1.
\]
于是
\[
 \Delta(T(a,b,c))=(a+1)(b+1)(c+1)
 \left(\frac1{a+1}+\frac1{b+1}+\frac1{c+1}-1\right).
\]
代入 $(2,2,1),(3,2,1),(4,2,1)$，依次得到 $3,2,1$。先给水平链编号，最后给短臂叶子编号，前面的主子式都为链的正行列式，最后一个也为正，故三图正定。

\textbf{(c)--(e)} 若 $A_\Gamma$ 正定，则任何子图也必须正定。这里对子图允许删边：若子图非正定，先取非零向量 $x$ 使其二次型不正，再改为 $d=|x|$；由于非对角项非正，二次型不会增加。将 $d$ 在其余顶点上延伸为零，增加边只会再减去非负项。因此找到这样的子图就足以排除正定性。

圈图取所有 $d_i=1$，则 $A_\Gamma d=0$。若有重边，两个端点上的矩阵为 $\bigl(\begin{smallmatrix}2&-r\\-r&2\end{smallmatrix}\bigr)$，当 $r\ge2$ 时在 $(1,1)$ 上不正，所以正定图必是没有重边的树。

四臂星形取中心权二、叶子权一，满足每个顶点的“两倍自身权等于相邻权之和”。有两个三价顶点时，取连接它们的整条链上的权为二，两端另外四个叶子的权为一，同样满足 $Ad=0$。若原树有更多分叉，取两分叉间最短路径，并截取两端各两条支边，仍得到该禁图。因此最大度数不超过三，三价顶点至多一个。

三个例外禁图的非零核向量如下。图的数字和连接关系按旧讲义第79--80页重画，数字都是顶点权，而不是顶点编号：
\begin{center}
\begin{tikzpicture}[scale=.68,every node/.style={circle,fill=black,inner sep=1.7pt},every label/.style={fill=none,draw=none,font=\small,inner sep=1pt}]
\node[label=below:1] (e60) at (-2,0) {};\node[label=below:2] (e61) at (-1,0) {};\node[label=below:3] (e62) at (0,0) {};\node[label=below:2] (e63) at (1,0) {};\node[label=below:1] (e64) at (2,0) {};\node[label=right:2] (e65) at (0,1) {};\node[label=above:1] (e66) at (0,2) {};
\draw (e60)--(e61)--(e62)--(e63)--(e64);\draw (e62)--(e65)--(e66);
\node[fill=none] at (4,0.7) {$\widetilde E_6$};
\foreach \i/\w in {0/1,1/2,2/3,3/4,4/3,5/2,6/1}{\node[label=below:\w] (e7\i) at (\i-3,-2.7) {};}
\foreach \i [evaluate=\i as \j using int(\i+1)] in {0,...,5}{\draw (e7\i)--(e7\j);}
\node[label=above:2] (e77) at (0,-1.7) {};\draw (e73)--(e77);\node[fill=none] at (4,-2.1) {$\widetilde E_7$};
\foreach \i/\w in {0/1,1/2,2/3,3/4,4/5,5/6,6/4,7/2}{\node[label=below:\w] (e8\i) at (\i-4,-5.2) {};}
\foreach \i [evaluate=\i as \j using int(\i+1)] in {0,...,6}{\draw (e8\i)--(e8\j);}
\node[label=above:3] (e88) at (1,-4.2) {};\draw (e85)--(e88);\node[fill=none] at (5,-4.6) {$\widetilde E_8$};
\end{tikzpicture}
\end{center}
逐点检查 $2d_i=\sum_{j\sim i}d_j$ 即得 $Ad=0$，故三者行列式为零。前两类禁图也可带权画成
\begin{center}
\begin{tikzpicture}[scale=.68,every node/.style={circle,fill=black,inner sep=1.7pt},every label/.style={fill=none,draw=none,font=\small,inner sep=1pt}]
\foreach \i in {0,...,4}{\node[label={\i*72:1}] (c\i) at (\i*72:1) {};}
\foreach \i/\j in {0/1,1/2,2/3,3/4,4/0}{\draw (c\i)--(c\j);}
\node[fill=none] at (0,-1.8) {圈：全为一};
\node[label=below:2] (l) at (3.2,0) {};\node[label=below:2] (r) at (6,0) {};
\node[label=above:1] (ll) at (2.3,.8) {};\node[label=below:1] (ld) at (2.3,-.8) {};
\node[label=above:1] (rr) at (6.9,.8) {};\node[label=below:1] (rd) at (6.9,-.8) {};
\draw (ll)--(l)--(ld);\draw[dashed] (l)--(r);\draw (rr)--(r)--(rd);
\node[fill=none] at (4.7,-1.8) {$\widetilde D_n$：中间链全为二};
\end{tikzpicture}
\end{center}
右图的两个分叉点重合时成为 $\widetilde D_4$：中心权二，四个叶子权一；圈的两个顶点由双边相连时成为 $\widetilde A_1$。

\textbf{(f)} 没有三价顶点的树是 $A_n$。有唯一三价顶点时是 $T(a,b,c)$，设 $1\le a\le b\le c$。正定等价于
\[
 \frac1{a+1}+\frac1{b+1}+\frac1{c+1}>1.
\]
若 $a\ge2$，左侧至多一，故 $a=1$。若 $b=1$，任意 $c\ge1$ 都成立，给出 $D_{c+3}$；若 $b=2$，则必须 $c=2,3,4$，给出 $E_6,E_7,E_8$；若 $b\ge3$，左侧又至多一。分类完成。

\textbf{(g)} 对上述每个禁图，已写出正核向量 $d$。把重边按条计数，有恒等式
\[
 x^{\mathsf T}Ax=
 \sum_{\{i,j\}\text{为边}}d_i d_j
 \left(\frac{x_i}{d_i}-\frac{x_j}{d_j}\right)^2.
\]
展开后，$x_i^2$ 的系数为 $\sum_{j\sim i}d_j/d_i=2$，交叉项为每条边的 $-2x_ix_j$，故恒等式成立。它证明半正定；连通性又使核恰为 $\R d$。

反之，设连通图的 $A$ 半正定且奇异。取非零核向量 $x$，则
$|x|^{\mathsf T}A|x|\le x^{\mathsf T}Ax=0$，故 $|x|$ 也在核中。若其某坐标为零，核方程迫使所有邻点坐标为零，连通性将迫使全部为零。因此有严格正的核向量。一个具有正核向量的禁图不能真包含在更大的连通半正定图中：增加其内部边立即使该向量的二次型为负；增加相邻顶点时，把旧向量延伸为零虽仍有零二次型，却不再满足核方程，而半正定矩阵的零二次型向量必须在核中，矛盾。

如果该图没有任何禁图，前面的分类已经证明它属于正定的有限型，与奇异性矛盾；所以它恰是一个禁图。由此全部仿射型为圈 $\widetilde A_{r-1}$（$r\ge3$）、双边 $\widetilde A_1$、$\widetilde D_n$（$n\ge4$）以及 $\widetilde E_6,\widetilde E_7,\widetilde E_8$。
\end{solution}

\begin{exercise}\label{hw:5.4}
\textbf{MIT 18.712（2010）Homework 10 原题5.4；旧PDF第80页。}
设 $Q$ 的顶点集为 $D$。若 $Q$ 的不可分解有限维表示只有有限个同构类，称 $Q$ 为有限型。令 $b_{ij}$ 是从 $i$ 到 $j$ 的箭头数。Gabriel定理说：连通箭图是有限型，当且仅当忘掉箭头方向后得到Dynkin图。本题只证明必要性。
\begin{enumerate}[label=(\alph*)]
\item 证明有限型迫使相应二次型在 $x_i\in\mathbb Q$、$x_i\ge0$ 且不全为零的向量上为正。原题印刷的二次型是
\[
 q_{\text{印}}(x)=\sum_i x_i^2-\frac12\sum_{i,j}b_{ij}x_ix_j.
\]
提示先取非负整数 $n_i$，考虑所有维数为 $n_i$ 的表示组成的空间 $W$，使 $\prod_i\GL_{n_i}(k)$ 在 $W\oplus k$ 上具有有限轨道，利用原题5.2，再推广到任意整数和有理数。
\item 推出二次型正定。原题提示利用 $\mathbb Q$ 在 $\R$ 中稠密。
\item 排除自环，再利用原题5.3得到Gabriel定理的必要性。
\end{enumerate}
\end{exercise}
\begin{remark}
按原题给定的\textbf{有向}箭头数 $b_{ij}$，用于Gabriel定理的二次型应为
\[
 q_Q(x)=\sum_i x_i^2-\sum_{i,j}b_{ij}x_ix_j.
\]
原印刷式的 $1/2$ 是多余的。若改用无向对称邻接矩阵 $R$，才可写成 $q_Q(x)=\sum_i x_i^2-\tfrac12\sum_{i,j}R_{ij}x_ix_j$（此处无自环）。例如Kronecker箭图的 $b_{12}=2$，原印刷式成为 $x_1^2+x_2^2-x_1x_2$，正定但图并非有限型，不能据此证明(c)。以下使用修正后的 $q_Q$，并完整证明原题意图中的必要性。
\end{remark}
\begin{solution}
箭图的一个表示给每个顶点 $i$ 一个向量空间 $V_i$，给每条箭头 $i\to j$ 一个线性映射 $V_i\to V_j$；同构是在各顶点换基，并与所有箭头映射相容。固定非零维数向量 $n=(n_i)$，令
\[
 W=\bigoplus_{a:i\to j}\Hom_k(k^{n_i},k^{n_j}),\qquad
 H=\prod_{i:n_i>0}\GL_{n_i}(k).
\]
于是 $\dim W=\sum_{i,j}b_{ij}n_i n_j$，$\dim H$ 在本题的变量计数意义下是 $\sum_i n_i^2$。$H$ 以 $T_a\mapsto g_jT_ag_i^{-1}$ 作用，轨道恰是这些表示的同构类。

任一有限维表示可通过反复分解为较小维数的非零直和项，写成不可分解表示的直和。若不可分解同构类只有有限个，则固定总维数时，各类的重数有统一有限上界，因此只产生有限个同构类，$W$ 上只有有限轨道。

为了获得严格不等式，而非仅有非负性，选一个 $n_j>0$，使 $H$ 还在额外的一维空间上按 $z\mapsto(\det g_j)z$ 作用。每个 $W$ 轨道的稳定子都含有共同标量 $g_i=tI_{n_i}$；它在额外坐标上乘以 $t^{n_j}$。代数闭性保证非零数都有 $n_j$ 次根，所以该稳定子在额外坐标的非零值上作用传递。每个旧轨道恰贡献 $z=0$ 和 $z\ne0$ 两类轨道。因此 $W\oplus k$ 仍只有有限轨道。用原题5.2的有理作用版本，得到
\[
 \sum_{i,j}b_{ij}n_i n_j+1\le\sum_i n_i^2,
 \qquad\text{即 }q_Q(n)\ge1.
\]
若 $n$ 是任意非零整数向量，由箭头数非负可知
\[
 q_Q(n)\ge q_Q(|n|)>0.
\]
给有理向量乘以公共分母，再用二次齐次性，证明(a)。

\textbf{(b)} 连续性和有理点稠密首先只证明 $q_Q$ 半正定，\textbf{不能直接证明正定}。若它奇异，其矩阵有有理系数，所以高斯消元给出非零的有理核向量，与(a)矛盾。因此它确实正定。这一步补足了原提示中未写出的关键论证。

\textbf{(c)} 若顶点 $i$ 上有自环，令 $x_i=1$、其余为零，则 $q_Q(x)=1-b_{ii}\le0$，与正定性矛盾。无自环时，忘掉方向后的邻接矩阵满足
\[
 q_Q(x)=\frac12x^{\mathsf T}(2I-R)x.
\]
原题5.3的分类于是表明底图是Dynkin图。本题只得到Gabriel定理的“有限型必为Dynkin型”方向；反方向并未由这段维数论证证明。
\end{solution}

\begin{exercise}\label{hw:5.5}
\textbf{MIT 18.712（2010）Homework 10 原题5.5；旧PDF第80--81页。}
设 $G\ne\{1\}$ 是 $\mathrm{SU}(2)$ 的有限子群，$V$ 是由该嵌入得到的二维表示，$V_i$（$i\in I$）是全部不可约表示。令 $r_{ij}$ 为 $V_i$ 在 $V\otimes V_j$ 中的重数。
\begin{enumerate}[label=(\alph*)]
\item 证明 $r_{ij}=r_{ji}$。
\item McKay图 $M(G)$ 以 $i\in I$ 为顶点，并在 $i,j$ 之间连 $r_{ij}$ 条边。证明该图连通。原题提示利用原题3.26。
\item 证明它是原题5.3中的仿射Dynkin图。先证 $A=(2\delta_{ij}-r_{ij})$ 半正定而非正定，再证没有自环。原题提示：令 $F=\sum_i x_i\chi_{V_i}$，计算 $\langle(2-\chi_V)F,F\rangle$；非循环的有限 $\mathrm{SU}(2)$ 子群包含 $-I$。
\item 原题3.24中的各有限子群对应哪一种图？
\item 利用McKay图求这些群的全部不可约表示维数，证明它们就是原图中的顶点数字，并与原题3.24所得 $\mathrm{SO}(3)$ 子群比较。
\end{enumerate}
\end{exercise}
\begin{solution}
\textbf{(a)} $V$ 上的交替形式 $\omega(u,v)=\det(u,v)$ 在 $G$ 下不变且非退化，给出 $V\cong V^*$。由张量积与对偶的自然同态关系，
\[
 \Hom_G(V_i,V\otimes V_j)\cong\Hom_G(V\otimes V_i,V_j).
\]
Maschke定理使两边维数分别为 $r_{ij},r_{ji}$，故相等。也可用实值特征标 $\chi_V$ 和特征标内积验证这一等式。

\textbf{(b)} 嵌入保证 $V$ 忠实。原题3.26说明每个不可约表示出现在某个 $V^{\otimes n}$ 中。此处还可直接补一个特征标证明：$g\in\mathrm{SU}(2)$ 的特征值为 $e^{\mathrm i\theta},e^{-\mathrm i\theta}$，所以 $\chi_V(g)=2\cos\theta$，且 $\chi_V(g)=2$ 当且仅当 $g=1$。设不可约特征标 $\psi$ 与所有 $\chi_V^n$ 正交。按 $\chi_V$ 的有限个不同取值 $a$ 分组，得到对所有 $n\ge0$ 的等式
\[
 \sum_a a^n\sum_{g:\chi_V(g)=a}\overline{\psi(g)}=0.
\]
取 $n=0,\ldots,t-1$，Vandermonde矩阵可逆，各内层和都为零；但 $a=2$ 的内层和为 $\psi(1)=\dim\psi>0$，矛盾。于是每个不可约表示都能从平凡表示出发，沿连续张量乘以 $V$ 的分解到达，正好给出McKay图中的一条路，故图连通。

\textbf{(c)} 在不可约特征标正交基下，乘以 $\chi_V$ 的矩阵就是 $R=(r_{ij})$。对实向量 $x$ 和 $F=\sum_i x_i\chi_{V_i}$，
\[
 x^{\mathsf T}Ax=\langle(2-\chi_V)F,F\rangle
 =\frac1{|G|}\sum_{g\in G}(2-\chi_V(g))|F(g)|^2\ge0.
\]
取 $d_i=\dim V_i$，由 $\dim(V\otimes V_j)=2d_j$ 得到 $Rd=2d$，所以 $Ad=0$，不是正定矩阵。上式还说明，等号只可能出现在 $F$ 支持于单位元时；这样的类函数空间是一维的。因此 $\ker A$ 一维，正核向量 $d$ 将唯一确定全部维数，差一个共同倍数。

还须排除自环。如果 $G$ 非循环，由原题3.24的旋转群分类及 $\mathrm{SU}(2)\to\mathrm{SO}(3)$ 的核为 $\{I,-I\}$，它必包含 $-I$。具体说，如果不包含 $-I$，它没有二阶元素，因而阶数为奇数；其像是奇数阶旋转群，分类说明该像循环，而原群嵌入同一旋转轴的逆像（一个圆群），也只能是循环群。中央元 $-I$ 在每个不可约 $V_i$ 上按某个 $\varepsilon_i=\pm1$ 作用，在 $V$ 上按 $-1$ 作用。因此 $V\otimes V_i$ 的中央符号为 $-\varepsilon_i$，不能包含 $V_i$，即 $r_{ii}=0$。

若 $G$ 循环，设其阶数为 $m\ge2$，可将生成元写成 $\operatorname{diag}(\zeta,\zeta^{-1})$。全部不可约表示为 $\chi_j:g\mapsto\zeta^j$，$j\in\Z/m\Z$，且
\[
 V\otimes\chi_j=\chi_{j+1}\oplus\chi_{j-1}.
\]
当 $m>2$ 时得到 $m$ 个顶点的圈；当 $m=2$ 时两项相同，得到两个顶点之间的双边。都没有自环。现在原题5.3(g)的分类适用，得出仿射Dynkin图。

\textbf{(d)--(e)} 将正核向量按互素整数归一化，记为 $\delta$。原题5.3所画图至少有一个数字一，且平凡表示的维数也是一；又因 $d=c\delta$ 的所有坐标都是整数，$\delta$ 的某个坐标为一，故 $c$ 为正整数。平凡表示的坐标一迫使 $c=1$。因此各不可约维数确实等于所画数字。

这里把原题3.24的分类明确写成对应表；$\mathrm{Dic}_r$ 表示阶数 $4r$ 的二元二面体群，$r\ge2$，其 $r=2$ 的情形是四元数群。$2T,2O,2I$ 分别是正四面体、正八面体、正二十面体旋转群在 $\mathrm{SU}(2)$ 中的逆像。
\begin{center}\small
\begin{tabular}{lll}
\toprule
群 & McKay图 & 不可约表示维数（含重数）\\
\midrule
$C_m$, $m\ge2$ & $\widetilde A_{m-1}$ & $m$ 个一\\
$\mathrm{Dic}_r$ & $\widetilde D_{r+2}$ & 四个一，$r-1$ 个二\\
$2T$，阶二十四 & $\widetilde E_6$ & $1,1,1,2,2,2,3$\\
$2O$，阶四十八 & $\widetilde E_7$ & $1,1,2,2,2,3,3,4$\\
$2I$，阶一百二十 & $\widetilde E_8$ & $1,2,2,3,3,4,4,5,6$\\
\bottomrule
\end{tabular}\end{center}
为使图的判定可复核，而不只凭阶数猜测，再核对一维表示数。二元二面体群有表示
\[
 \langle a,b\mid a^{2r}=1,\ b^2=a^r,\ bab^{-1}=a^{-1}\rangle.
\]
交换化后 $a^2=1,b^2=a^r$，得到阶四的交换群，所以一维表示恰有四个。非循环图中只有 $\widetilde D_n$ 有四个维数一；其权平方和为 $4+4(n-3)=4(n-2)$，由群阶确定 $n=r+2$。

对三个旋转商群，可分别取生成元
\[
\begin{array}{c|cc|c}
 &\overline x&\overline y&\operatorname{ord}(\overline x\overline y)\\\hline
 A_4&(12)(34)&(123)&3\\
 S_4&(12)&(234)&4\\
 A_5&(12)(34)&(135)&5
\end{array}
\]
前一组生成四元Klein子群及其三循环，故生成 $A_4$；第二组的乘积为四循环，结合该换位可得到全部相邻换位，故生成 $S_4$。第三组所生成子群的阶被二、三、五整除，故为三十或六十；但 $A_5$ 由三循环生成，任何到二阶群的同态都使三循环取单位值，所以没有指数二子群，阶只能为六十。令 $\overline z=(\overline x\overline y)^{-1}$，则三种阶数为 $(2,3,s)$，$s=3,4,5$，且乘积为单位元。

取这些旋转的半角提升 $x,y,z$，使 $x^2=y^3=z^s=-I$；再改变 $x$ 的符号，可使 $xyz=-I$ 而不改变 $x^2$。这些提升生成整个二元群，因为它们的像生成旋转群，且 $x^2=-I$ 包含投影核。它们满足必要关系
\[
 x^2=y^3=z^s=xyz,\qquad s=3,4,5.
\]
对一维表示用加法记号记录这些关系，有 $x=y+z$、$y=2z$，故 $x=3z$、$(6-s)z=0$。全部取值由一个 $6-s$ 次单位根确定，一维表示数至多分别为三、二、一。另一方面，$A_4/V_4\cong C_3$ 的三种角色、$S_4$ 的平凡与符号角色、$A_5$ 的平凡角色，沿商映射拉回，分别给出三、二、一种互异一维表示。因此上界恰好达到，唯一对应 $\widetilde E_6,\widetilde E_7,\widetilde E_8$。这只使用提升生成元的必要关系与商群下界，不需要未经证明的完整二元三角群呈示。相应权平方和为二十四、四十八、一百二十，与群阶再次一致。

最后比较旋转群。一个表示能下降到 $G/\{\pm I\}$，当且仅当 $-I$ 在其中作用为 $+1$。McKay图相邻顶点的中央符号相反；从平凡顶点出发，距离为偶数的顶点就是可下降的表示。因而 $A_4,S_4,A_5$ 的维数分别为
\[
 (1,1,1,3),\qquad (1,1,2,3,3),\qquad(1,3,3,4,5).
\]
二面体旋转群 $D_r$ 在 $r$ 为奇数时有两个一维表示和 $(r-1)/2$ 个二维表示；$r$ 为偶数时有四个一维表示和 $r/2-1$ 个二维表示。这与二元二面体图中偶中央符号的顶点相符。对于循环群，若其阶为偶数，下降的恰是中央符号为正的半数角色；若其阶为奇数，则它不含 $-I$，到旋转群的映射本来就是同构。
\end{solution}

\originalsection{Homework 11：循环箭图、根格与组成列}
本次作业对应旧讲义原题5.39--5.41，以下同样为AI辅助独立编写的解答。根格和扩张的定义只补到能够读懂并核算这些题的程度。

\begin{exercise}\label{hw:5.39}
\textbf{MIT 18.712（2010）Homework 11 原题5.39；旧PDF第96--97页。}
令 $Q_n$ 为 $n$ 个顶点和 $n$ 条箭头组成的有向圈。$Q_1$ 的不可分解表示由Jordan标准形分类。对 $Q_2$，即一对映射 $A:V\to W$、$B:W\to V$，求类似分类。
\begin{enumerate}[label=(\alph*)]
\item 对 $n\ge1$，考虑下列表示，并证明它们不可分解且两两不同构：
\begin{enumerate}[label=\arabic*)]
\item $E_{n,\lambda}$：$V=W=\C^n$，$A=J_n(\lambda)$ 是Jordan块，$B=I$，$\lambda\in\C$；
\item $E_{n,\infty}$：由 $E_{n,0}$ 交换 $V,W$ 和 $A,B$ 得到；
\item $H_n$：$V=\C^n,W=\C^{n-1}$，基为 $v_i,w_i$，且 $Av_i=w_i$、$Bw_i=v_{i+1}$（$i<n$），$Av_n=0$；
\item $K_n$：由 $H_n$ 交换 $V,W$ 和 $A,B$ 得到。
\end{enumerate}
\item 若 $AB$ 不幂零，证明表示含有一个直和项 $E_{n,\lambda}$，其中 $\lambda\ne0$。
\item 若 $AB$ 幂零，对 $V\oplus W$ 上的算子 $X(v,w)=(Bw,Av)$，证明 $X$ 幂零，并有由齐次Jordan链组成的基（每条链的首向量在 $V$ 或 $W$ 中）。据此证明(a)的列表完整。
\item 把分类推广到Kronecker箭图：两条箭头都从顶点一指向顶点二。
\item 把分类推广到 $n>2$ 的圈图任意定向。
\end{enumerate}
\end{exercise}
\begin{solution}
这里的同构保持两个顶点的名称，不能擅自交换 $V$ 和 $W$。这是 $E_{n,0}$ 与 $E_{n,\infty}$、$H_n$ 与 $K_n$ 需要分别列出的原因。

\textbf{(a)} 若 $B=I$，同构的两个换基矩阵必须相同，故 $E_{n,\lambda}$ 的同构问题就是 $J_n(\lambda)$ 的相似问题。其自同态代数为
\[
 \C[J_n(\lambda)]\cong\C[t]/(t^n).
\]
这个代数中的元素可逆当且仅当常数项非零；幂等元只能为零或一。因此表示不可分解。交换两顶点后，同样证明 $E_{n,\infty}$ 不可分解。

$H_n$ 在 $V\oplus W$ 上形成一条链
\[
 v_1\longmapsto w_1\longmapsto v_2\longmapsto\cdots
 \longmapsto w_{n-1}\longmapsto v_n\longmapsto0,
\]
其中箭头交替由 $A,B$ 给出。与 $X$ 交换的算子在单Jordan块上是 $X$ 的多项式；还须保持 $V,W$，所以只能含偶次项。其自同态代数为 $\C[X^2]$，仍是截断多项式代数，没有非平凡幂等元，因此 $H_n$ 不可分解。$K_n$ 同理。不同 $n$ 由维数区分；$H_n,K_n$ 的两个顶点维数不同；$E_{n,0}$ 的 $B$ 可逆而 $A$ 不可逆，$E_{n,\infty}$ 则反过来；其余 $E_{n,\lambda}$ 由 $BA$ 的特征值区分。这证明两两不同构，包括 $n=1$ 的两个顶点简单表示 $H_1,K_1$。

\textbf{(b)} 记 $T=BA$、$S=AB$，则 $AT=SA$、$BS=TB$。对非零特征值 $\lambda$，$A,B$ 分别在 $T,S$ 的广义 $\lambda$ 特征空间之间映射；由于 $T,S$ 在这些空间上可逆，$A,B$ 都是同构。各广义特征空间按顶点配对，给出箭图表示的直和分解。在其中一个非零特征值块上，用 $B$ 识别两空间后，表示写成 $(A,B)=(T,I)$。再对 $T$ 取Jordan标准形，各块给出 $E_{n,\lambda}$。若 $AB$ 不幂零，至少有一个这样的非零特征值块，故得到所求直和项。

\textbf{(c)} 若 $(AB)^r=0$，则 $(BA)^{r+1}=B(AB)^rA=0$，所以
\[
 X^2=\begin{pmatrix}BA&0\\0&AB\end{pmatrix}
\]
也幂零，从而 $X$ 幂零。还需说明Jordan基能保持顶点分次，而不能只引用不带分次的Jordan定理。

令 $\ell$ 是 $X$ 的最大链长。可取齐次向量 $u\in V$ 或 $W$，使 $X^{\ell-1}u\ne0$：把任一可用向量分成两顶点分量，至少一个分量仍满足条件。取齐次线性泛函 $\varphi$，使
\[
 \varphi(X^{\ell-1}u)=1,\qquad
 \varphi(X^ju)=0\quad(0\le j<\ell-1).
\]
这里“齐次”指它在与 $X^{\ell-1}u$ 不同的顶点分量上为零；同一分量上的条件可在线性无关的链向量上指定后延拓。定义
\[
 P(z)=\sum_{i=0}^{\ell-1}\varphi(X^{\ell-1-i}z)X^iu.
\]
直接计算得 $PX=XP$、$P(X^ju)=X^ju$，且 $P$ 保持两个顶点分量。因此 $P$ 是到这条链的表示投影，$\ker P$ 是齐次不变补空间。对补空间归纳，获得所需齐次链基。

奇数长链 $2n-1$ 的首向量在 $V$ 时给出 $H_n$，在 $W$ 时给出 $K_n$。偶数长链 $2n$ 的首向量在 $W$ 时，$B$ 可逆而 $A$ 为幂零Jordan块，给出 $E_{n,0}$；首向量在 $V$ 时给出 $E_{n,\infty}$。结合(b)，列表完整。

\textbf{(d)} 现在 $A,B$ 都是 $U\to W$ 的映射，同构是同一对源、靶换基作用于两矩阵。全部不可分解块如下：
\begin{itemize}
\item $P_d$，$d\ge0$：$\dim(U,W)=(d,d+1)$，取基 $u_1,\ldots,u_d$、$w_1,\ldots,w_{d+1}$，令 $Au_i=w_i$、$Bu_i=w_{i+1}$；
\item $I_d$，$d\ge0$：$\dim(U,W)=(d+1,d)$，令 $Au_i=w_i$（$i\le d$）、$Au_{d+1}=0$，$Bu_1=0$、$Bu_{i+1}=w_i$（$i\le d$）；
\item $R_{d,\lambda}$，$d\ge1,\lambda\in\C$：$U=W=\C^d$、$A=J_d(\lambda)$、$B=I$；
\item $R_{d,\infty}$，$d\ge1$：$U=W=\C^d$、$A=I$、$B=J_d(0)$。
\end{itemize}
$P_0=(0,\C)$ 与 $I_0=(\C,0)$ 都须保留。下面证明完备性，不把“矩阵束标准形”当作一个没有解释的口号。

第一步是维数不等的情形。对表示 $Y=(U,W;A,B)$，若 $AU+BU\ne W$，靶空间中的一个补空间立即给出 $P_0$ 的直和项；若 $\ker A\cap\ker B\ne0$，在源空间取补空间立即给出 $I_0$ 的直和项。因而不可分解表示除这两个简单块外，满足相应的满射或单射条件。

当 $AU+BU=W$ 时，定义反射
\[
 S^+Y=(T,U;p_1,p_2),\qquad
 T=\{(u_1,u_2)\in U\oplus U:Bu_1=Au_2\}.
\]
若 $\dim(U,W)=(a,b)$，则 $\dim S^+Y=(2a-b,a)$。反方向构造为
\[
 S^-Y=(W,C;\alpha,\beta),\qquad
 C=(W\oplus W)/\{(Au,Bu):u\in U\},
\]
其中 $\alpha(w)=[(0,-w)]$、$\beta(w)=[(w,0)]$。当 $\ker A\cap\ker B=0$ 时，其维数为 $(b,2b-a)$。

这两构造在上述子类中互为逆。事实上，$S^-S^+Y$ 的商空间由满射 $(u_1,u_2)\mapsto Bu_1-Au_2$ 识别为 $W$，其核恰为 $T$，且新箭头在这一识别下正好是 $A,B$。反过来，$S^+S^-Y$ 的拉回由单射 $u\mapsto(Au,Bu)$ 识别为 $U$，也恢复 $A,B$。同样的构造用于同态，因此它们给出互逆的表示范畴等价，保持直和、不可分解性和自同态代数。

若不可分解 $Y$ 有 $b>a$，记 $d=b-a>0$。反复用 $S^+$，维数从 $(a,b)$ 变为 $(a-d,a)$，差仍为 $d$，但总维数下降；必在出现负维数前达到例外简单块 $P_0$。所以 $d=1$，而原表示由 $P_0$ 反复用 $S^-$ 唯一恢复。代入上述基向量公式，$S^-P_j\cong P_{j+1}$，故它是某个 $P_j$。对偶地，$a>b$ 的不可分解表示恰为 $I_j$。这些块的自同态代数由反射保持，等于 $\C$，所以都不可分解。

第二步处理 $a=b$。先补一个只用于这里的扩张计算。对两个表示 $Y=(U,W;A,B)$、$Z=(U',W';A',B')$，顶点上的线性分裂使扩张的两箭头为块上三角矩阵，右上角是任意两个 $U\to W'$ 的映射；改变分裂所产生的差是
\[
 (s,t)\longmapsto(tA-A's,\ tB-B's),
\]
其中 $s:U\to U'$、$t:W\to W'$。因此该线性映射的核为 $\Hom(Y,Z)$，余核为 $\operatorname{Ext}^1(Y,Z)$，后者表示扩张在固定子模和商模后的等价类。由秩零化度公式，若 $\dim Y=(a,b)$、$\dim Z=(c,d)$，则
\[
 \dim\Hom(Y,Z)-\dim\operatorname{Ext}^1(Y,Z)=ac+bd-2ad.
\]

令 $R=(\C,\C;0,1)$。从 $P_j$ 的显式箭头公式计算同态到 $R$：靶空间的线性泛函必须在 $w_1,\ldots,w_j$ 上为零，只可自由指定 $w_{j+1}$；源空间的泛函随之确定。因此 $\dim\Hom(P_j,R)=1$。上式右侧也为一，故
\[
 \operatorname{Ext}^1(P_j,R)=0\qquad(j\ge0).
\]

设 $Y$ 平衡且不可分解，$A$ 不可逆。取 $0\ne u\in\ker A$。若 $Bu=0$，它给出 $I_0$ 直和项，与平衡的不可分解性矛盾；所以 $Bu\ne0$，由 $(u,Bu)$ 给出子表示 $R\subset Y$。令 $Z=Y/R$。若 $Z$ 有 $P_j$ 直和项，刚证的扩张消失使这一项可以提升并从 $Y$ 中分裂出来，矛盾。因此 $Z$ 没有 $P_j$ 项；它的维数仍平衡，而已分类的不平衡项中，只有 $I_j$ 具有正的 $\dim U-\dim W$。维数差的可加性于是也排除了 $I_j$ 项，$Z$ 的各不可分解项均平衡。

若 $Z$ 的某个不可分解项 $Z'$ 的箭头 $A'$ 可逆，则到 $R$ 的同态满足 $tA'=0$，从而 $t=0$，再由第二条箭头相容性得到 $s=0$。Euler式的右侧为零，故 $\operatorname{Ext}^1(Z',R)=0$，又使 $Z'$ 从 $Y$ 中分裂出来，矛盾。对顶点维数归纳，$Z$ 中各项的另一箭头 $B'$ 必可逆。扩张中 $B$ 的块对角线是 $1,B'$，故 $B$ 可逆。归纳初始维数一时结论直接成立。因此平衡不可分解表示至少有一个箭头可逆。

若 $B$ 可逆，将它化为 $I$ 后只剩 $A$ 的相似分类，不可分解性迫使它是单Jordan块 $J_d(\lambda)$。若 $A$ 可逆、$B$ 不可逆，将 $A$ 化为 $I$ 后，$B$ 必为单幂零Jordan块。得到恰好 $R_{d,\lambda},R_{d,\infty}$。它们的自同态代数是截断多项式代数，因此不可分解。维数差、箭头的可逆性及 $B^{-1}A$ 的Jordan参数区分列表中的各项。以上证明了完整分类；任意表示反复分解后，是这些块的直和。

\textbf{(e)} 对任意定向的 $n$ 边圈（$n>2$），分类可用整数直线上的有限区间和一个非零Jordan参数统一写出。把圈的顶点标为 $\Z/n\Z$，选择一次绕圈的正方向；在整数直线上，每条相邻边继承它投到圈上的箭头方向。

\textbf{区间块。} 对整数 $a\le b$，给每个整数 $\ell\in[a,b]$ 一个基向量 $u_\ell$。顶点 $i$ 的空间为所有 $\ell\equiv i\pmod n$ 的 $u_\ell$ 张成的空间。区间内相邻基向量按该边的实际方向以系数一相连；跨出区间的箭头作用为零。记此表示为 $G[a,b]$。同时把 $a,b$ 加上 $n$ 不改变表示，故可规定 $0\le a<n$。其维数向量可直接逐项数得
\[
 \dim G[a,b]_i=\#\{\ell\in[a,b]:\ell\equiv i\pmod n\}.
\]

\textbf{环块。} 对 $d\ge1$ 和 $\lambda\in\C^*$，在所有顶点放置 $\C^d$。沿圈运输基：顺箭头用该映射，逆箭头用逆映射。令绕圈一周的运输为 $J_d(\lambda)$；可把除一条固定边外的全部映射化为 $I$，最后一条若顺所选方向，取 $J_d(\lambda)$，若逆该方向，取 $J_d(\lambda)^{-1}$。记为 $B_{d,\lambda}$。这里必须 $\lambda\ne0$，否则所需逆映射不存在，零参数已经由区间块包含。

这两类块就是全部不可分解表示，参数除区间的共同平移外没有重复。下面先验证各块，再证明列表的完备性。

先验证各个实际块。一个区间块只有区间两端所对应的边可能不是同构；若两端落在同一条圈边，只有一条边可能非同构，其矩阵为幂零Jordan块，其余矩阵为 $I$。若两端落在不同边，把其余同构箭头逐一收缩，得到两个顶点和两条箭头。两箭头构成有向二圈时，收缩后的表示正是(a)--(c)中的某个 $H_j,K_j,E_{j,0},E_{j,\infty}$；两箭头平行时，正是(d)中的某个 $P_j,I_j$。收缩同构箭头只是用该同构识别两端空间，表示同态也必须按同一识别运输，所以收缩与反向展开保持自同态代数和直和分解。这证明每个区间块不可分解，包括某些顶点空间为零的情形（$0\to0$ 也是同构）。

若区间长度 $b-a+1$ 不是 $n$ 的倍数，维数向量中较大的那一段循环连续顶点确定端点残类，结合总维数确定 $[a,b]$。若长度为 $dn$，所有维数相同，但唯一非同构箭头的位置确定端点残类，幂零Jordan块的大小确定 $d$。因此不同区间参数确实不同构。环块所有箭头可逆，将它们运输到一个顶点后，同态就是与周游算子 $J_d(\lambda)$ 交换的矩阵，仍为截断多项式代数；所以环块不可分解。周游算子的相似类区分 $(d,\lambda)$，并使它们与至少有一个非同构箭头的区间块区分开。若改选反方向绕圈，环参数相应变为 $\lambda^{-1}$；固定正方向后没有这种歧义。

\textbf{完备性的第一步：展开为分次Kronecker表示。}
记原顶点空间为 $W_i$，边 $i$--$i+1$ 的映射为 $f_i$。在每条边中插入辅助空间 $U_i$，使两条辅助箭头都从 $U_i$ 指向相邻原顶点：
\[
 A_i:U_i\to W_i,\qquad B_i:U_i\to W_{i+1}.
\]
若原箭头为 $i\to i+1$，取 $U_i=W_i,A_i=I,B_i=f_i$；若原箭头为 $i+1\to i$，取 $U_i=W_{i+1},A_i=f_i,B_i=I$。每条原边因此指定了一条必须为同构的辅助箭头。

反之，只要这些指定箭头为同构，就能沿它们识别辅助空间与原顶点空间，收缩回原表示。表示同态也必须按同样方式识别，所以展开与收缩保持同态、直和及不可分解性。一张块对角映射可逆，当且仅当它的每个直和块可逆，故展开表示的任何直和分解中，各项仍满足指定同构条件。

令 $U=\bigoplus_iU_i,W=\bigoplus_iW_i$，合并所有 $A_i,B_i$ 为两张 $U\to W$ 的映射。这是(d)中的Kronecker表示，但保留 $\Z/n\Z$ 分次：
\[
 A(U_i)\subset W_i,\qquad B(U_i)\subset W_{i+1}.
\]
称 $A$ 的度为零，$B$ 的度为一；同态与分解都要求保持每一度。这里\textbf{不能}假定分次不可分解表示忘掉分次后仍不可分解。下面把(d)的论证在分次范畴中重新核对。

\textbf{第二步：不平衡的分次块。}
若 $AU+BU\ne W$，在各 $W_i$ 中分别取像的补空间，得到某个靶度上的一维简单直和项。若 $\ker A\cap\ker B\ne0$，取齐次核向量和齐次补空间，得到某个源度上的一维简单直和项。故除这些简单表示外，分次不可分解表示满足(d)所需的满射、单射条件。

正反射仍取拉回，但赋予分次
\[
 T_i=\{(u_i,u_{i+1})\in U_i\oplus U_{i+1}:Bu_i=Au_{i+1}\},
 \qquad S^+Y=(T,U;p_1,p_2).
\]
两投影分别为度零、度一。逆反射的商分次则必须平移为
\[
 C_i=(W_{i-1}\oplus W_i)/
       \{(Au,Bu):u\in U_{i-1}\},\qquad S^-Y=(W,C;\alpha,\beta),
\]
其中 $\alpha(w)=[(0,-w)]\in C_i$（$w\in W_i$），$\beta(w)=[(w,0)]\in C_{i+1}$。两箭头又为度零、度一。互逆识别是
\[
 (u_{i-1},u_i)\longmapsto Bu_{i-1}-Au_i\in W_i,
 \qquad u_i\longmapsto(Au_i,Bu_i)\in W_i\oplus W_{i+1},
\]
它们保持度，同态上的对应也保持度。因此(d)的反射等价现在确实保持\textbf{分次}不可分解性。

记 $a=\dim U,b=\dim W$。若 $b>a$，反复 $S^+$ 使总维数严格下降，而差 $b-a$ 保持不变；必须停在某个靶度上的简单表示，故差为一。逆反射唯一恢复一个齐次 $P_d$，其基可取
\[
 \deg u_j=\deg w_j=s+j-1,\qquad Au_j=w_j,\quad Bu_j=w_{j+1}.
\]
若 $a>b$，反复 $S^-$ 后停在某个源度上的简单表示，得到齐次 $I_d$，其基可取
\[
 \deg u_j=\deg w_j=s-(j-1),\qquad Au_j=w_j,\quad Bu_{j+1}=w_j,
\]
端点按(d)的公式取零。两者都在展开后二分圈的整数覆盖上形成一条有限链。由反射从简单块恢复的唯一性，没有遗漏新的分次参数。

\textbf{第三步：平衡块仍至少有一箭头可逆。}
先说明为什么(d)中的扩张分裂可保留分次。若两个分次表示的扩张忘掉分次后有表示截面 $s$，把 $s$ 分解为各个度的线性映射，取度零部分 $s_0$。由于 $A$ 度零、$B$ 度一，从相容关系中取相应度部分仍有 $s_0A=As_0,s_0B=Bs_0$；商映射 $q$ 度零，从 $qs=I$ 取度零部分又得 $qs_0=I$。所以该扩张也有保持分次的截面。

现设 $a=b=m>0$，$Y$ 分次不可分解而 $A$ 不可逆。取齐次 $0\ne u\in\ker A$。若 $Bu=0$，已得到源简单直和项，矛盾；因此 $u,Bu$ 构成分次子表示 $R=(\C,\C;0,1)$，源、靶的度相差一。令 $Z=Y/R$，其总维数仍平衡，且更小。

忘掉分次后，(d)已经证明 $\operatorname{Ext}^1(P_j,R)=0$。所以 $Z$ 中的任一分次 $P_j$ 直和项都能先提升为不分次的直和项，再用度零截面得到分次分裂，与 $Y$ 的分次不可分解性矛盾。$Z$ 没有靶多一维的项；结合第二步的不平衡分类与总维数差为零，也没有源多一维的 $I_j$ 项，故它的各个分次不可分解项均平衡。

若某一项 $Z'$ 的 $A'$ 可逆，(d)的计算给出忘分次后的 $\Hom(Z',R)=0$；总维数平衡使Euler式右侧为零，故 $\operatorname{Ext}^1(Z',R)=0$。取度零截面又使 $Z'$ 从 $Y$ 中分次分裂，矛盾。于是各 $A'$ 均不可逆。对 $m$ 归纳，它们的 $B'$ 均可逆；扩张中的 $B$ 有块对角线 $1,B'$，因而可逆。$m=1$ 时商为零，论证直接给出 $R$。因此平衡且分次不可分解的表示，至少有一个全空间箭头 $A,B$ 可逆。这一步使用扩张消失与度零截面，始终没有把分次不可分解偷换为不分次不可分解。

\textbf{第四步：齐次Jordan链与周游块。}
若 $B$ 可逆，用它识别两空间后只剩
\[
 T=B^{-1}A,\qquad T(U_i)\subset U_{i-1}.
\]
$T^n$ 为度零。取充分大的 $N$，其Fitting分解
\[
 U=\ker(T^{nN})\oplus\im(T^{nN})
\]
保持分次和 $T$，分别是幂零部分和可逆部分。具体说，取 $N$ 使 $\ker T^{2nN}=\ker T^{nN}$ 且像也稳定；若 $x=T^{nN}y$ 又在核中，则 $T^{2nN}y=0$，所以 $T^{nN}y=0$，即 $x=0$。核、像维数和由秩零化度公式等于总维数，因而确为直和；$T$ 在像上满射，故可逆。将各部分沿 $B$ 送到 $W$，得到分次表示的直和。

在幂零部分，对(c)的齐次Jordan投影作同一构造，但现在度取模 $n$。取最大链 $u,Tu,\ldots,T^{\ell-1}u$ 的齐次首向量，以及只在末向量所属度上非零的泛函 $\varphi$，使它在末向量上为一、其余链向量上为零。则
\[
 P(z)=\sum_{j=0}^{\ell-1}\varphi(T^{\ell-1-j}z)T^ju
\]
与 $T$ 交换、在该链上为恒等，并保持度：某项系数非零时，$z$ 与 $T^ju$ 的度相同。其分次核是补空间，归纳得到齐次Jordan链分解。分次不可分解性只允许一条链；给每个 $u_j$ 配上 $w_j=Bu_j$，箭头为 $Bu_j=w_j,Au_j=w_{j+1}$，末端为零，正是一条二分圈覆盖上的有限链。

在可逆部分，$T$ 在相邻各度之间为同构，所有度的维数相同。在固定度 $U_0$ 上对 $T^n$ 作Jordan分解，再把各块沿 $T$ 运送到全部度。绕一周时块由 $T^n$ 保持，故运输相容，给出分次直和。不可分解时恰有一个非零Jordan块，并且 $A,B$ 全可逆，得到周游块。

若 $B$ 不可逆，第三步已保证 $A$ 可逆，改用 $T=A^{-1}B$，其度为一。同样分为齐次幂零链与可逆周游块；可逆部分会同时使 $B$ 可逆，所以当前不可分解情形只能是幂零链。至此分次Kronecker的不可分解项全部得到。

\textbf{第五步：收缩回原圈。}
上述链沿展开后二分圈的整数覆盖前进或后退。每条指定同构箭头在每个直和项上仍为同构，所以链端不能在该边处截断：若端点是箭头的源，产生一个非零核基向量；若端点是靶，产生一个不在像中的基向量；链重复绕圈时其他基向量的像落在别的位置，不能补救。因而指定边上的基向量必须成对匹配。收缩这些匹配对后，一条有限链恰成为原圈整数覆盖上的有限区间 $[a,b]$，箭头方向继承原定向，系数为一。这正是 $G[a,b]$；链上所有非零系数都可通过逐个缩放基向量化为一，没有额外标量参数。

周游块收缩后各原箭头仍可逆。按已经固定的原圈正方向运输，绕一周为单Jordan块 $J_d(\lambda)$，$\lambda\ne0$，得到 $B_{d,\lambda}$。使用度为负一的 $T$ 时，先计算的一周方向可能相反；最后按固定正方向定义参数，必要时改成 $\lambda^{-1}$，即可与本题开头的环块定义一致。

任一原表示先展开，再反复分解为分次不可分解项，最后收缩；有限总维数保证分解过程终止。以上证明每项都在区间、环块列表中，而前面已逐块证明不可分解与参数不重复，故列表完备。这里所需的是全部不可分解表示的分类，无需另引一般直和分解唯一性定理。全同向时，区间块就是带顶点分次的幂零Jordan链；非零环块是各顶点维数相同的一周Jordan块，与(a)--(c)的二圈答案一致。
\end{solution}

\begin{exercise}\label{hw:5.40}
\textbf{MIT 18.712（2010）Homework 11 原题5.40；旧PDF第97页。}
令 $L\subset\tfrac12\Z^8$ 为满足下述条件的向量格：坐标或者全为整数，或者全为非整数的半整数；所有坐标之和为偶整数。
\begin{enumerate}[label=(\alph*)]
\item 设 $e_i$ 是标准正交基，令
\[
 \alpha_i=e_i-e_{i+1}\ (1\le i\le6),\quad
 \alpha_7=e_6+e_7,\quad
 \alpha_8=-\frac12\sum_{i=1}^8e_i.
\]
证明它们是 $L$ 的整数基。
\item 用通常内积证明 $L$ 中的根形成 $E_8$ 型根系；计算上述基的内积。
\item 证明 $E_7$ 和 $E_6$ 根格分别可由 $L$ 中前两个、前三个坐标相等的向量组成。
\item 分类并计数坐标形式，证明 $E_6,E_7,E_8$ 的根数分别为 $72,126,240$。
\end{enumerate}
\end{exercise}
\begin{solution}
这里的格是向量空间中由一组实线性无关向量的整数线性组合组成的群；本题中的根指平方长度为二的格向量。下面还核对反射性质，说明它们确实构成根系，而不只是一组长度相同的向量。

\textbf{(a)} 先令
\[
 D_r=\{(z_1,\ldots,z_r)\in\Z^r:\textstyle\sum_i z_i\text{为偶数}\}.
\]
向量 $e_i-e_{i+1}$（$1\le i<r$）和 $e_{r-1}+e_r$ 生成 $D_r$：相邻差生成所有 $e_i-e_r$，而 $(e_{r-1}+e_r)-(e_{r-1}-e_r)=2e_r$；将一个总和为偶数的向量减去适当的 $e_i-e_r$ 后，剩下的就是 $2e_r$ 的整数倍。这也证明它们是基，因为向量数为 $r$ 且线性无关。

取 $x\in L$，令 $t=-2x_8\in\Z$，并设 $y=x-t\alpha_8$。则 $y_8=0$。当 $x$ 的坐标为整数时 $t$ 为偶数；当坐标为非整数半整数时 $t$ 为奇数。因此所有 $y_i$ 都是整数，且 $\sum_i y_i=\sum_i x_i+4t$ 为偶数。前七个坐标遂在 $D_7$ 中，由 $\alpha_1,\ldots,\alpha_7$ 生成。加回 $t\alpha_8$，得到 $L$ 被所给八个向量生成。前七个向量在前七个坐标上线性无关，$\alpha_8$ 的第八坐标不为零，故八个向量线性无关，确为整数基。

\textbf{(b)} 所有 $\alpha_i$ 平方长度为二。不同基向量之间，非零内积恰为
\[
 (\alpha_i,\alpha_{i+1})=-1\ (1\le i\le5),\qquad
 (\alpha_5,\alpha_7)=-1,\qquad(\alpha_7,\alpha_8)=-1;
\]
其余为零，特别是 $(\alpha_6,\alpha_7)=0$。对应的图是三臂长度 $(4,1,2)$ 的 $E_8$：
\begin{center}
\begin{tikzpicture}[scale=.85,every node/.style={circle,draw,inner sep=2pt},every label/.style={font=\small}]
\foreach \i in {1,...,6}{\node (r\i) at (\i,0) {$\i$};}
\foreach \i [evaluate=\i as \j using int(\i+1)] in {1,...,5}{\draw (r\i)--(r\j);}
\node (r7) at (5,-.9) {$7$};\node (r8) at (5,-1.8) {$8$};\draw (r5)--(r7)--(r8);
\end{tikzpicture}
\end{center}
图上的数字是本题 $\alpha_i$ 的下标，与5.3中不带编号的图同构。

整数坐标且总和为偶数时，$\sum_i x_i^2\equiv\sum_i x_i\equiv0\pmod2$。半整数坐标可写为 $x_i=z_i+\tfrac12$，其中 $\sum_i z_i$ 为偶数，因而
\[
 (x,x)=\sum_i(z_i^2+z_i)+2
\]
也是偶整数。$L$ 对加减封闭，所以由极化公式知 $(x,y)\in\Z$。非零向量的长度平方因此至少为二。

设 $\Phi=\{\alpha\in L:(\alpha,\alpha)=2\}$。它有限，包含上述一组基，故张成 $\R^8$；同一直线上的根只有 $\pm\alpha$。对任一根，反射
\[
 s_\alpha(x)=x-(x,\alpha)\alpha
\]
保持 $L$，因为内积为整数；它又保持长度，所以置换 $\Phi$。且 $2(\beta,\alpha)/(\alpha,\alpha)=(\beta,\alpha)\in\Z$，满足晶体根系条件。因此 $\Phi$ 是约化晶体根系。

还要核对所给基是简单根系。若一个根的整数展开同时有正、负系数，写成 $\beta=\beta_+-\beta_-$，其中两部分是互不重叠的基向量正整数和。不同基向量内积非正，所以 $(\beta_+,\beta_-)\le0$；两个非零格向量的平方长度又都至少为二，于是
\[
 (\beta,\beta)=(\beta_+,\beta_+)+(\beta_-,\beta_-)-2(\beta_+,\beta_-)\ge4,
\]
与根长二矛盾。因此每个根的系数全非负或全非正，所给基确是简单根系。其Cartan图为 $E_8$，故根系类型为 $E_8$。

\textbf{(c)} 在 $x_1=x_2$ 的子空间中，取
\[
 e_3-e_4,\ e_4-e_5,\ e_5-e_6,\ e_6-e_7,\ e_7-e_8,\ e_7+e_8,
 \ -\frac12\sum_{i=1}^8e_i.
\]
前六个是最后六个坐标上 $D_6$ 的基。若共同坐标为 $a$，减去 $-2a$ 倍的最后一个向量，前两坐标化为零，剩下六个坐标为整数且总和为偶数。因此所写七个向量是这一子格的整数基。其内积图的三臂长度为 $(3,1,2)$，即 $E_7$。

在 $x_1=x_2=x_3$ 的子空间中，同样取
\[
 e_4-e_5,\ e_5-e_6,\ e_6-e_7,\ e_7-e_8,\ e_7+e_8,
 \ -\frac12\sum_{i=1}^8e_i.
\]
前五个生成最后五个坐标上的 $D_5$。减去 $-2a$ 倍最后一个向量后，前三坐标为零，最后五个坐标为整数且总和为偶数；总和的变化为 $-8a$，总是偶整数。所写六个向量遂为整数基，内积图的三臂长度为 $(2,1,2)$，即 $E_6$。这些子格中的根也构成根系：根的反射保持相应相等坐标子空间，前述根系检查原样适用；上述正、负系数矛盾也同样证明所写子格基是简单根系。

\textbf{(d)} 在 $E_8$ 中，整数根必须恰有两个非零坐标，各为 $\pm1$，因此为 $\pm e_i\pm e_j$，共 $4\binom82=112$ 个。半整数根的每个坐标绝对值至少为 $1/2$；长度平方为二迫使所有坐标都为 $\pm1/2$。总和为偶整数等价于负号个数为偶数，因此有 $2^7=128$ 个，合计 $240$。

对 $E_7$ 条件 $x_1=x_2$，整数根或者全部支撑在最后六个坐标上，共 $4\binom62=60$ 个；或者为 $\pm(e_1+e_2)$，另有两个。半整数根的前两符号相同；它们有两种选法，最后六个符号必须有偶数个负号，各有 $2^5$ 种。因此根数为
\[
 60+2+2\cdot2^5=126.
\]
对 $E_6$ 条件 $x_1=x_2=x_3$，整数根只能支撑在最后五个坐标上，共 $4\binom52=40$ 个。半整数根的前三符号全正时，后五个负号数须偶；前三全负时，后五个负号数须奇。两种情形各有 $2^4$ 种，故根数为
\[
 40+2\cdot2^4=72.
\]
两个计数都包含正根和负根，不能再乘二。
\end{solution}

\begin{exercise}\label{hw:5.41}
\textbf{MIT 18.712（2010）Homework 11 原题5.41；旧PDF第97页。}
设 $Q$ 为Dynkin箭图，$V_\alpha$ 是对应正根 $\alpha$ 的不可分解表示。简单根 $\alpha_i$ 对应的表示 $S_i=V_{\alpha_i}$ 只在顶点 $i$ 放置一维空间，其余顶点空间为零。
\begin{enumerate}[label=(\alph*)]
\item 若 $i$ 为源点（没有箭头进入），证明对任意表示 $V$，有 $\operatorname{Ext}^1(V,S_i)=0$；若 $i$ 为汇点（没有箭头离开），证明 $\operatorname{Ext}^1(S_i,V)=0$。
\item 对给定的箭头定向，构造 $V_\alpha$ 的Jordan--H\"older组成列。
\end{enumerate}
\end{exercise}
\begin{solution}
$\operatorname{Ext}^1(U,W)$ 的元素是短正合列 $0\to W\to E\to U\to0$ 的扩张类；它等于零表示所有这种列都能在表示范畴中分裂，而不只是向量空间能够分裂。

\textbf{(a)} 若 $i$ 是源点，考虑 $0\to S_i\to E\to V\to0$。除 $i$ 外，$E_j=V_j$；在 $i$ 取任一向量空间补空间 $L_i$，使 $E_i=(S_i)_i\oplus L_i$。从 $i$ 发出的所有箭头都在 $(S_i)_i$ 上为零，因为 $S_i$ 是子表示；从其他顶点进入 $i$ 的箭头则不存在。于是 $L_i$ 与其余全部 $E_j$ 一起构成一个子表示，与 $S_i$ 互补，故扩张分裂。

若 $i$ 是汇点，考虑 $0\to V\to E\to S_i\to0$。除 $i$ 外仍有 $E_j=V_j$，在 $i$ 取 $V_i$ 的一维补空间。由于没有离开 $i$ 的箭头，只在 $i$ 放置这个补空间就构成子表示，并同构于 $S_i$。所以扩张也分裂。这分别说明源点简单表示是内射对象、汇点简单表示是投射对象；本题并不需要一般投射覆盖或内射包理论。

\textbf{(b)} 底图是树，任意定向都无有向圈。先取一个汇点 $i_1$，删去它，再取剩余箭图的汇点 $i_2$，如此得到顺序 $i_1,\ldots,i_N$。等价地，每条箭头都从这个顺序中较后的顶点指向较前的顶点。

写 $\alpha=\sum_i n_i\alpha_i$，题设中的根对应意味着 $\dim(V_\alpha)_i=n_i$。定义 $U_r\subset V_\alpha$，使它在已选顶点 $i_1,\ldots,i_r$ 上等于全部顶点空间，在其余顶点为零。已选顶点不会向未选顶点发出箭头，因此各 $U_r$ 都是子表示，并有
\[
 0=U_0\subset U_1\subset\cdots\subset U_N=V_\alpha,
 \qquad U_r/U_{r-1}\cong S_{i_r}^{\oplus n_{i_r}}.
\]
在每个 $(V_\alpha)_{i_r}$ 中取任意完整旗标
$0\subset L_1\subset\cdots\subset L_{n_{i_r}}=(V_\alpha)_{i_r}$，用它把 $U_{r-1}\subset U_r$ 细分。细分时，流向更早顶点的箭头落在已经完整加入的空间中，进入当前顶点的未选空间为零，所以每一步仍为子表示。相邻商依次为
\[
 \underbrace{S_{i_1},\ldots,S_{i_1}}_{n_{i_1}\text{个}},\quad
 \underbrace{S_{i_2},\ldots,S_{i_2}}_{n_{i_2}\text{个}},\quad\ldots,\quad
 \underbrace{S_{i_N},\ldots,S_{i_N}}_{n_{i_N}\text{个}}.
\]
零重数的段略去。这就是对该定向的一条组成列，长度为 $\sum_i n_i$。该构造适用于任意无圈箭图表示，不依赖 $V_\alpha$ 不可分解；不可分解性也不意味着组成列只有一个商项。
\end{solution}

% 中文题面译自 MIT 18.712 Fall 2010 takehome assignment 第1页。
% 未见随题公开官方答案；下列解答为编者独立推导。
\originalsection{期末作业第1题}
\begin{exercise}\label{final:1}
\textbf{MIT 18.712（2010）期末作业原题1；官方期末PDF第1页.}
设 $Q$ 是有限箭图，即有限有向图。域 $k$ 上的路径代数 $A(Q)$ 以所有沿箭头方向的路径为一组基，包括各顶点上长度为零的路径；乘法为路径拼接，不能拼接时乘积为零。
\begin{enumerate}[label=(\roman*)]
\item 把上三角矩阵代数表示为某个 $A(Q)$。
\item 证明：$A(Q)$ 有限维，当且仅当 $Q$ 无有向圈。
\item 当 $Q$ 无有向圈时，把左正则模 $A(Q)$ 分解为不可分解模的直和。
\item 给出使 $A(Q)\cong A(Q)^{\mathrm{op}}$ 成立的箭图条件，并据此举出一个与其反代数不同构的代数。
\end{enumerate}
\end{exercise}
\begin{solution}
以下约定乘积 $pq$ 表示先走 $q$、再走 $p$，与线性映射的复合次序一致。改变这一约定只会把所画箭图反向，不改变题目的内容。记顶点 $i$ 的零长路径为 $e_i$，则
\[
 e_i^2=e_i,\qquad e_i e_j=0\ (i\ne j),\qquad 1=\sum_i e_i.
\]

\textbf{(i)} 对 $n\times n$ 上三角矩阵，取有向直线
\[
1\longleftarrow2\longleftarrow\cdots\longleftarrow n.
\]
每对 $i\le j$ 恰有一条从 $j$ 到 $i$ 的路径 $p_{ij}$。令 $p_{ij}\mapsto E_{ij}$，其中 $E_{ij}$ 是矩阵单位。路径乘积满足
\[
 p_{ij}p_{ab}=\begin{cases}p_{ib},&j=a,\\0,&j\ne a,\end{cases}
\]
这正是 $E_{ij}E_{ab}=\delta_{ja}E_{ib}$。该映射把一组基映为一组基，因此是代数同构。

\textbf{(ii)} 若存在正长有向圈 $c$，则 $c,c^2,c^3,\ldots$ 是互不相同的路径，故在路径基中线性无关，$A(Q)$ 无限维。反之，若没有有向圈，一条有向路径不能重复顶点，否则重复部分给出有向圈。若 $Q$ 有 $s$ 个顶点，所有路径长度至多 $s-1$；箭头数有限，所以这样的路径只有有限条，$A(Q)$ 有限维。

\textbf{(iii)} 由正交幂等元得到
\[
 A(Q)=\bigoplus_{i\in Q_0}A(Q)e_i.
\]
这里 $A(Q)e_i$ 的路径基是所有以 $i$ 为起点的路径。直和性可直接由路径的起点唯一看出。任一左模自同态 $T$ 由 $T(e_i)$ 确定，因为 $T(ae_i)=aT(e_i)$；又有 $T(e_i)=e_iT(e_i)$，故
\[
 \End_{A(Q)}(A(Q)e_i)\cong(e_iA(Q)e_i)^{\mathrm{op}}.
\]
无有向圈时，从 $i$ 回到 $i$ 的路径只有 $e_i$，所以右侧是 $k$。若 $A(Q)e_i$ 有非平凡直和分解，投影到一个直和项将给出一个不等于 $0,1$ 的幂等自同态，而域 $k$ 没有这种幂等元。因此各项均不可分解，所写直和就是所求分解。其各项维数等于以相应顶点为起点的路径数。

\textbf{(iv)} 记 $Q^{\mathrm{op}}$ 为把所有箭头反向的箭图。反向阅读路径给出
\[
 A(Q)^{\mathrm{op}}\cong A(Q^{\mathrm{op}}).
\]
因此，一个适用于任意有限箭图的充分条件是 $Q\cong Q^{\mathrm{op}}$：存在顶点置换，把从 $i$ 到 $j$ 的箭头与从置换后的 $j$ 到置换后的 $i$ 的箭头一一对应。

对无有向圈的箭图，这一条件也是必要的。为证明这一点，令 $J$ 为所有正长路径张成的理想。它是幂零理想，且 $A(Q)/J\cong k^{Q_0}$，故 $J$ 是代数的根基。具体说，幂零理想在任意简单模上的像必须为零，否则迭代作用仍等于整个简单模，与幂零性矛盾；另一方面，$k^{Q_0}$ 的简单模足以分离其元素。因此根基恰为 $J$。代数同构保持根基，并置换商代数 $k^{Q_0}$ 的坐标幂等元。空间
\[
 \overline e_j(J/J^2)\overline e_i
\]
的维数恰为从 $i$ 到 $j$ 的箭头数，因为 $J^2$ 由长度至少二的路径张成。于是代数同构也恢复箭头数；$A(Q)\cong A(Q)^{\mathrm{op}}$ 将迫使 $Q\cong Q^{\mathrm{op}}$。

例如取下面的分叉箭图：
\begin{center}
\begin{tikzpicture}[>=Stealth,every node/.style={circle,draw,inner sep=2pt},scale=.85]
\node (a) at (0,0) {$1$};\node (b) at (1.4,.65) {$2$};\node (c) at (1.4,-.65) {$3$};
\draw[->] (a)--(b);\draw[->] (a)--(c);
\node[draw=none] at (3,0) {$Q^{\mathrm{op}}:$};
\node (aa) at (4.2,0) {$1$};\node (bb) at (5.6,.65) {$2$};\node (cc) at (5.6,-.65) {$3$};
\draw[->] (bb)--(aa);\draw[->] (cc)--(aa);
\end{tikzpicture}
\end{center}
左图有一个出度为二的顶点，右图没有这样的顶点，故两图不同构，因而 $A(Q)\not\cong A(Q)^{\mathrm{op}}$。这也是一个五维代数的例子：三条零长路径加两条箭头构成其基。
\end{solution}


\originalsection{期末作业第2题}
\begin{exercise}\label{final:2}
\textbf{MIT 18.712 (2010) 期末作业第2题.}
设 $\mathbb F_q$ 为有 $q$ 个元素的有限域. 分类二维空间 $\mathbb F_q^2$ 的仿射变换群
$\GL_2(\mathbb F_q)\ltimes\mathbb F_q^2$ 的所有不可约复表示, 并求出它们的特征标.
\end{exercise}

以下解答由中文编者撰写. 原题见课程期末作业第1页; 所用理论见公开旧讲义第4.24节与第4.26节, 印刷页68--75及76--77. 原讲义的第4.24节先假设特征不为2, 本题本身没有这个限制. 下面用二次扩域的乘法模型重新证明所需结论, 因而也包含偶特征和 $q=2$.

\begin{solution}
令 $F=\mathbb F_q$, $K=\GL_2(F)$, $A=F^2$, 把 $A$ 的元素写成列向量. 记仿射群为 $\Gamma$, 其元素写成 $(g,v)$, 对列向量 $x$ 的作用是 $x\mapsto gx+v$. 因而
\[
 (g,v)(h,w)=(gh,v+gw),\qquad
 |K|=(q^2-1)(q^2-q)=q(q-1)(q^2-1),\qquad |\Gamma|=q^2|K|.
\]
先完整列出 $K$ 的不可约表示, 再分析非零平移特征对应的稳定子. 这样得到的参数和公式可以在任意给定 $(g,v)$ 上直接求值.

\textbf{一、$K$ 的参数与共轭类型.}
写 $\widehat{F^\times}$ 为乘法群 $F^\times$ 的复特征群. 再取二次扩域 $E=\mathbb F_{q^2}$, 把 $E$ 看成二维 $F$ 向量空间. 对 $z\in E^\times$, 乘法 $x\mapsto zx$ 给出一个矩阵 $m_z\in K$. 它的迹和行列式分别是
\[
 \operatorname{Tr}_{E/F}(z)=z+z^q,\qquad N_{E/F}(z)=z^{q+1}.
\]
这个构造不要求选取平方根, 所以适用于偶特征.

$K$ 有如下四种共轭类型:
\[
 S_a=aI,\qquad
 J_a=\begin{pmatrix}a&1\\0&a\end{pmatrix},\qquad
 D_{a,b}=\begin{pmatrix}a&0\\0&b\end{pmatrix},\qquad E_z=m_z.
\]
其中 $a,b\in F^\times$, $D_{a,b}$ 中 $a\ne b$, 无序对 $\{a,b\}$ 只计一次; $E_z$ 中 $z\in E^\times\setminus F^\times$, 对 $\{z,z^q\}$ 也只计一次. 这四种情形分别对应标量矩阵、非平凡Jordan块、两个不同的域内特征值、不可约二次特征多项式. 有理标准形说明没有遗漏, 并给出上述参数的唯一性.

各类大小及类数如下:
\begin{center}
\begin{tabular}{@{}lll@{}}
\toprule
类型 & 每个共轭类的大小 & 共轭类数 \\
\midrule
$S_a$ & $1$ & $q-1$ \\
$J_a$ & $q^2-1$ & $q-1$ \\
$D_{a,b}$ & $q(q+1)$ & $(q-1)(q-2)/2$ \\
$E_z$ & $q(q-1)$ & $q(q-1)/2$ \\
\bottomrule
\end{tabular}
\end{center}
确实, 四种中心化子的阶依次为 $|K|$, $q(q-1)$, $(q-1)^2$, $q^2-1$. 对最后一种, 与 $m_z$ 交换的 $F$ 线性变换恰是 $E$ 线性变换, 因而可逆者恰为 $E^\times$. 类数合计为 $q^2-1$.

\textbf{二、$K$ 的全部不可约特征标.}
它们分成以下四族. 此处 $\alpha,\beta\in\widehat{F^\times}$, 而 $\theta\in\widehat{E^\times}$. 记 $\theta^q(z)=\theta(z^q)$.
\begin{center}
\begin{tabular}{@{}llll@{}}
\toprule
表示 & 参数 & 维数 & 个数 \\
\midrule
$L_\alpha$ & $\alpha$ & $1$ & $q-1$ \\
$\mathrm{St}_\alpha$ & $\alpha$ & $q$ & $q-1$ \\
$P_{\alpha,\beta}$ & $\{\alpha,\beta\}$, $\alpha\ne\beta$ & $q+1$ & $(q-1)(q-2)/2$ \\
$C_\theta$ & $\{\theta,\theta^q\}$, $\theta\ne\theta^q$ & $q-1$ & $q(q-1)/2$ \\
\bottomrule
\end{tabular}
\end{center}
除表中写出的无序参数等同以外, 没有其他同构关系. 对非分裂参数, 条件 $\theta\ne\theta^q$ 必须保留.

为了给出每个共轭类上的值, 令 $\varphi=\theta|_{F^\times}$. 特征标表拆成两部分写为
\begin{center}
\begin{tabular}{@{}lll@{}}
\toprule
特征标 & $S_a$ & $J_a$ \\
\midrule
$\chi_{L_\alpha}$ & $\alpha(a^2)$ & $\alpha(a^2)$ \\
$\chi_{\mathrm{St}_\alpha}$ & $q\alpha(a^2)$ & $0$ \\
$\chi_{P_{\alpha,\beta}}$ & $(q+1)\alpha(a)\beta(a)$ & $\alpha(a)\beta(a)$ \\
$\chi_{C_\theta}$ & $(q-1)\varphi(a)$ & $-\varphi(a)$ \\
\bottomrule
\end{tabular}
\end{center}
\begin{center}
\begin{tabular}{@{}lll@{}}
\toprule
特征标 & $D_{a,b}$ & $E_z$ \\
\midrule
$\chi_{L_\alpha}$ & $\alpha(ab)$ & $\alpha(N_{E/F}z)$ \\
$\chi_{\mathrm{St}_\alpha}$ & $\alpha(ab)$ & $-\alpha(N_{E/F}z)$ \\
$\chi_{P_{\alpha,\beta}}$ & $\alpha(a)\beta(b)+\alpha(b)\beta(a)$ & $0$ \\
$\chi_{C_\theta}$ & $0$ & $-\theta(z)-\theta(z^q)$ \\
\bottomrule
\end{tabular}
\end{center}
这些参数也可以明确用整数给出. 取 $E^\times$ 的生成元 $\gamma$, 则 $\gamma^{q+1}$ 生成 $F^\times$. 可令
\[
 \alpha_t(\gamma^{q+1})=e^{2\pi it/(q-1)}\quad(0\le t<q-1),\qquad
 \theta_j(\gamma)=e^{2\pi ij/(q^2-1)}\quad(0\le j<q^2-1).
\]
$C_\theta$ 的参数就是把不被 $q+1$ 整除的 $j$ 按 $j\sim qj\pmod{q^2-1}$ 成对归类. $\theta^q=\theta$ 的参数共有 $q-1$ 个, 所以余下的参数对共有 $q(q-1)/2$ 个. 当 $q=2$ 时 $\widehat{F^\times}$ 只有平凡特征, 上述 $\alpha_0$ 公式的分母为1, 仍有意义.

下面证明表中表示的存在、不可约性和完整性.

一维表示定义为 $L_\alpha(g)=\alpha(\det g)$. 取上三角子群 $B\le K$, 定义
\[
 \xi_{\alpha,\beta}\begin{pmatrix}a&c\\0&b\end{pmatrix}=\alpha(a)\beta(b),\qquad
 P_{\alpha,\beta}=\Ind_B^K\xi_{\alpha,\beta}.
\]
这里暂时也允许 $\alpha=\beta$. $K/B$ 是 $F^2$ 中直线的集合, 共 $q+1$ 条. 诱导特征标只对 $g$ 的不变直线求和, 每条直线贡献 $g$ 在该直线上和商空间上的两个标量所对应的 $\alpha,\beta$ 值. 四种类型的不变直线数分别为 $q+1,1,2,0$, 于是得到 $P_{\alpha,\beta}$ 那一行, 包括 $\alpha=\beta$ 的情形.

为核不可约性, 令 $w=\left(\begin{smallmatrix}0&1\\1&0\end{smallmatrix}\right)$. 上三角群的双陪集分解是 $K=B\sqcup BwB$, 且 $B\cap wBw^{-1}$ 为对角群. 用本书前文的Frobenius互反与Mackey双陪集分解, 有
\begin{equation}\label{eq:final2-principal-inner}
 \langle\chi_{P_{\alpha,\beta}},\chi_{P_{\mu,\nu}}\rangle_K
 =\delta_{\alpha,\mu}\delta_{\beta,\nu}
  +\delta_{\alpha,\nu}\delta_{\beta,\mu}.
\end{equation}
具体地说, 双陪集 $B$ 对应在 $B$ 上直接比较两个一维特征, 给第一项; $BwB$ 对应在对角群上比较 $\alpha(a)\beta(b)$ 与 $\mu(b)\nu(a)$, 给第二项. 对角线上的 $a,b$ 可以独立取值, 所以只有表中两个匹配条件贡献1. 这个论证不使用特征不为2的条件.

当 $\alpha\ne\beta$ 时, 式\eqref{eq:final2-principal-inner}的自内积为1, 故 $P_{\alpha,\beta}$ 不可约, 而且只在交换参数时同构. 当 $\alpha=\beta$ 时, Frobenius互反说明 $L_\alpha$ 在 $P_{\alpha,\alpha}$ 中恰出现一次. 自内积为2, 所以
\[
 P_{\alpha,\alpha}=L_\alpha\oplus\mathrm{St}_\alpha,
 \qquad \dim\mathrm{St}_\alpha=q,
\]
其中 $\mathrm{St}_\alpha$ 不可约. 用 $P_{\alpha,\alpha}$ 的特征标减去 $L_\alpha$ 即得表中 $\mathrm{St}_\alpha$ 的行. 同一内积公式和Frobenius互反还给出
$\langle\chi_{\mathrm{St}_\alpha},\chi_{\mathrm{St}_\mu}\rangle_K=\delta_{\alpha,\mu}$, 并说明它们与 $L_\mu$、参数不同的主系列都正交. 因而这些表示之间没有遗漏的同构.

剩下的 $C_\theta$ 用虚表示构造. 虚表示的特征标是实际表示特征标的整数线性组合; 若其自内积为1, 则在不可约特征标基下只有一个系数为 $1$ 或 $-1$. 若它在单位元的值为正, 系数必为 $1$, 因而它确实是一个不可约表示的特征标.

令 $T=\{m_z:z\in E^\times\}\le K$, 并取 $\theta\ne\theta^q$. 诱导特征标公式给出 $Y_\theta=\Ind_T^K\theta$ 的值:
\[
 \chi_{Y_\theta}(S_a)=q(q-1)\varphi(a),\qquad
 \chi_{Y_\theta}(J_a)=\chi_{Y_\theta}(D_{a,b})=0,\qquad
 \chi_{Y_\theta}(E_z)=\theta(z)+\theta(z^q).
\]
第一式来自中心标量及指数 $[K:T]=q(q-1)$; 后一式来自 $E_z$ 的共轭类与 $T$ 相交于 $m_z,m_{z^q}$, 且其中心化子就是 $T$. 中间两种类型不能共轭进入 $T$.

考虑虚表示
\begin{equation}\label{eq:final2-cuspidal-virtual}
 \mathcal C_\theta
 =\mathrm{St}_1\otimes P_{\varphi,1}-P_{\varphi,1}-Y_\theta.
\end{equation}
即使 $\varphi=1$, 这里的 $P_{\varphi,1}$ 仍是已经定义的实际诱导表示, 无须假定它不可约. 用已算出的三行特征标相乘、相减, 就得到表中 $\chi_{C_\theta}$ 的四个值. 尤其单位元的值为 $q-1>0$.

还要真正核其自内积. 对另一满足 $\kappa\ne\kappa^q$ 的参数 $\kappa$, 记 $\omega=\kappa|_{F^\times}$, 并设
$u=\delta_{\theta,\kappa}+\delta_{\theta,\kappa^q}$. 循环群特征的正交性给出
\[
 \sum_{z\in E^\times}
 (\theta(z)+\theta^q(z))\overline{(\kappa(z)+\kappa^q(z))}
 =2(q^2-1)u.
\]
在子群 $F^\times$ 上, 这个被加数变为 $4\varphi(a)\overline{\omega(a)}$. 结合四种类的大小, 标量类与Jordan类的总贡献为
\[
 \bigl((q-1)^2+q^2-1\bigr)
 \sum_{a\in F^\times}\varphi(a)\overline{\omega(a)}
 =2q(q-1)\sum_{a\in F^\times}\varphi(a)\overline{\omega(a)}.
\]
非分裂类的贡献则为
\[
 \frac{q(q-1)}2\left(
 2(q^2-1)u-4\sum_{a\in F^\times}\varphi(a)\overline{\omega(a)}\right).
\]
两式中的域内求和抵消, 再除以 $|K|=q(q-1)(q^2-1)$, 得到
\begin{equation}\label{eq:final2-cuspidal-inner}
 \langle\chi_{\mathcal C_\theta},\chi_{\mathcal C_\kappa}\rangle_K
 =\delta_{\theta,\kappa}+\delta_{\theta,\kappa^q}.
\end{equation}
所以式\eqref{eq:final2-cuspidal-virtual}确实给出不可约 $C_\theta$, 且参数正好按Frobenius对归类. 当 $q>2$ 时, 维数 $q-1$ 与其余三族的维数不同; 当 $q=2$ 时, $C_\theta$ 在 $J_1$ 上的值为 $-1$, 而唯一的 $L_1$ 在这里的值为1, 所以也不与它同构.

全部表示的维数平方和为
\begin{align*}
 &(q-1)+(q-1)q^2
 +\frac{(q-1)(q-2)}2(q+1)^2
 +\frac{q(q-1)}2(q-1)^2\\
 &\hspace{3em}=q(q-1)^2(q+1)=|K|.
\end{align*}
已构造的表示两两不同, 所以正则表示的维数平方和定理说明它们构成 $K$ 的完整不可约列表. 这里没有使用“所有一维表示都是行列式特征”作为前提; 对 $q=2$ 而言, 这句话本来就是错误的, 因为 $C_\theta$ 也是一维表示.

\textbf{三、平移特征的两个轨道及稳定子.}
选定一个非平凡加法特征 $\psi:F\to\C^\times$. 若 $q=p^f$, 可取
$\psi(t)=\exp(2\pi i\operatorname{Tr}_{F/\mathbb F_p}(t)/p)$.
对行向量 $y\in F^{1\times2}$, 定义
\[
 \lambda_y(v)=\psi(yv),\qquad v\in A.
\]
这些是 $A$ 的全部 $q^2$ 个特征. 若 $y\ne y'$, 则非零 $F$ 线性映射 $v\mapsto(y-y')v$ 的像为 $F$, 因而 $\lambda_y\ne\lambda_{y'}$. $K$ 在这些特征上的作用为 $g\cdot\lambda_y=\lambda_{yg^{-1}}$. 非零行向量彼此可由可逆矩阵变换, 所以只有两个轨道: $\{0\}$ 与所有非零行向量.

零轨道的稳定子是整个 $K$. 它给出上面四族表示向 $\Gamma$ 的提升, 即
\[
 \widetilde U(g,v)=U(g),\qquad U\in\Irr(K),\qquad
 \chi_{\widetilde U}(g,v)=\chi_U(g).
\]
因此平移子群在这些表示上作用平凡, 维数也不变.

取非零轨道的代表 $y_0=(1,0)$. 稳定子为
\[
 H=\left\{h(b,d)=\begin{pmatrix}1&0\\b&d\end{pmatrix}:
 b\in F,\ d\in F^\times\right\},\qquad |H|=q(q-1).
\]
因为 $y_0h=y_0$, 这也等价于 $y_0h^{-1}=y_0$. 其乘法为
$h(b,d)h(c,e)=h(b+dc,de)$, 即一维仿射群.

$H$ 的不可约表示也可完整算出. 正规子群
$N=\{h(b,1):b\in F\}$ 是加法群, 而 $F^\times$ 对 $\widehat N$ 的作用有两个轨道: 平凡特征与全部非平凡特征. 平凡轨道给出
\[
 \eta_\tau(h(b,d))=\tau(d),\qquad \tau\in\widehat{F^\times}.
\]
非平凡轨道的稳定子为平凡群, 给出唯一的另一个表示
\[
 R=\Ind_N^H\psi,\qquad \dim R=q-1.
\]
这里 $\psi$ 在 $N$ 上的含义是 $\psi(h(b,1))=\psi(b)$. 其特征标为
\begin{equation}\label{eq:final2-H-character}
 r(b,d)=\chi_R(h(b,d))=
 \begin{cases}
 q-1,& d=1,\ b=0,\\
 -1,& d=1,\ b\ne0,\\
 0,& d\ne1.
 \end{cases}
\end{equation}
确实, 当 $d\ne1$ 时, $h(b,d)$ 不能共轭进 $N$, 诱导特征标为0. 当 $d=1$ 时, 诱导公式给出
$\sum_{t\in F^\times}\psi(tb)$, 它在 $b=0$ 时为 $q-1$, 在 $b\ne0$ 时为 $-1$.
$R$ 在 $N$ 上的不同特征各出现一次, 且 $F^\times$ 传递置换这些一维空间, 所以 $R$ 不可约. $\eta_\tau$ 两两不同, 又与 $R$ 不同; 即使 $q=2$, $R$ 在非平凡平移上的值为 $-1$, 仍与 $\eta_1$ 有别. 最后
\[
 (q-1)\cdot1^2+(q-1)^2=q(q-1)=|H|,
\]
证明没有其他不可约表示.

\textbf{四、$\Gamma$ 的完整分类与其不可约性.}
因为 $H$ 固定 $\lambda_{y_0}$, 对任意 $\sigma\in\Irr(H)$, 公式
\[
 \rho_\sigma(h,v)=\lambda_{y_0}(v)\sigma(h)
 \qquad (h\in H,\ v\in A)
\]
定义 $H\ltimes A$ 的表示. 乘法成立的原因是
$\lambda_{y_0}(hw)=\lambda_{y_0}(w)$. 定义
\[
 V_\sigma=\Ind_{H\ltimes A}^{\Gamma}\rho_\sigma.
\]
把 $\sigma=\eta_\tau$ 的表示记为 $Q_\tau$, 把 $\sigma=R$ 的表示记为 $Q_R$. 则所有不可约表示为
\begin{equation}\label{eq:final2-all-affine}
 \widetilde L_\alpha,\quad
 \widetilde{\mathrm{St}}_\alpha,\quad
 \widetilde P_{\alpha,\beta},\quad
 \widetilde C_\theta,\quad
 Q_\tau\ (\tau\in\widehat{F^\times}),\quad Q_R.
\end{equation}
前四族的参数等同如上; 后两族没有进一步的同构关系. 新表示的维数为
\[
 \dim Q_\tau=q^2-1,\qquad
 \dim Q_R=(q^2-1)(q-1).
\]

为证明这一分类, 将任意复表示按平移特征分解为权空间
$V=\bigoplus_y V_y$, 其中 $v\in A$ 在 $V_y$ 上作用为 $\psi(yv)$. 群 $K$ 按 $y\mapsto yg^{-1}$ 置换这些空间. 若 $V$ 不可约, 它的非零权空间只能落在一个轨道: 否则每个轨道对应的权空间和都是真不变子空间. 零轨道的情况恰是 $K$ 的提升表示.

对于非零轨道, $V_{y_0}$ 是 $H$ 的表示. 若它有真非零 $H$ 不变子空间, 将该子空间由 $K$ 搬到各个权空间, 其直和就给出 $V$ 的真非零 $\Gamma$ 不变子空间. 因而 $V_{y_0}$ 必须是不可约 $\sigma$. 反过来, $V_\sigma$ 在每个非零权空间上的维数为 $\dim\sigma$, $y_0$ 权空间上的 $H$ 作用正是 $\sigma$. 任何 $\Gamma$ 不变子空间也分解为权空间; 在 $y_0$ 权空间上由 $\sigma$ 的不可约性只能取零或全部, 再由轨道传递性在其余权空间上同时取零或全部. 所以 $V_\sigma$ 不可约.

取 $K/H$ 的一组代表, 把 $V_{y_0}$ 搬到全部非零权空间, 自然得到 $V_\sigma$ 到 $V$ 的同构. 这也证明每个非零轨道表示都由这种诱导构造获得. 其 $y_0$ 权空间能恢复 $\sigma$, 因而不同 $\sigma$ 给出不同表示. 零轨道与非零轨道则可由平移子群的作用区分. 这是本题所需半直积分类定理的直接证明.

再作完整性检验. 前四族的平方和为 $|K|$, 后两族的平方和为
\[
 (q^2-1)^2\bigl((q-1)+(q-1)^2\bigr)
 =(q^2-1)^2q(q-1)=(q^2-1)|K|.
\]
合计正好为 $q^2|K|=|\Gamma|$. 不可约表示总数为
$(q^2-1)+(q-1)+1=q^2+q-1$.

\textbf{五、所有仿射元素上的特征标公式.}
前四族的公式已经给出: 忽略 $v$, 在 $g$ 的共轭类型上代入 $K$ 的特征标表. 为求后两族, 定义
\[
 W_g=\{y\in F^{1\times2}:yg=y\},\qquad
 k_g=\dim_F W_g,\qquad
 \epsilon_g(v)=
 \begin{cases}1,&v\in\im(g-I),\\0,&v\notin\im(g-I).\end{cases}
\]
这里 $\im(g-I)$ 是 $F^2$ 中的列向量像空间. 对任意 $\sigma\in\Irr(H)$, 诱导特征标为以下完全可计算的有限和:
\begin{equation}\label{eq:final2-affine-induced-character}
 \chi_{V_\sigma}(g,v)=\frac1{|H|}
 \sum_{\substack{t\in K\\tgt^{-1}\in H}}
 \psi(y_0tv)\chi_\sigma(tgt^{-1}).
\end{equation}
例如可遍历 $K$ 的全部矩阵, 对满足条件者使用式\eqref{eq:final2-H-character}或 $\eta_\tau$ 直接求和. 下面化成更简短的闭式.

先说明式\eqref{eq:final2-affine-induced-character}的来源与归一化. 对 $(s,u)\in\Gamma$, 有
\[
 (s,u)^{-1}(g,v)(s,u)
 =\bigl(s^{-1}gs,\ s^{-1}(v+(g-I)u)\bigr).
\]
若 $s^{-1}gs\in H$, 则行向量 $y=y_0s^{-1}$ 满足 $yg=y$, 所以 $y(g-I)u=0$. 诱导公式中的平移特征因而与 $u$ 无关, 对 $u$ 求和抵消分母里的 $|A|$. 再令 $t=s^{-1}$, 便得到式\eqref{eq:final2-affine-induced-character}.

对每个满足 $yg=y$ 的非零行向量 $y$, 恰有 $|H|$ 个 $t$ 满足 $y_0t=y$. 它们彼此相差左乘 $H$ 元素, 故 $\chi_\sigma(tgt^{-1})$ 的值相同. 于是式\eqref{eq:final2-affine-induced-character}也等于对 $W_g\setminus\{0\}$ 的一次求和. 当 $\sigma=\eta_\tau$ 时, $tgt^{-1}=h(b,d)$ 中的 $d$ 必为 $\det g$, 所以
\[
 \chi_{Q_\tau}(g,v)
 =\tau(\det g)\sum_{y\in W_g\setminus\{0\}}\psi(yv).
\]
对整个 $W_g$ 求和, 若 $v\in\im(g-I)$, 每项都是1, 和为 $q^{k_g}$. 若 $v\notin\im(g-I)$, 由于 $W_g$ 恰是 $\im(g-I)$ 的行向量零化空间, 存在 $y\in W_g$ 使 $yv\ne0$. 将该 $y$ 乘以 $F$ 中适当的标量, 可保证 $\psi(yv)\ne1$. 因而 $y\mapsto\psi(yv)$ 是 $W_g$ 的非平凡加法特征, 其和为0. 去掉零向量的贡献1, 得
\begin{equation}\label{eq:final2-affine-linear-little}
 \chi_{Q_\tau}(g,v)
 =\tau(\det g)\bigl(q^{k_g}\epsilon_g(v)-1\bigr).
\end{equation}
若 $k_g=0$, 则 $\im(g-I)=F^2$, 右端为0, 所以该式也包含没有特征值1的情形.

当 $\sigma=R$ 时, 式\eqref{eq:final2-H-character}还给出
\begin{equation}\label{eq:final2-affine-nonlinear-little}
 \chi_{Q_R}(g,v)=
 \begin{cases}
 (q-1)\bigl(q^2\mathbf1_{\{v=0\}}-1\bigr),&g=I,\\
 1-q\epsilon_g(v),&g\ne I,\ (g-I)^2=0,\\
 0,&\text{其他情形}.
 \end{cases}
\end{equation}
确实, 当 $g=I$ 时, 稳定子表示的每项迹为 $q-1$, 再乘所有非零平移特征的和即可. 若 $g\ne I$ 且 $(g-I)^2=0$, 则 $g$ 为非平凡幺幂矩阵, $k_g=1$. 每个相关共轭矩阵都是 $h(b,1)$, $b\ne0$, 故 $R$ 的特征标值恒为 $-1$, 结果是式\eqref{eq:final2-affine-linear-little}中求和的负值. 其余情形中, 若没有特征值1则求和集合为空; 若有特征值1且 $\det g\ne1$, 式\eqref{eq:final2-H-character}给出的每项为0. 这穷尽了二维矩阵的所有可能性.

特别地, 对非平凡幺幂 $g$, $Q_\tau$ 的值与 $\tau$ 无关, 在 $v\in\im(g-I)$ 时为 $q-1$, 否则为 $-1$; $Q_R$ 的值分别为 $1-q$ 和1. 对特征值为 $1,d$、$d\ne1$ 的矩阵, $Q_\tau$ 的值分别为 $(q-1)\tau(d)$ 和 $-\tau(d)$, 而 $Q_R$ 的值总为0. 这些值都由上述闭式公式直接得到.

\textbf{六、$q=2$ 的边界检验.}
此时 $K$ 置换三个非零向量, 得到 $K\simeq S_3$. 其三个不可约表示是 $L_1$、$\mathrm{St}_1$ 和唯一的 $C_\theta$, 维数分别为 $1,2,1$. 在 $I,J_1,E_z$ 上的值分别为
\[
 \chi_{L_1}=(1,1,1),\qquad
 \chi_{\mathrm{St}_1}=(2,0,-1),\qquad
 \chi_{C_\theta}=(1,-1,1).
\]
最后一个正是 $S_3$ 的符号表示. $H\simeq C_2$ 的两个不可约表示都是一维的, 分别为 $\eta_1$ 与 $R$. 因而 $\Gamma$ 的维数列表为 $1,1,2,3,3$, 平方和为24. 仿射群忠实作用于 $F_2^2$ 的四个点, 阶为24, 故 $\Gamma\simeq S_4$.

还可逐类核对公式. 非零纯平移在四个点上是两个不交的换位. 非平凡幺幂 $g$ 满足 $\im(g-I)=\ker(g-I)$: 若 $v\in\im(g-I)$, 则 $(g,v)$ 与 $(g,0)$ 共轭, 在四个点上是一个换位; 若 $v\notin\im(g-I)$, 则 $(g,v)^2=(I,(g-I)v)$ 是非零平移, 所以 $(g,v)$ 是一个4轮换. 非分裂 $g$ 的阶为3且 $g-I$ 可逆, 所有 $(g,v)$ 都与 $(g,0)$ 共轭, 在四个点上是3轮换. 按这些五类, 式\eqref{eq:final2-affine-linear-little}与式\eqref{eq:final2-affine-nonlinear-little}分别给出
\begin{center}
\begin{tabular}{@{}llllll@{}}
\toprule
 & 单位元 & 换位 & 双换位 & 3轮换 & 4轮换 \\
\midrule
$\chi_{Q_1}$ & $3$ & $1$ & $-1$ & $0$ & $-1$ \\
$\chi_{Q_R}$ & $3$ & $-1$ & $-1$ & $0$ & $1$ \\
\bottomrule
\end{tabular}
\end{center}
这与 $S_4$ 的两个三维不可约特征标一致, 也检查了偶特征下的符号、平移项及边界维数. 至此, 式\eqref{eq:final2-all-affine}给出全部表示, 两张 $K$ 的表与式\eqref{eq:final2-affine-linear-little}、\eqref{eq:final2-affine-nonlinear-little}给出它们在每一个仿射元素上的特征标.
\end{solution}

\originalsection{期末作业第3题}
\begin{exercise}\label{final:3}
MIT 18.712 (2010) 期末作业第3题. 设 $\pi:\mathrm{SU}(2)\to\mathrm{SO}(3)$ 为通常的二重覆盖, $A_5\subset\mathrm{SO}(3)$ 为正二十面体的旋转群, $\Gamma=\pi^{-1}(A_5)$. 计算从 $\Gamma$ 的所有循环子群诱导得到的表示的不可约分解. 此群对应仿射 Dynkin 图 $\widetilde E_8$.
\end{exercise}
\begin{solution}
以下解答为编者独立推导的中文解答, 不是官方公布的答案. 我们先构造不可约表示, 再把诱导重数化成循环子群上的特征值重数. 最后的表和公式同时给出每个循环子群、每个一维特征标的答案; 任意循环子群表示的答案再按直和相加.

\textbf{群元素与循环子群.}
记 $z=-I$, 则 $\ker\pi=\{1,z\}$, $|\Gamma|=120$. 任意 $g\in\mathrm{SU}(2)$ 的特征值可写成 $e^{\mathrm i\theta},e^{-\mathrm i\theta}$, 其像是旋转角为 $2\theta$ 的空间旋转. 因而 $A_5$ 中阶数为 $2,3,5$ 的旋转轴在 $\Gamma$ 中分别给出最大循环子群 $H_4,H_6,H_{10}$. 选生成元 $g_m$ 使自然二维表示 $V=\C^2$ 上的特征值为 $\zeta_m,\zeta_m^{-1}$, 其中 $\zeta_m=e^{2\pi\mathrm i/m}$. 取
\[
 H_m=\langle g_m\rangle\quad(m=4,6,10),\qquad
 H_2=\langle z\rangle,\quad H_3=\langle g_6^2\rangle,\quad
 H_5=\langle g_{10}^2\rangle,\quad H_1=\{1\}.
\]
这些给出循环子群的全部共轭类型. 理由是 $A_5$ 的非平凡元素阶数只有 $2,3,5$, 每个旋转轴的旋转组成一个循环群; 而非中心的 $\mathrm{SU}(2)$ 元素有两个不同特征值, 与它交换的矩阵必须保持两条特征直线, 所以其中心化子是该轴对应的圆群. 与 $\Gamma$ 相交恰为上述最大循环群. 同一种轴的子群在 $A_5$ 中共轭, 提升共轭元便得到 $\Gamma$ 中的共轭. 阶数为 $3,5$ 的子群是相应最大循环群中唯一的同阶子群, 阶数为 $2$ 的子群则唯一等于 $\langle z\rangle$. 按阶数 $1,2,3,4,5,6,10$ 排列, 实际子群数分别为 $1,1,10,15,6,10,6$, 共 $49$ 个循环子群.

为计算特征标, 令 $\varphi=(1+\sqrt5)/2$, $\tau=\varphi-1$. 群中共轭类可由 $t=\tr_V(g)$ 区分, 对应的大小如下:
\[
\begin{array}{c|rrrrrrrrr}
t&2&-2&0&1&-1&\varphi&-\varphi&\tau&-\tau\\ \hline
\text{类大小}&1&1&30&20&20&12&12&12&12
\end{array}
\]
这里 $A_5$ 的双换位有 $15$ 个, 三循环有 $20$ 个, 五循环有 $24$ 个. 五循环的中心化子阶数为 $5$, 所以分成两个各有 $12$ 个元素的共轭类, 分别对应旋转角 $2\pi/5$ 与 $4\pi/5$ (同时包括反向旋转). 提升后, 非零迹的两种提升分别形成一个类: 相应中心化子阶数为 $6$ 或 $10$, 类大小为 $20$ 或 $12$. 双换位的两个提升迹都为零, 在 $\Gamma$ 中中心化子阶数为 $4$, 合成一个大小为 $30$ 的类. 这也直接核对了九个类的总大小为 $120$.

\textbf{全部九个不可约表示的构造.}
若 $g$ 在 $V$ 上的特征值为 $\lambda,\lambda^{-1}$, 则它在 $\Sym^nV$ 上的特征值为 $\lambda^{n-2r}$, $0\le r\le n$. 因此其特征标 $P_n(t)$ 满足
\[
 P_0=1,\quad P_1=t,\quad P_{n+1}=tP_n-P_{n-1}.
\]
我们用记号
\[
 \mathbf1=\Sym^0V,\quad A=\Sym^2V,\quad C=\Sym^3V,
 \quad D=\Sym^4V,\quad E=\Sym^5V.
\]
这些实际表示的维数依次是 $1,3,4,5,6$; 另有自然表示 $V$ 的维数为 $2$. 为了取得余下的表示, 注意 $A$ 因 $z$ 的作用为 $1$ 而下降到 $A_5$. 旋转角 $2\theta$ 的三维旋转表示的特征标为 $1+2\cos(2\theta)=t^2-1=P_2(t)$, 因而它的复化与 $A$ 同构. 将这个 $A_5$ 表示与 $S_5$ 中由奇置换共轭给出的自同构复合, 再拉回 $\Gamma$, 得到三维表示 $A'$. 此自同构交换两个五循环类. 例如对 $c=(12345)$, 把 $c$ 变成 $c^2$ 的正规化置换在其五个位置上相当于 $r\mapsto2r$ ($r\in\Z/5\Z$), 在四个非零位置构成四循环, 故为奇置换; 所以奇共轭确实交换两类.

另取 $A_5$ 在五个字母上的置换表示中坐标和为零的四维子空间, 再拉回 $\Gamma$, 记作 $B$. 它的特征标是固定字母数减 $1$: 在恒等元、双换位、三循环、五循环处分别为 $4,0,1,-1$. 最后暂在虚表示环中定义
\[
 V'=V\otimes B-E.
\]
不能仅从这个差式断言 $V'$ 是实际表示; 下面的正交计算会证明这一点.

用递推式展开 $P_n$, 并用 $\varphi^2=\varphi+1$, $\varphi\tau=1$ 简化, 得到完整特征标表:
\begin{center}
\begin{tabular}{c|rrrrrrrrr}
\toprule
 & $2$&$-2$&$0$&$1$&$-1$&$\varphi$&$-\varphi$&$\tau$&$-\tau$\\
\midrule
$\mathbf1$&1&1&1&1&1&1&1&1&1\\
$V$&2&$-2$&0&1&$-1$&$\varphi$&$-\varphi$&$\tau$&$-\tau$\\
$V'$&2&$-2$&0&1&$-1$&$-\tau$&$\tau$&$-\varphi$&$\varphi$\\
$A$&3&3&$-1$&0&0&$\varphi$&$\varphi$&$-\tau$&$-\tau$\\
$A'$&3&3&$-1$&0&0&$-\tau$&$-\tau$&$\varphi$&$\varphi$\\
$B$&4&4&0&1&1&$-1$&$-1$&$-1$&$-1$\\
$C$&4&$-4$&0&$-1$&1&1&$-1$&$-1$&1\\
$D$&5&5&1&$-1$&$-1$&0&0&0&0\\
$E$&6&$-6$&0&0&0&$-1$&1&1&$-1$\\
\bottomrule
\end{tabular}
\end{center}
表头是自然表示的迹 $t$, 不是表示自身的维数或元素阶数. 所有行的数值均为实数. 将类大小记为 $s_t$, 对任意两行 $X,Y$ 计算
\[
 \inner{\chi_X}{\chi_Y}
 =\frac1{120}\sum_t s_t\chi_X(t)\chi_Y(t).
\]
所得九行的 Gram 矩阵就是单位矩阵. 这项计算只涉及表内九项, 也可逐行核对: 例如
\[
 \inner{\chi_V}{\chi_V}
 =\frac{8+40+24(\varphi^2+\tau^2)}{120}=1,
 \qquad \varphi^2+\tau^2=3.
\]
对其余行同样代入 $\varphi+\tau=\sqrt5$, $\varphi-\tau=1$, $\varphi\tau=1$, 便得对角项 $1$、非对角项 $0$. 对实际表示而言, 特征标范数为 $1$ 表明表示不可约. 对虚表示 $V'$, 若其不可约展开为 $\sum_i n_iX_i$, 则 $n_i\in\Z$ 且 $\sum_i n_i^2=1$, 所以只有一项, 系数为 $\pm1$. 又 $\chi_{V'}(1)=2>0$, 故系数为 $1$, $V'$ 确为实际不可约表示. 因而上述九个不可约表示两两不等价. 最后
\[
 1^2+2^2+2^2+3^2+3^2+4^2+4^2+5^2+6^2=120
\]
说明它们已经穷尽全部不可约表示. 特别要注意, 下一阶对称幂不再不可约: 表中逐项相加给出 $\Sym^6V\cong A'\oplus B$.

题面中的 $\widetilde E_8$ 也可由此核对, 无须预先调用图的分类定理. 将每一行乘以 $\chi_V=t$, 再按上表分解, 得
\begin{align*}
V\otimes\mathbf1&=V,& V\otimes V&=\mathbf1\oplus A,& V\otimes A&=V\oplus C,\\
V\otimes C&=A\oplus D,& V\otimes D&=C\oplus E,& V\otimes E&=D\oplus B\oplus A',\\
V\otimes B&=E\oplus V',& V\otimes V'&=B,& V\otimes A'&=E.
\end{align*}
所以 McKay 图是一条 $\mathbf1,V,A,C,D,E,B,V'$ 的链, 并在 $E$ 处接一个端点 $A'$, 即仿射 $E_8$ 图. 这里每条边的含义已经由实际张量积分解证明, 不靠图名来推测重数.

\textbf{三个最大循环子群上的完整权重.}
对 $m=4,6,10$, 定义非负整数 $a_{X,m}(r)$ 为 $g_m$ 在 $X$ 上的特征值 $\zeta_m^r$ 的重数, $0\le r<m$. 下列三个表列出所有这些整数. 对 $\Sym^nV$, 直接把整数 $n,n-2,\ldots,-n$ 按模 $m$ 计数即可. 对 $A'$, 奇共轭在五阶旋转上把角 $2\pi/5$ 换成 $4\pi/5$, 所以 $H_{10}$ 上的三个权重为 $0,4,6$; 在 $H_4,H_6$ 上与 $A$ 相同. 对 $B$, 则用五点置换的循环结构减去一份平凡特征值. 对 $V'$, 从 $B$ 的权重分别加 $1$、减 $1$, 再减去 $E$ 的权重, 得
\[
 a_{V',m}(r)=a_{B,m}(r-1)+a_{B,m}(r+1)-a_{E,m}(r),
\]
其中下标按模 $m$ 理解. 这是从上面的虚表示构造计算特征值重数, 而其为非负整数则由下表实际核实.

\begin{center}
\begin{tabular}{c|rrrr}
\toprule
$X\backslash r$ ($m=4$)&0&1&2&3\\ \midrule
$\mathbf1$&1&0&0&0\\
$V$&0&1&0&1\\
$V'$&0&1&0&1\\
$A$&1&0&2&0\\
$A'$&1&0&2&0\\
$B$&2&0&2&0\\
$C$&0&2&0&2\\
$D$&3&0&2&0\\
$E$&0&3&0&3\\ \bottomrule
\end{tabular}
\qquad
\begin{tabular}{c|rrrrrr}
\toprule
$X\backslash r$ ($m=6$)&0&1&2&3&4&5\\ \midrule
$\mathbf1$&1&0&0&0&0&0\\
$V$&0&1&0&0&0&1\\
$V'$&0&1&0&0&0&1\\
$A$&1&0&1&0&1&0\\
$A'$&1&0&1&0&1&0\\
$B$&2&0&1&0&1&0\\
$C$&0&1&0&2&0&1\\
$D$&1&0&2&0&2&0\\
$E$&0&2&0&2&0&2\\ \bottomrule
\end{tabular}
\end{center}
\begin{center}
\begin{tabular}{c|rrrrrrrrrr}
\toprule
$X\backslash r$ ($m=10$)&0&1&2&3&4&5&6&7&8&9\\ \midrule
$\mathbf1$&1&0&0&0&0&0&0&0&0&0\\
$V$&0&1&0&0&0&0&0&0&0&1\\
$V'$&0&0&0&1&0&0&0&1&0&0\\
$A$&1&0&1&0&0&0&0&0&1&0\\
$A'$&1&0&0&0&1&0&1&0&0&0\\
$B$&0&0&1&0&1&0&1&0&1&0\\
$C$&0&1&0&1&0&0&0&1&0&1\\
$D$&1&0&1&0&1&0&1&0&1&0\\
$E$&0&1&0&1&0&2&0&1&0&1\\ \bottomrule
\end{tabular}
\end{center}
每行之和都是相应表示的维数. 例如 $E=\Sym^5V$ 在 $H_{10}$ 上的权重为 $5,3,1,-1,-3,-5$, 故权重 $5$ 出现两次, 这解释了表中唯一的重数 $2$ 的中间位置. 一般也可从特征标表以离散 Fourier 变换重算每项:
\[
 a_{X,m}(r)=\frac1m\sum_{k=0}^{m-1}
 \chi_X(g_m^k)\zeta_m^{-rk},\qquad
 \tr_V(g_m^k)=2\cos\frac{2\pi k}{m}.
\]
这是因为 $m^{-1}\sum_{k=0}^{m-1}\zeta_m^{(s-r)k}$ 在 $s\equiv r\pmod m$ 时为 $1$, 否则为 $0$.

\textbf{所有诱导表示的答案.}
对 $H_d$ 的生成元 $h_d$, 记一维特征标为 $\eta_{d,j}(h_d)=\zeta_d^j$, $0\le j<d$. 对最大子群取 $h_d=g_d$. 对 $d=2,3,5$, 分别选 $(m,d)=(4,2),(6,3),(10,5)$, 并取 $h_d=g_m^{m/d}$. Frobenius 互反给出
\[
 [\Ind_{H_d}^{\Gamma}\eta_{d,j}:X]
 =[\Res_{H_d}^{\Gamma}X:\eta_{d,j}]
 =\sum_{\substack{0\le r<m\\r\equiv j\ (\mathrm{mod}\ d)}} a_{X,m}(r).
\]
因此对 $d=4,6,10$, 第 $j$ 列就是 $\Ind_{H_d}^{\Gamma}\eta_{d,j}$ 的九个不可约重数; 对 $d=2,3,5$, 将上表中下标模 $d$ 相同的列相加就是答案. 为使较小子群的结果也能直接读取, 列出折叠后的三个表:
\begin{center}
\begin{tabular}{c|rr|rrr|rrrrr}
\toprule
&\multicolumn{2}{c|}{$d=2$}&\multicolumn{3}{c|}{$d=3$}&\multicolumn{5}{c}{$d=5$}\\
$X\backslash j$&0&1&0&1&2&0&1&2&3&4\\ \midrule
$\mathbf1$&1&0&1&0&0&1&0&0&0&0\\
$V$&0&2&0&1&1&0&1&0&0&1\\
$V'$&0&2&0&1&1&0&0&1&1&0\\
$A$&3&0&1&1&1&1&0&1&1&0\\
$A'$&3&0&1&1&1&1&1&0&0&1\\
$B$&4&0&2&1&1&0&1&1&1&1\\
$C$&0&4&2&1&1&0&1&1&1&1\\
$D$&5&0&1&2&2&1&1&1&1&1\\
$E$&0&6&2&2&2&2&1&1&1&1\\ \bottomrule
\end{tabular}
\end{center}
这里每一列始终按行 $\mathbf1,V,V',A,A',B,C,D,E$ 解读为直和重数. 例如
\begin{align*}
\Ind_{H_{10}}^{\Gamma}\eta_{10,0}&=\mathbf1\oplus A\oplus A'\oplus D,\\
\Ind_{H_{10}}^{\Gamma}\eta_{10,5}&=2E,\\
\Ind_{H_5}^{\Gamma}\eta_{5,0}&=\mathbf1\oplus A\oplus A'\oplus D\oplus2E,\\
\Ind_{H_3}^{\Gamma}\eta_{3,1}&=V\oplus V'\oplus A\oplus A'\oplus B\oplus C\oplus2D\oplus2E.
\end{align*}
从平凡子群诱导则是正则表示, 即
\[
 \Ind_{H_1}^{\Gamma}\mathbf1
 =\mathbf1\oplus2V\oplus2V'\oplus3A\oplus3A'
 \oplus4B\oplus4C\oplus5D\oplus6E.
\]
上述七种阶数 $1,2,3,4,5,6,10$ 的子群及其全部特征标都已列全. 若 $U=\bigoplus_{j=0}^{d-1}b_j\eta_{d,j}$ 是任意 $H_d$ 表示, 则 $\Ind U$ 中 $X$ 的重数为上式乘 $b_j$ 后对 $j$ 求和.

最后说明选择生成元与共轭的影响. 将子群共轭, 并同时搬运它的特征标, 不改变诱导表示的同构类型. 若固定一个子群而改取生成元 $\widetilde h=h^u$, $(u,d)=1$, 则以 $\widetilde h$ 标为 $j$ 的特征标, 在旧生成元下标为 $u^{-1}j\pmod d$; 因而只需把相应列置换. 特别地, 表中各列关于 $j\mapsto-j$ 对称, 但对 $H_{10}$, 改取 $g_{10}^3$ 必须按上述规则重新标列, 不能仍把自然表示的权重标成 $\pm1$.

所有表项非负, 每列按不可约维数加权之和分别为 $120/d$, 与诱导表示维数相符. 对阶数 $2$ 的折叠, 无论从 $H_4,H_6$ 还是 $H_{10}$ 进行, 都得到同一张表: $z$ 在 $\mathbf1,A,A',B,D$ 上作用为 $1$, 在 $V,V',C,E$ 上作用为 $-1$. 这些检查同时覆盖中心元、平凡特征标、重复特征值和奇数阶子群的边界情形.
\end{solution}

\originalsection{期末作业第4题}
\begin{exercise}\label{final:4}
\textbf{MIT 18.712（2010）期末作业原题4；官方期末PDF第1页.}
设 $V$ 是 $\mathfrak{sl}_2(\C)$ 的二维向量表示。求 $V^{\otimes n}$ 中各不可约表示的重数。
\end{exercise}
\begin{solution}
本题所需的李代数知识在此补齐。$\mathfrak{sl}_2(\C)$ 是迹为零的 $2\times2$ 复矩阵组成的李代数，括号为 $[X,Y]=XY-YX$。取
\[
 e=\begin{pmatrix}0&1\\0&0\end{pmatrix},\quad
 f=\begin{pmatrix}0&0\\1&0\end{pmatrix},\quad
 h=\begin{pmatrix}1&0\\0&-1\end{pmatrix},
\]
则 $[h,e]=2e,[h,f]=-2f,[e,f]=h$。李代数在张量积上的作用是
\[
 X(u_1\otimes\cdots\otimes u_n)
 =\sum_{j=1}^n u_1\otimes\cdots\otimes Xu_j\otimes\cdots\otimes u_n.
\]
下文只证明本题需要的不可约表示与分解，不把这段增补视为覆盖旧讲义的整个李代数专题。

记 $L_m=\Sym^m V$，$m\ge0$。把它写成变量 $x,y$ 的 $m$ 次齐次多项式空间，取基 $v_j=x^{m-j}y^j$，$0\le j\le m$，则
\[
 ev_j=jv_{j-1},\qquad fv_j=(m-j)v_{j+1},\qquad hv_j=(m-2j)v_j.
\]
不存在的端点基向量按零处理。$h$ 的特征值两两不同；任一不变子空间经 $h$ 的多项式投影含有某个 $v_j$，而 $e,f$ 又把所有基向量连在一起，故 $L_m$ 不可约。它的维数为 $m+1$，权为 $m,m-2,\ldots,-m$，每个权的重数为一。

这些也是所有有限维不可约表示。事实上，$e$ 把 $h$ 的广义特征空间的权提高二，$f$ 则降低二；有限谱使 $e,f$ 都幂零，这里不需要预先假定 $h$ 可对角化。$\ker e$ 非零且在 $h$ 下不变，所以可取 $ev=0,hv=\lambda v$。由交换关系归纳得到
\[
 e f^jv=j(\lambda-j+1)f^{j-1}v,\qquad
 h f^jv=(\lambda-2j)f^jv.
\]
设 $f^m v\ne0,f^{m+1}v=0$。代入 $j=m+1$ 得到 $\lambda=m$。所生成的子空间有基 $v,fv,\ldots,f^mv$，在 $e,f,h$ 下不变；不可约性使它等于整个表示，并由上面的作用公式与 $L_m$ 同构。

为在本题中获得完全可约性，无需引用一般李代数的完全可约定理。直接考虑乘法映射
\[
 \mu:L_m\otimes V\longrightarrow L_{m+1},\qquad p\otimes z\longmapsto pz.
\]
它满射且与 $e,f,h$ 的作用相容。映射
\[
 \delta(p)=\frac1{m+1}\left(\frac{\partial p}{\partial x}\otimes x+
               \frac{\partial p}{\partial y}\otimes y\right)
\]
也是表示同态；对 $e,f,h$ 逐一应用乘积求导法则即可验证。Euler齐次公式给出 $\mu\delta=\id$。当 $m\ge1$ 时，映射
\[
 L_{m-1}\longrightarrow\ker\mu,\qquad
 p\longmapsto xp\otimes y-yp\otimes x
\]
是单射且与作用相容，因为 $x\otimes y-y\otimes x$ 在 $\mathfrak{sl}_2$ 下不变。比较维数，它的像就是 $\ker\mu$。因此
\[
 L_m\otimes V\cong L_{m+1}\oplus L_{m-1}\quad(m\ge1),
 \qquad L_0\otimes V\cong L_1.
\]
从 $V^{\otimes0}=L_0$ 出发归纳，$V^{\otimes n}$ 是若干 $L_m$ 的直和。

现在数权空间。$V$ 的两个基向量的权为 $1,-1$。张量基中若有 $r$ 个权为 $-1$ 的因子，总权就是 $n-2r$，所以权 $m$ 的重数为
\[
 a_{n,m}=\binom n{(n-m)/2}
\]
（下标非整数或不在 $[0,n]$ 时按零处理）。对 $m\ge0$，每个最高权为 $m,m+2,m+4,\ldots$ 的不可约直和项都向权 $m$ 贡献一维。若 $M_{n,m}$ 是 $L_m$ 的重数，则
\[
 a_{n,m}=M_{n,m}+M_{n,m+2}+\cdots.
\]
相邻两式相减，得到最终答案
\[
 \boxed{\displaystyle
 M_{n,m}=\binom n r-\binom n{r-1},\qquad r=\frac{n-m}{2},
 \quad 0\le m\le n,\quad m\equiv n\pmod2.}
\]
其余 $m$ 的重数为零，约定 $\binom n{-1}=0$。等价地，对上述允许的 $m$，
\[
 M_{n,m}=\frac{m+1}{n-r+1}\binom n r.
\]
例如 $V^{\otimes2}=L_2\oplus L_0$，$V^{\otimes3}=L_3\oplus2L_1$，
$V^{\otimes4}=L_4\oplus3L_2\oplus2L_0$；最后一个分解的维数为 $5+3\cdot3+2=16$，与 $2^4$ 一致。
\end{solution}

\originalchapter{来源与范围对应}
\originalsection{正文主线}
本表中的页码是公开旧讲义PDF页序, 与其印刷页码一致; 末页109为OCW条款. 分章PDF分别重复完整PDF相应正文, 另附条款页, 所以仅归档完整旧PDF以去重. 中文正文没有嵌入英文原页.
\begin{longtable}{p{.27\textwidth} p{.27\textwidth} p{.38\textwidth}}
\toprule 本书内容&旧讲义定位&编排与边界\\\midrule\endhead
第1章群与表示&§1.1--1.3、§1.7--1.8, 页5--14&群论必要语言与置换矩阵为衔接补充; 不展开一般代数分类\\
第2章完全可约&§1.3、§2.1、§3.1、§3.6, 页7--10、23、33--34、38--39&Schur、Maschke、平均内积与重数唯一性完整证明\\
第3章特征标&§2.6、§3.2--3.5, 页27--28、34--38&以Hom空间平均为统一证明; 置换轨道计数为补充\\
第4章正则与群代数&§2.2--2.5、§3.1、§3.7、§4.2, 页24--27、33、39--40、48--49&专门重证复群代数分块; 加入中心投影与反演; 不展开一般有限维代数根基\\
第5章经典表&§3.3、§3.8, 页35--36、40--42&循环群、$S_3$、$D_8$、$Q_8$、$A_4$, 补齐表的构造和完整性检查\\
第6章新表示&§1.10、§3.4、§3.9、§4.6--4.7, 页17--18、36--37、42--43、54&张量、对偶、对称与外平方、直积与虚表示\\
第7章诱导&§4.8--4.11, 页54--58&张量模型与函数模型互证; 两方向互反、传递性、投影公式\\
第8章双陪集&§4.9--4.11、§4.26, 页55--58、76--77&双陪集分解及判据为编者展开; 二面体群统一分类\\
第9章形式与$S_4$&§4.1、§3.8及§4.11, 页47--48、40--42、57--58&指标完整证明; $S_4$用商群与置换空间构造, 不调用Young图分类\\\bottomrule
\end{longtable}
前九章有编者例题与练习, 不冒充官方作业编号. 课程题解部分则以Homework和原题号保留身份, 共52道作业题及4道期末题. 旧讲义中的未布置练习并未全部收录, 官方未公开的题卷也不据书目或编号臆造.

\originalsection{课程题目与拓展范围}
Homework 1--4的题面来自旧PDF页10--15、18--22, 包含结合代数、Weyl代数、张量、Dehn不变量、$\mathfrak{sl}_2$及可解李代数. Homework 5--9来自页30--31、34、43--47、51、55、57--59、63、68, 包含扩张、群论应用、诱导、对称群及一般线性群局部题解. Homework 10--11来自页78--81、96--97, 涉及箭图与根系. 每题解答保留全部小问; 字符抽取损坏、原题笔误及需要补明的底域条件在对应位置说明.

期末作业单独官方PDF为2页, 第1页有4道题, 第2页为OCW条款. 本书分别处理路径代数、仿射群 $\GL_2(\mathbb F_q)\ltimes\mathbb F_q^2$、二元二十面体群的循环子群诱导, 以及 $\mathfrak{sl}_2$ 张量幂. 全部答案由编者推导并交叉复核; 官网无独立Solutions文件, 不能将本书答案称为官方标准答案.

课程题解为了处理原题而加入的局部理论, 不等于主源完整专题的系统覆盖. 李代数、箭图、根系、$\GL_2(\mathbb F_q)$和有限群半直积在题解中按需要展开; 范畴论、Gabriel定理全套证明、一般代数投射覆盖、完整Young图与Schur--Weyl理论、Artin定理、Frobenius整除性和Burnside可解性仍非本书完整专题.

\begin{thebibliography}{9}
\bibitem{old} Pavel Etingof, Oleg Golberg, Sebastian Hensel, Tiankai Liu, Alex Schwendner, Elena Yudovina, Dmitry Vaintrob. \textit{Introduction to Representation Theory}. 公开旧讲义, 2011-02-01, MIT OCW 18.712, Fall 2010. \href{https://ocw.mit.edu/courses/18-712-introduction-to-representation-theory-fall-2010/24d8b3fa2ce48e48ee6c2d8d5e3562f6_MIT18_712F10_replect.pdf}{109页完整旧PDF}.
\bibitem{course} MIT OpenCourseWare. \href{https://ocw.mit.edu/courses/18-712-introduction-to-representation-theory-fall-2010/pages/lecture-notes/}{18.712官方讲义目录}, 含旧版分章文件及2016出版书转载限制说明.
\bibitem{license} MIT OpenCourseWare. \href{https://ocw.mit.edu/pages/privacy-and-terms-of-use/}{使用条款}; \href{https://creativecommons.org/licenses/by-nc-sa/4.0/}{CC BY-NC-SA 4.0}.
\end{thebibliography}
\end{document}
